% Lesson 10 — Grammar: Phrasal verbs and prepositional phrases — Anglais Tle A — Cours signé OnBuch+
% Programme officiel MINESEC. Compiler avec : tectonic EN10-grammar-phrasal-verbs-and-prepositional-phrases.tex
\newcommand{\DOCMATIERE}{Anglais}
\newcommand{\DOCNIVEAU}{Tle A}
\newcommand{\DOCMODULE}{Chapter 1 : Family and Social Life}
\newcommand{\DOCLECON}{EN10}
\newcommand{\DOCTITRE}{Lesson 10 — Grammar: Phrasal verbs and prepositional phrases}
% OnBuch+ — préambule « cours signés » ENRICHI (compilé avec tectonic / XeLaTeX)
% Système visuel « Pop » d'OnBuch+ : fond crème (jamais blanc), bordures encre
% épaisses, ombres « dures » décalées, titres Archivo Black en capitales.
% Version MATHÉMATIQUES : graphes (pgfplots/tikz), boîtes pédagogiques
% (définition, propriété, méthode, attention, le savais-tu, exercice résolu,
% exercice, TP, bilan), page de garde et signature OnBuch+ ancrée en bas.
%
% Macros attendues dans main.tex AVANT \input{preamble} :
%   \DOCMATIERE  \DOCNIVEAU  \DOCMODULE  \DOCLECON (numéro)  \DOCTITRE
\documentclass[11pt]{article}

\usepackage{fontspec}
\usepackage[english,french]{babel} % français = langue principale ; l'anglais via \eng{..} / textebox
\usepackage[a4paper,margin=2cm,top=3.4cm,bottom=2.8cm,headheight=1.4cm,footskip=1.3cm]{geometry}
\usepackage{amsmath,amssymb,amsfonts}
\usepackage{mathtools} % \xmapsto, \coloneqq... souvent utilisés par les modèles
\usepackage{cancel} % simplifications barrées (bilans de forces)
% \abs{z} (module, valeur absolue) et \norm{u} (norme), utilisés par les modèles
\providecommand{\abs}{}\let\abs\relax\DeclarePairedDelimiter{\abs}{\lvert}{\rvert}
\providecommand{\norm}{}\let\norm\relax\DeclarePairedDelimiter{\norm}{\lVert}{\rVert}
\usepackage[table]{xcolor} % [table] : fournit aussi \rowcolors (alternance de lignes)
\usepackage{pagecolor}
\usepackage{tikz}
\usetikzlibrary{arrows.meta,positioning,calc,shapes.geometric,shapes.misc,shapes.arrows,decorations.pathmorphing,decorations.pathreplacing,decorations.markings,patterns,shadows,mindmap,trees,fit,backgrounds,matrix,intersections,angles,quotes}
\usepackage{pgfplots}
\usepackage[version=4]{mhchem} % équations nucléaires \ce{^{235}_{92}U}
\usepackage[european,siunitx]{circuitikz} % schémas électriques
\pgfplotsset{compat=1.18}
\usepgfplotslibrary{fillbetween,groupplots} % aires entre courbes (calcul intégral)
\usepackage{enumitem}
\usepackage{fancyhdr}
% [most] charge toutes les bibliothèques tcolorbox (listings, minted, raster,
% documentation...) et gonfle inutilement la mémoire TeX sur un document long
% (~300 boîtes) : on ne charge que ce qui est réellement utilisé.
\usepackage[skins,breakable]{tcolorbox}
\usepackage{graphicx}
\usepackage{colortbl}
\usepackage{booktabs}
\usepackage{tabularx}
\usepackage{diagbox} % case d'en-tête diagonale des tableaux à double entrée
% Colonnes extensibles centrée / gauche / droite (C, L, R), souvent utilisées par les modèles
\newcolumntype{C}{>{\centering\arraybackslash}X}
\newcolumntype{L}{>{\raggedright\arraybackslash}X}
\newcolumntype{R}{>{\raggedleft\arraybackslash}X}
\usepackage{array}
\usepackage{multicol}
\usepackage{multirow}
\usepackage{siunitx}
% parse-numbers=false : accepte aussi \SI{2,5 \times 10^{-3}}{...}
\sisetup{locale=FR,output-decimal-marker={,},group-minimum-digits=4,parse-numbers=false}
\DeclareSIUnit\ohm{\ensuremath{\symup{\Omega}}} % U+2126 absent de la police
\DeclareSIUnit\division{div} % réglages d'oscilloscope (ms/div, V/div)
\DeclareSIUnit\tr{tr} % tours (vitesse de rotation, ex. \SI{1500}{\tr\per\minute})
\DeclareSIUnit\molar{M} % concentration molaire (ex. \SI{3}{\molar}, \SI{1}{\micro\molar})
\DeclareSIUnit\an{an} % année (le modèle confond parfois avec \year, un compteur TeX) — ex. \SI{5}{\kilo\meter\per\an}
\DeclareSIUnit\FCFA{FCFA} % franc CFA (ex. \SI{500}{\FCFA\per\kilo\gram})
\DeclareSIUnit\degreeBrix{\textdegree Brix} % teneur en sucre d'un sirop/jus (réfractométrie)
\DeclareSIUnit\month{mois} % le modèle utilise parfois \month, un compteur TeX, comme unité de durée
\usepackage{qrcode}
% \degree : commande fréquemment utilisée par les modèles pour un angle ou une
% température en dehors de \SI{}{} (non fournie par les paquets chargés).
\providecommand{\degree}{^{\circ}}
% \qed (fin de démonstration), utilisé par les modèles sans amsthm
\providecommand{\qed}{\hfill$\square$}
% \barre : double barre des valeurs interdites dans un tableau de variations (tkz-tab)
\providecommand{\barre}{\ensuremath{\|}}
% environnement proof (amsthm non chargé) utilisé par les modèles
\ifdefined\proof\else\newenvironment{proof}[1][Démonstration]{\par\noindent\textit{#1.}\ }{\qed\par}\fi
\usepackage{lastpage}
\usepackage{hyperref}
\hypersetup{hidelinks}

% ============================================================
% Palette « Pop » OnBuch+ — mêmes hex que la classe P (plus_ui.dart)
% ============================================================
\definecolor{poporange}{HTML}{F59321}
\definecolor{popdark}{HTML}{E07A0C}
\definecolor{popcream}{HTML}{FFF6E8}
\definecolor{popink}{HTML}{1C1714}
\definecolor{poppurple}{HTML}{7A5AE0}
\definecolor{popgreen}{HTML}{2E9E6A}
\definecolor{poppink}{HTML}{E0457B}
\definecolor{popblue}{HTML}{2F7DE1}
\definecolor{popgold}{HTML}{C98A12}
\definecolor{popmuted}{HTML}{8A7F76}
\definecolor{popline}{HTML}{EADFD2}
\definecolor{popgray}{HTML}{8A7F76}
\colorlet{popgrey}{popgray}
% Alias tolérants (le modèle nomme parfois les couleurs différemment)
\colorlet{popwhite}{white}
\colorlet{popblack}{popink}
\colorlet{popred}{poppink}
\colorlet{popyellow}{popgold}
% Teintes claires (fonds de tableaux/encadrés) — le modèle nomme parfois
% les couleurs avec un suffixe L (ex. popblueL) sans les avoir définies.
\colorlet{poporangeL}{poporange!20}
\colorlet{popdarkL}{popdark!20}
\colorlet{poppurpleL}{poppurple!20}
\colorlet{popgreenL}{popgreen!20}
\colorlet{poppinkL}{poppink!20}
\colorlet{popblueL}{popblue!20}
\colorlet{popgoldL}{popgold!20}
\colorlet{popredL}{popred!20}
\colorlet{popgrayL}{popgray!20}
% Teintes claires pour fonds de tableaux / zones
\colorlet{popblueL}{popblue!10!white}
\colorlet{poporangeL}{poporange!14!white}
\colorlet{popgreenL}{popgreen!12!white}
\colorlet{poppurpleL}{poppurple!12!white}
\colorlet{poppinkL}{poppink!12!white}
\colorlet{popgoldL}{popgold!14!white}

\pagecolor{popcream}

% ============================================================
% Typographie
% ============================================================
\defaultfontfeatures{Ligatures=TeX}
\setmainfont{PlusJakartaSans-Medium.ttf}[Path=../../../fonts/,BoldFont=PlusJakartaSans-Bold.ttf]
% Symboles, API et emojis absents de la police principale : repli sur DejaVu Sans (caractères actifs)
\newfontface\symfont{DejaVuSans.ttf}[Path=../../../fonts/]
\makeatletter
\newcommand\symactive[1]{\catcode#1=13 \begingroup\lccode`\~=#1 \lowercase{\endgroup\def~}{{\symfont\char#1}}}
\newcommand\symblank[1]{\catcode#1=13 \begingroup\lccode`\~=#1 \lowercase{\endgroup\def~}{}}
\newcommand\symhyph[1]{\catcode#1=13 \begingroup\lccode`\~=#1 \lowercase{\endgroup\def~}{\mbox{-}}}
\newcount\symc
\newcommand\symrange[2]{\symc=#1 \@whilenum\symc<#2 \do{\expandafter\symactive\expandafter{\the\symc}\advance\symc by 1 }}
\newcommand\symblankrange[2]{\symc=#1 \@whilenum\symc<#2 \do{\expandafter\symblank\expandafter{\the\symc}\advance\symc by 1 }}
\makeatother
\symrange{"0250}{"02FF}\symactive{"2713}\symactive{"2714}\symactive{"2717}\symactive{"2718}\symactive{"2610}\symactive{"2611}\symactive{"2612}
\symhyph{"2011}\symhyph{"2E1A}\symblankrange{"1F300}{"1FAFF}\symblank{"FE0F}
% Maths en sans-serif (Fira Math) avec chiffres et lettres droites de Plus
% Jakarta, pour que les formules se fondent dans le texte.
\usepackage[math-style=ISO,bold-style=ISO]{unicode-math}
\setmathfont{FiraMath-Regular.otf}[Path=../../../fonts/]
\setmathfont{PlusJakartaSans-Medium.ttf}[Path=../../../fonts/,range={up/num,up/{Latin,latin},\mathup}]
\usepackage{icomma} % virgule décimale sans espace en mode math
% Caractères mathématiques Unicode absents des polices de texte (Plus Jakarta,
% Archivo Black) : rendus par leur équivalent mathématique, en texte comme en
% formule — sinon ils s'affichent en carré vide (ex. « Arithmétique dans ℤ »).
\usepackage{newunicodechar}
\newunicodechar{ℝ}{\ensuremath{\mathbb{R}}}
\newunicodechar{ℤ}{\ensuremath{\mathbb{Z}}}
\newunicodechar{ℂ}{\ensuremath{\mathbb{C}}}
\newunicodechar{ℕ}{\ensuremath{\mathbb{N}}}
\newunicodechar{ℚ}{\ensuremath{\mathbb{Q}}}
\newunicodechar{𝔻}{\ensuremath{\mathbb{D}}}
\newunicodechar{α}{\ensuremath{\alpha}}
\newunicodechar{β}{\ensuremath{\beta}}
\newunicodechar{γ}{\ensuremath{\gamma}}
\newunicodechar{δ}{\ensuremath{\delta}}
\newunicodechar{ε}{\ensuremath{\varepsilon}}
\newunicodechar{θ}{\ensuremath{\theta}}
\newunicodechar{λ}{\ensuremath{\lambda}}
\newunicodechar{μ}{\ensuremath{\mu}}
\newunicodechar{σ}{\ensuremath{\sigma}}
\newunicodechar{τ}{\ensuremath{\tau}}
\newunicodechar{φ}{\ensuremath{\varphi}}
\newunicodechar{ψ}{\ensuremath{\psi}}
\newunicodechar{ω}{\ensuremath{\omega}}
\newunicodechar{Σ}{\ensuremath{\Sigma}}
% majuscules produites par \MakeUppercase dans les titres (α-aminés -> Α)
\newunicodechar{Α}{\ensuremath{\alpha}}
\newunicodechar{Β}{\ensuremath{\beta}}
\newunicodechar{Γ}{\ensuremath{\Gamma}}
\newunicodechar{Δ}{\ensuremath{\Delta}}
\newunicodechar{Ε}{\ensuremath{\varepsilon}}
\newunicodechar{Θ}{\ensuremath{\Theta}}
\newunicodechar{Λ}{\ensuremath{\Lambda}}
\newunicodechar{Μ}{\ensuremath{\mu}}
\newunicodechar{Τ}{\ensuremath{\tau}}
\newunicodechar{Φ}{\ensuremath{\Phi}}
\newunicodechar{Ψ}{\ensuremath{\Psi}}
\newunicodechar{Ω}{\ensuremath{\Omega}}
\newunicodechar{↦}{\ensuremath{\mapsto}}
\newunicodechar{∈}{\ensuremath{\in}}
\newunicodechar{∉}{\ensuremath{\notin}}
\newunicodechar{∧}{\ensuremath{\wedge}}
\newunicodechar{⇔}{\ensuremath{\Leftrightarrow}}
\newunicodechar{⟺}{\ensuremath{\Longleftrightarrow}}
\newunicodechar{⇒}{\ensuremath{\Rightarrow}}
\newunicodechar{⟹}{\ensuremath{\Longrightarrow}}
\newunicodechar{∘}{\ensuremath{\circ}}
\newunicodechar{∪}{\ensuremath{\cup}}
\newunicodechar{∩}{\ensuremath{\cap}}
\newunicodechar{≡}{\ensuremath{\equiv}}
\newunicodechar{⊂}{\ensuremath{\subset}}
\newunicodechar{⊥}{\ensuremath{\perp}}
\newunicodechar{∃}{\ensuremath{\exists}}
\newunicodechar{∀}{\ensuremath{\forall}}
\newunicodechar{∛}{\ensuremath{\sqrt[3]{\,}}}
\newunicodechar{∜}{\ensuremath{\sqrt[4]{\,}}}
\newunicodechar{⃗}{\ensuremath{{}^{\to}}}
\newunicodechar{ⁿ}{\ifmmode^{n}\else\textsuperscript{\ensuremath{n}}\fi}
\newunicodechar{ˣ}{\ifmmode^{x}\else\textsuperscript{\ensuremath{x}}\fi}
\newunicodechar{ᵇ}{\ifmmode^{b}\else\textsuperscript{\ensuremath{b}}\fi}
\newunicodechar{ʳ}{\ifmmode^{r}\else\textsuperscript{\ensuremath{r}}\fi}
\newunicodechar{ᵘ}{\ifmmode^{u}\else\textsuperscript{\ensuremath{u}}\fi}
\newunicodechar{ʸ}{\ifmmode^{y}\else\textsuperscript{\ensuremath{y}}\fi}
\newunicodechar{ᵃ}{\ifmmode^{a}\else\textsuperscript{\ensuremath{a}}\fi}
\newunicodechar{ᵉ}{\ifmmode^{e}\else\textsuperscript{\ensuremath{e}}\fi}
\newunicodechar{ᵏ}{\ifmmode^{k}\else\textsuperscript{\ensuremath{k}}\fi}
\newunicodechar{ᵖ}{\ifmmode^{p}\else\textsuperscript{\ensuremath{p}}\fi}
\newunicodechar{ᴺ}{\ifmmode^{N}\else\textsuperscript{\ensuremath{N}}\fi}
\newunicodechar{ᶜ}{\ifmmode^{c}\else\textsuperscript{\ensuremath{c}}\fi}
\newunicodechar{ʰ}{\ifmmode^{h}\else\textsuperscript{\ensuremath{h}}\fi}
\newunicodechar{ᵠ}{\ifmmode^{q}\else\textsuperscript{\ensuremath{q}}\fi}
\newunicodechar{ₙ}{\ifmmode_{n}\else\textsubscript{\ensuremath{n}}\fi}
\newunicodechar{ₐ}{\ifmmode_{a}\else\textsubscript{\ensuremath{a}}\fi}
\newunicodechar{ᵢ}{\ifmmode_{i}\else\textsubscript{\ensuremath{i}}\fi}
\newunicodechar{ᵣ}{\ifmmode_{r}\else\textsubscript{\ensuremath{r}}\fi}
\newunicodechar{ₖ}{\ifmmode_{k}\else\textsubscript{\ensuremath{k}}\fi}
\newunicodechar{ᵦ}{\ifmmode_{\beta}\else\textsubscript{\ensuremath{\beta}}\fi}
\newunicodechar{ₓ}{\ifmmode_{x}\else\textsubscript{\ensuremath{x}}\fi}
\newfontfamily\popdisplay{ArchivoBlack-Regular.ttf}[Path=../../../fonts/]
\newfontfamily\poplogo{SpaceGrotesk-Bold.ttf}[Path=../../../fonts/]
\newfontfamily\popsemi{PlusJakartaSans-SemiBold.ttf}[Path=../../../fonts/]
\newfontfamily\popxbold{PlusJakartaSans-ExtraBold.ttf}[Path=../../../fonts/]
\newfontfamily\popxboldls{PlusJakartaSans-ExtraBold.ttf}[Path=../../../fonts/,LetterSpace=3.0]

\setlist[enumerate]{leftmargin=1.7em,itemsep=5pt,topsep=4pt}
\setlist[itemize]{leftmargin=1.7em,itemsep=4pt,topsep=4pt,label={\color{poporange}\textbullet}}
\setlength{\parindent}{0pt}
\setlength{\parskip}{5pt}
\renewcommand{\arraystretch}{1.25}

% pgfplots : style Pop par défaut
\pgfplotsset{
  every axis/.append style={axis line style={popink,line width=1pt},tick style={popink},
    grid style={popline,line width=0.5pt},label style={font=\small},tick label style={font=\footnotesize}},
  popcurve/.style={poporange,line width=1.8pt,no markers,smooth},
}
% Même style disponible dans un \draw TikZ simple (hors pgfplots)
\tikzset{popcurve/.style={poporange,line width=1.8pt}}
\tikzset{popgrid/.style={popline,very thin}}
\colorlet{popcurve}{poporange} % le nom du style est parfois utilisé comme couleur

% ============================================================
% Pill / squiggle
% ============================================================
\newcommand{\poppill}[3]{%
  \tikz[baseline=-0.55ex]\node[draw=popink,line width=1.2pt,rounded corners=2.6pt,%
    fill=#2,inner xsep=6.5pt,inner ysep=3pt]{%
    \popxboldls\fontsize{7}{7}\selectfont\color{#3}\MakeUppercase{#1}};%
}
\newcommand{\popsquiggle}[1]{%
  \tikz[baseline=-0.4ex]{
    \draw[#1,line width=2.6pt,line cap=round]
      plot [smooth] coordinates {(0,0.09) (0.34,0.01) (0.68,0.21) (1.02,0.01) (1.36,0.10)};
  }%
}
% Mot-clé surligné (marqueur orange sous le texte)
\newcommand{\cle}[1]{{\popxbold\color{popdark}#1}}

% ============================================================
% En-tête / pied
% ============================================================
\pagestyle{fancy}
\fancyhf{}
\renewcommand{\headrulewidth}{0pt}
\renewcommand{\footrulewidth}{0pt}
\lhead{\raisebox{-0.15ex}{\poplogo\fontsize{13}{13}\selectfont\color{popink}OnBuch\color{poporange}+}}
\rhead{\poppill{\DOCMATIERE\ \textbullet\ \DOCNIVEAU\ \textbullet\ Leçon \DOCLECON}{white}{popink}}
\lfoot{\popxbold\fontsize{7.6}{7.6}\selectfont\color{popmuted}\MakeUppercase{Cours signé OnBuch+}\ \textbullet\ {\popsemi\DOCTITRE}}
\rfoot{\tikz[baseline=-0.6ex]\node[rounded corners=8.5pt,fill=popink,minimum height=17pt,minimum width=17pt,inner xsep=5pt,inner sep=0]{\popxbold\fontsize{8}{8}\selectfont\color{popcream}\thepage\,/\,\pageref*{LastPage}};}

% ============================================================
% Titres
% ============================================================
\newcommand{\chaptitre}[2]{%
  \begin{tcolorbox}[enhanced,colback=poporange,colframe=popink,boxrule=2.6pt,arc=11pt,
      drop shadow=popink,left=15pt,right=15pt,top=12pt,bottom=11pt,before skip=3pt,after skip=18pt]
    \poppill{#2}{popink}{popcream}\par\vspace{9pt}
    {\raggedright\hyphenpenalty=10000\exhyphenpenalty=10000\popdisplay\fontsize{22}{24}\selectfont\color{popcream}\MakeUppercase{#1}\par}
  \end{tcolorbox}
}
% Section numérotée : pastille orange avec le numéro + titre Archivo
\newcounter{popsec}
\newcommand{\coursec}[1]{%
  \stepcounter{popsec}\setcounter{popsub}{0}%
  \par\vspace{16pt}\noindent
  \begin{minipage}{\linewidth}
  \tikz[baseline=-0.7ex]\node[draw=popink,line width=1.6pt,fill=poporange,rounded corners=4pt,minimum size=22pt,inner sep=0]{\popdisplay\fontsize{12}{12}\selectfont\color{popcream}\thepopsec};%
  \hspace{8pt}{\popdisplay\fontsize{15}{17}\selectfont\color{popink}\MakeUppercase{#1}}\par
  \vspace{4pt}\noindent{\color{poporange}\rule{36pt}{3.6pt}}
  \end{minipage}\par\nopagebreak\vspace{8pt}
}
\newcounter{popsub}
\newcommand{\courssub}[1]{%
  \stepcounter{popsub}%
  \par\vspace{9pt}\noindent{\popxbold\fontsize{11.5}{13}\selectfont\color{popink}{\color{poporange}\thepopsec.\thepopsub}\hspace{6pt}#1}\par\nopagebreak\vspace{4pt}
}

% ============================================================
% Boîtes pédagogiques — même recette : fond blanc, bordure épaisse colorée,
% ombre dure encre, badge plein. Titre optionnel [#1].
% ============================================================
\tcbset{popbase/.style={breakable,enhanced,colback=white,boxrule=2.3pt,arc=9pt,
  shadow={3pt}{-3pt}{0pt}{popink},left=14pt,right=14pt,top=11pt,bottom=11pt,before skip=11pt,after skip=15pt,
  fontupper=\popsemi\fontsize{10.6}{14.5}\selectfont\color{popink},
  fontlower=\popsemi\fontsize{10.6}{14.5}\selectfont\color{popink}}}
% Coche dessinée (le glyphe ✓ n'existe pas dans les polices du document)
\newcommand{\popcheck}[1]{\tikz[baseline=-0.5ex]\draw[#1,line width=1.6pt,line cap=round,line join=round](-0.13,0.0)--(-0.04,-0.09)--(0.13,0.1);}
\providecommand{\coche}{\popcheck{popgreen}}
\providecommand{\resultat}[1]{\par\smallskip\noindent\fcolorbox{popgreen}{popgreenL}{\parbox{\dimexpr\linewidth-2\fboxsep-2\fboxrule\relax}{\textbf{Résultat.} #1}}\par\smallskip}
\newcommand{\popboxhead}[3]{\poppill{#1}{#2}{white}\ifx\relax#3\relax\else\hspace{6pt}{\popxbold\fontsize{10.6}{12}\selectfont\color{popink}#3}\fi\par\vspace{7pt}}

\newtcolorbox{aretenir}[1][]{popbase,colframe=popgreen,before upper={\popboxhead{À retenir}{popgreen}{#1}}}
% Texte en anglais (document de lecture, dialogue, extrait) : cadre
% d'étude. Un titre optionnel (source, auteur) s'affiche dans l'en-tête.
\newtcolorbox{textebox}[1][]{popbase,colframe=popblue,before upper={\popboxhead{Text}{popblue}{#1}\begin{otherlanguage}{english}},after upper={\end{otherlanguage}}}
% Marques ✓ / ✗ (forme correcte / fautive) souvent utilisées par les modèles
\providecommand{\cross}{\textcolor{poppink}{\ensuremath{\times}}}
\providecommand{\xmark}{\cross}\providecommand{\crossmark}{\cross}\providecommand{\cmark}{\ensuremath{\checkmark}}
% Ligne de réponse / justification à compléter (inventée par les modèles)
\providecommand{\justification}[1]{\par\noindent\textit{Justification~:} \dotfill\par}
% \textipa (package tipa non chargé) : la police ne couvre pas l'API -> simple passage
\providecommand{\textipa}[1]{#1}
% Soulignement (ulem non chargé) : \uline, \ul
\providecommand{\uline}[1]{\underline{#1}}\providecommand{\ul}[1]{\underline{#1}}
% \glossaire{\item ...} inventée par les modèles : liste de glossaire
\providecommand{\glossaire}[1]{\begin{itemize}#1\end{itemize}}
% \alert (beamer) inventée par les modèles : texte mis en évidence
\providecommand{\alert}[1]{\textbf{\textcolor{poppink}{#1}}}
% variantes ASCII d'ordinaux écrites en macro
\providecommand{\iere}{\textsuperscript{re}}\providecommand{\ieme}{\textsuperscript{e}}\providecommand{\ere}{\textsuperscript{re}}\providecommand{\eme}{\textsuperscript{e}}\providecommand{\er}{\textsuperscript{er}}\providecommand{\ie}{\textsuperscript{e}}
% Ordinaux français écrits en macro par les modèles : 1\ière, 3\ième
\newcommand{\ière}{\textsuperscript{re}}\newcommand{\ième}{\textsuperscript{e}}
% \exemple (inventée par les modèles) : étiquette « Exemple : »
\providecommand{\exemple}{\textbf{Exemple~:}\ }
% \square utilisable aussi hors mode math (cases à cocher)
\AtBeginDocument{\let\squareMath\square\renewcommand{\square}{\ensuremath{\squareMath}}}
% Mot ou phrase en anglais à l'intérieur d'un paragraphe français
\newcommand{\eng}[1]{\ifmmode\text{\textit{#1}}\else\begingroup\selectlanguage{english}\itshape #1\endgroup\fi}
\let\esp\eng % alias toléré
% flèches d'intonation inventées par les modèles
\providecommand{\textsearrow}{}\renewcommand{\textsearrow}{\ensuremath{\searrow}}\providecommand{\textnearrow}{}\renewcommand{\textnearrow}{\ensuremath{\nearrow}}
% \fr{..} : mot français (inventé par les modèles)
\providecommand{\fr}[1]{\textit{#1}}
\newtcolorbox{exemplebox}[1][]{popbase,colframe=poporange,before upper={\popboxhead{Exemple}{poporange}{#1}}}
\newtcolorbox{definition}[1][]{popbase,colframe=popblue,before upper={\popboxhead{Définition}{popblue}{#1}}}
\newtcolorbox{propriete}[1][]{popbase,colframe=poppurple,before upper={\popboxhead{Propriété}{poppurple}{#1}}}
\newtcolorbox{methode}[1][]{popbase,colframe=popgold,before upper={\popboxhead{Méthode}{popgold}{#1}}}
\newtcolorbox{attention}[1][]{popbase,colframe=poppink,before upper={\popboxhead{Attention}{poppink}{#1}}}
\newtcolorbox{experience}[1][]{popbase,colframe=popblue,colback=popblueL,before upper={\popboxhead{Expérience}{popblue}{#1}}}
% « Le savais-tu ? » : carte violette pleine (contexte, culture, Cameroun)
\newtcolorbox{savaistu}[1][]{popbase,colback=poppurple,colframe=popink,
  fontupper=\popsemi\fontsize{10.4}{14.2}\selectfont\color{white},
  before upper={\poppill{Le savais-tu\,?}{white}{poppurple}\ifx\relax#1\relax\else\hspace{6pt}{\popxbold\color{white}#1}\fi\par\vspace{7pt}}}
% Exercice résolu : énoncé en haut, \tcblower puis la solution sur fond crème
\newcounter{popexo}\newcounter{popexr}
\newtcolorbox{exoresolu}[1][]{popbase,colframe=poporange,colbacklower=popcream,
  segmentation style={popink,line width=1.2pt,dashed},
  before upper={\stepcounter{popexr}\popboxhead{Exercice résolu \thepopexr}{poporange}{#1}},
  before lower={\poppill{Solution}{popink}{popcream}\par\vspace{6pt}}}
% Exercice (non résolu en place) — la correction est dans la partie Corrigés
\newtcolorbox{exercice}[1][]{popbase,colframe=popink,
  before upper={\stepcounter{popexo}\popboxhead{Exercice \thepopexo}{popink}{#1}}}
% Corrigé
\newtcolorbox{corrige}[1][]{popbase,colframe=popgreen,colback=popgreenL,
  before upper={\popboxhead{Corrigé}{popgreen}{#1}}}
% Objectifs / prérequis / bilan
\newtcolorbox{objectifs}{popbase,colframe=popink,colback=poporangeL,
  before upper={\popboxhead{Objectifs de la leçon}{popink}{}}}
\newtcolorbox{prerequis}{popbase,colframe=popink,colback=popblueL,
  before upper={\popboxhead{Prérequis}{popblue}{}}}
\newtcolorbox{bilan}{popbase,colframe=popink,colback=white,boxrule=2.8pt,
  before upper={\popboxhead{Fiche bilan}{poporange}{L'essentiel à connaître pour l'examen}}}
% Figure encadrée avec légende
\newtcolorbox{popfigure}[1][]{enhanced,breakable=false,colback=white,colframe=popline,boxrule=1.4pt,arc=8pt,
  left=10pt,right=10pt,top=10pt,bottom=8pt,before skip=10pt,after skip=12pt,halign=center,
  lower separated=false,halign lower=center,
  fontlower=\popsemi\fontsize{9}{11.5}\selectfont\color{popmuted},
  #1}
\newcommand{\legende}[1]{\par\vspace{6pt}{\popsemi\fontsize{9}{11.5}\selectfont\color{popmuted}#1}}

% ============================================================
% Signature OnBuch+
% ============================================================
\newcommand{\poplogoinline}[1]{{\poplogo\fontsize{#1}{#1}\selectfont\color{popink}OnBuch\color{poporange}+}}
\newcommand{\signatureonbuch}{%
  \begin{tcolorbox}[enhanced,colback=popink,colframe=popink,boxrule=2.2pt,arc=10pt,
      drop shadow=poporange,left=14pt,right=14pt,top=10pt,bottom=10pt,before skip=0pt,after skip=0pt]
    \tikz[baseline=-0.7ex]\node[circle,fill=popgreen,draw=popcream,line width=1.3pt,minimum size=22pt,inner sep=0]{\popcheck{white}};%
    \hspace{9pt}\begin{minipage}[c]{0.62\linewidth}
      {\popxbold\fontsize{12}{13}\selectfont\color{popcream}Signé\ \textbullet\ {\poplogo OnBuch\color{poporange}+}}\par\vspace{2pt}
      {\raggedright\popsemi\fontsize{8}{10.5}\selectfont\color{popline}\MakeUppercase{Équipe pédagogique OnBuch+}\\\MakeUppercase{Conforme au programme officiel MINESEC}\par}
    \end{minipage}\hfill
    {\poplogo\fontsize{22}{22}\selectfont\color{popcream}OnBuch\color{poporange}+}
  \end{tcolorbox}%
}

% ============================================================
% Page de garde
% #1 titre · #2 matière · #3 niveau · #4 module · #5 n° leçon · #6 durée ·
% #7 description / objectifs (texte ou itemize)
% ============================================================
\newcommand{\pagegarde}[7]{%
  \thispagestyle{empty}
  \newgeometry{a4paper,margin=1.7cm,top=1.6cm,bottom=1.4cm}%
  \begin{tikzpicture}[remember picture,overlay]
    \fill[poporange!18] ([xshift=-3.2cm,yshift=-3.6cm]current page.north east) circle (4.6cm);
    \fill[poppurple!14] ([xshift=2.4cm,yshift=7.6cm]current page.south west) circle (3.8cm);
  \end{tikzpicture}%
  {\poplogo\fontsize{30}{30}\selectfont\color{popink}OnBuch\color{poporange}+}\hfill
  \raisebox{0.6ex}{\poppill{Cours signé \textbullet\ Premium}{popink}{popcream}}\par
  \vspace{1pt}\popsquiggle{poppurple}\par
  \vspace{8pt}
  \begin{tcolorbox}[enhanced,colback=poporange,colframe=popink,boxrule=2.8pt,arc=13pt,
      drop shadow=popink,left=18pt,right=18pt,top=12pt,bottom=13pt,after skip=10pt]
    \poppill{#2 \textbullet\ #3}{popink}{popcream}\hspace{5pt}\poppill{#4}{white}{popink}\par\vspace{8pt}
    {\popdisplay\fontsize{14}{14}\selectfont\color{popink}LEÇON #5}\par\vspace{4pt}
    {\raggedright\hyphenpenalty=10000\exhyphenpenalty=10000\popdisplay\fontsize{25}{27}\selectfont\color{popcream}\MakeUppercase{#1}\par}
    \vspace{8pt}
    {\popxbold\fontsize{9.5}{9.5}\selectfont\color{popink}\MakeUppercase{Durée conseillée : #6}}
  \end{tcolorbox}
  \begin{tcolorbox}[enhanced,colback=white,colframe=popink,boxrule=2.2pt,arc=10pt,
      drop shadow=popink,left=16pt,right=16pt,top=10pt,bottom=10pt,after skip=8pt]
    \poppill{Dans cette leçon}{poppurple}{white}\par\vspace{5pt}
    \popsemi\fontsize{9.3}{12.6}\selectfont\color{popink}#7
  \end{tcolorbox}
  \noindent\begin{minipage}[t]{0.487\linewidth}
    \begin{tcolorbox}[enhanced,colback=popgreen,colframe=popink,boxrule=2.2pt,arc=10pt,
        drop shadow=popink,left=11pt,right=11pt,top=10pt,bottom=10pt,height=3.1cm,valign=center]
      \tcbox[colback=white,colframe=popink,boxrule=1.3pt,arc=5pt,boxsep=0pt,nobeforeafter,
        left=3pt,right=3pt,top=3pt,bottom=3pt,valign=center]{\qrcode[height=1.9cm]{https://play.google.com/store/apps/details?id=com.luvvix.onbuch&hl=fr}}%
      \hspace{8pt}\begin{minipage}[c]{0.5\linewidth}
        {\popdisplay\fontsize{11}{12.5}\selectfont\color{white}\MakeUppercase{Télécharge OnBuch}}\par\vspace{4pt}
        {\raggedright\popsemi\fontsize{8.4}{11}\selectfont\color{white}Gratuit \textbullet\ Google Play\\Tuteur IA, annales, concours, résultats\par}
      \end{minipage}
    \end{tcolorbox}
  \end{minipage}\hfill
  \begin{minipage}[t]{0.487\linewidth}
    \begin{tcolorbox}[enhanced,colback=poppurple,colframe=popink,boxrule=2.2pt,arc=10pt,
        drop shadow=popink,left=11pt,right=11pt,top=10pt,bottom=10pt,height=3.1cm,valign=center]
      \tikz[baseline=-0.7ex]\node[circle,fill=white,draw=popink,line width=1.3pt,minimum size=1.6cm,inner sep=0]{\popdisplay\fontsize{24}{24}\selectfont\color{poppurple}+};%
      \hspace{8pt}\begin{minipage}[c]{0.6\linewidth}
        {\popdisplay\fontsize{11}{12.5}\selectfont\color{white}\MakeUppercase{Passe à OnBuch+}}\par\vspace{4pt}
        {\raggedright\popsemi\fontsize{8.4}{11}\selectfont\color{white}Tous les cours signés, Tuteur IA illimité, fiches PDF — onglet OnBuch+ de l'app\par}
      \end{minipage}
    \end{tcolorbox}
  \end{minipage}
  \vspace{8pt}
  \popcontenu
  \vfill
  \signatureonbuch
  \restoregeometry
  \newpage
}

% Rangée « Ce cours contient » (page de garde)
\newcommand{\poptile}[3]{% #1 couleur #2 gros texte #3 légende
  \begin{tcolorbox}[enhanced,colback=#1,colframe=popink,boxrule=2pt,arc=9pt,shadow={2.5pt}{-2.5pt}{0pt}{popink},
      left=6pt,right=6pt,top=9pt,bottom=9pt,halign=center,valign=center,height=1.75cm,width=\linewidth,nobeforeafter]
    {\popdisplay\fontsize{12.5}{13}\selectfont\color{white}\MakeUppercase{#2}}\par\vspace{3pt}
    {\centering\popsemi\fontsize{7.8}{9.5}\selectfont\color{white}#3\par}
  \end{tcolorbox}}
\newcommand{\popcontenu}{%
  \par\noindent{\popxboldls\fontsize{7.5}{7.5}\selectfont\color{popmuted}CE COURS SIGNÉ CONTIENT}\par\vspace{7pt}
  \noindent\begin{minipage}[t]{0.235\linewidth}\poptile{popblue}{Cours}{complet, illustré, pas à pas}\end{minipage}\hfill
  \begin{minipage}[t]{0.235\linewidth}\poptile{poporange}{Méthodes}{et exercices résolus}\end{minipage}\hfill
  \begin{minipage}[t]{0.235\linewidth}\poptile{poppink}{Exercices}{type examen + corrigés}\end{minipage}\hfill
  \begin{minipage}[t]{0.235\linewidth}\poptile{popgreen}{Fiche bilan}{l'essentiel à retenir}\end{minipage}\par}

% Sommaire
\newtcolorbox{sommaire}{enhanced,colback=white,colframe=popink,boxrule=2.2pt,arc=10pt,
  drop shadow=popink,left=18pt,right=18pt,top=13pt,bottom=13pt,before skip=6pt,after skip=20pt,
  before upper={\poppill{Sommaire}{poppurple}{white}\par\vspace{9pt}\popsemi\fontsize{10.6}{14.5}\selectfont\color{popink}}}

% Signature de fin de cours — poussée tout en bas de la dernière page
\newcommand{\findecours}{%
  \par\vfill\nopagebreak
  \begin{center}\popsquiggle{poporange}\end{center}\vspace{4pt}
  \signatureonbuch
}
\AtBeginDocument{\def\llbracket{\mathopen{[\mkern-3.5mu[}}\def\rrbracket{\mathclose{]\mkern-3.5mu]}}} % intervalles d'entiers (glyphe ⟦ absent de la police)
% Glyphes absents de Fira Math : redéfinis à partir de glyphes disponibles
\AtBeginDocument{%
  \def\setminus{\mathbin{\backslash}}\let\smallsetminus\setminus
  \def\vdots{\mathord{\vbox{\baselineskip3.2pt\lineskiplimit0pt\kern4pt\hbox{.}\hbox{.}\hbox{.}}}}%
  \def\ddots{\mathinner{\mkern1mu\raise6.5pt\hbox{.}\mkern2mu\raise3.6pt\hbox{.}\mkern2mu\raise0.7pt\hbox{.}\mkern1mu}}%
  \def\triangle{\mathord{\scalebox{0.9}{$\symup{\Delta}$}}}%
  \def\nabla{\mathord{\raisebox{\depth}{\scalebox{1}[-1]{$\symup{\Delta}$}}}}%
  \def\checkmark{\popcheck{popdark}}%
  \def\star{\mathbin{\ast}}\def\bigstar{\mathord{\ast}}%
  \def\diamond{\mathbin{\scalebox{0.6}{$\Diamond$}}}%
  \def\bigcup{\mathop{\mathchoice{\raisebox{-0.3em}{\scalebox{1.7}{$\cup$}}}{\scalebox{1.25}{$\cup$}}{\cup}{\cup}}\displaylimits}%
  \def\bigcap{\mathop{\mathchoice{\raisebox{-0.3em}{\scalebox{1.7}{$\cap$}}}{\scalebox{1.25}{$\cap$}}{\cap}{\cap}}\displaylimits}%
  \def\lesseqqgtr{\mathrel{\lessgtr}}\def\bigoplus{\mathop{\oplus}\displaylimits}%
}
\providecommand{\ssi}{si et seulement si} % abréviation fréquente
\AtBeginDocument{\providecommand{\wideparen}{\overparen}} % arc de cercle

%% FIGFIX-BEGIN v3 — correctifs globaux des figures (docs/onbuch-generation/tools/figfix)
\usepackage{adjustbox}\usepackage{etoolbox}
% toute figure TikZ/pgfplots plus large que la ligne est réduite à la largeur disponible
\BeforeBeginEnvironment{tikzpicture}{\begin{adjustbox}{max width=\linewidth}}
\AfterEndEnvironment{tikzpicture}{\end{adjustbox}}
% légendes pgfplots lisibles par-dessus les courbes
\pgfplotsset{every axis/.append style={legend style={fill=white,fill opacity=0.92,text opacity=1,draw=black!15,inner sep=3pt}}}
% étiquettes d'arêtes (syntaxe edge["..."]) avec fond blanc
\tikzset{every edge quotes/.append style={fill=white,inner sep=1.2pt}}
% bulles de cartes mentales : texte à la ligne dans la bulle, bulle un peu plus grande
\tikzset{every concept/.append style={text width=2.0cm,align=center,minimum size=2.3cm,inner sep=2pt}}
%% FIGFIX-END
\begin{document}

\pagegarde{Lesson 10 — Grammar: Phrasal verbs and prepositional phrases}{Anglais}{Tle A}{Chapter 1 : Family and Social Life}{EN10}{2 h}{Cette leçon de grammaire présente les phrasal verbs (verbes à particule) et les prepositional phrases (groupes prépositionnels), structures essentielles de l'anglais authentique. À travers l'observation d'exemples, des règles claires, une comparaison avec le français et des exercices progressifs, les élèves apprennent à les identifier, les construire et les utiliser correctement à l'écrit comme à l'oral.
\vspace{4pt}
\begin{multicols}{2}\small
\begin{itemize}
\item À la fin de cette leçon, je sais définir un phrasal verb et le distinguer d'un verbe suivi d'une préposition ordinaire
\item À la fin de cette leçon, je sais classer les phrasal verbs en quatre types : intransitifs, transitifs séparables, transitifs inséparables et verbes à deux particules
\item À la fin de cette leçon, je sais placer correctement le complément (nom ou pronom) avec les phrasal verbs séparables
\item À la fin de cette leçon, je sais identifier et construire les prepositional phrases (préposition + groupe nominal) et les constructions verbe + préposition
\item À la fin de cette leçon, je sais éviter les erreurs fréquentes des francophones (traduction mot à mot, mauvais placement du pronom, préposition calquée du français)
\item À la fin de cette leçon, je sais employer les phrasal verbs courants liés à la vie familiale et sociale dans des phrases et de courts textes
\end{itemize}\end{multicols}}

\begin{sommaire}
\begin{multicols}{2}
\begin{enumerate}
\item Observation et découverte
\item Les quatre types de phrasal verbs
\item Phrasal verbs fréquents : vie familiale et sociale
\item Les prepositional phrases
\item Comparaison français/anglais et pièges
\item Entraînement guidé et réinvestissement
\item Activité d'intégration
\item Méthodes et exercices résolus
\item Exercices
\item Corrigés des exercices
\item Fiche bilan
\end{enumerate}
\end{multicols}
\end{sommaire}

\begin{objectifs}
\begin{itemize}
\item À la fin de cette leçon, je sais définir un phrasal verb et le distinguer d'un verbe suivi d'une préposition ordinaire
\item À la fin de cette leçon, je sais classer les phrasal verbs en quatre types : intransitifs, transitifs séparables, transitifs inséparables et verbes à deux particules
\item À la fin de cette leçon, je sais placer correctement le complément (nom ou pronom) avec les phrasal verbs séparables
\item À la fin de cette leçon, je sais identifier et construire les prepositional phrases (préposition + groupe nominal) et les constructions verbe + préposition
\item À la fin de cette leçon, je sais éviter les erreurs fréquentes des francophones (traduction mot à mot, mauvais placement du pronom, préposition calquée du français)
\item À la fin de cette leçon, je sais employer les phrasal verbs courants liés à la vie familiale et sociale dans des phrases et de courts textes
\item À la fin de cette leçon, je sais deviner le sens d'un phrasal verb inconnu grâce au contexte
\item À la fin de cette leçon, je sais transformer des phrases en remplaçant un verbe simple par un phrasal verb équivalent
\end{itemize}
\end{objectifs}

\begin{prerequis}
\begin{itemize}
\item Les prépositions de base en anglais (in, on, at, for, with, up, down...) vues au collège et en Seconde
\item La structure de la phrase anglaise : sujet + verbe + complément (Première)
\item Les pronoms compléments (me, him, her, it, us, them)
\item Le vocabulaire de la vie familiale et sociale (Chapter 1, leçons précédentes)
\item La distinction verbe transitif / verbe intransitif
\end{itemize}
\end{prerequis}

\newpage

\coursec{Observation et découverte}

\courssub{Une petite erreur qui change tout}

Pendant les vacances, Divine, élève de Terminale A à Bafoussam, discute sur Internet avec son correspondant britannique, Tom. Elle veut expliquer qu'elle garde ses petits frères pendant que sa mère est au marché. Fière de son anglais, elle écrit :

\eng{I will look the children after while my mother is at the market.}

Tom répond aussitôt, amusé :

\eng{You mean you will look AFTER the children! And careful: don't look FOR them, don't look them UP, just look AFTER them!}

Divine est perplexe. Elle connaît le verbe \eng{to look} (« regarder »). Alors pourquoi Tom insiste-t-il sur ce petit mot \eng{after} ? Et que signifient \eng{look for} et \eng{look up} ? Elle réalise que ces petits mots qui suivent les verbes anglais — \eng{up, after, for, off, out, back} — changent complètement le sens de la phrase. Sans eux, impossible de comprendre une conversation anglaise authentique ; mal placés, ils produisent des contresens amusants... ou embarrassants.

\textbf{Problématique :}\par
Comment maîtriser ces \cle{phrasal verbs} (verbes à particule) et ces \cle{prepositional phrases} (groupes prépositionnels) qui font toute la richesse — et la difficulté — de l'anglais parlé et écrit ?

\courssub{Lecture découverte : un matin dans une famille camerounaise}

Lisons d'abord ce court texte qui décrit le quotidien de la famille Ndi, à Bamenda. Lis-le une première fois pour le sens général, puis une seconde fois en soulignant tous les verbes suivis d'un petit mot.

\begin{textebox}[A morning with the Ndi family]
Every morning, Mrs Ndi gets up at 5 a.m., wakes the children up and helps them put on their uniforms. Her husband turns off the generator before they leave for school. Mrs Ndi looks after her elderly mother too: she cooks for her and takes care of the house.

Life has not always been easy. When her son dropped out of school two years ago, she never gave up on him. She encouraged him to go back to his studies, and last month he finally came back home with his GCE certificate. The whole family celebrated.

“I grew up in a small village near Bafoussam,” Mrs Ndi says, “and my parents always told me: never put off until tomorrow what you can do today. That is why I get on with my work every single day.”
\end{textebox}

\textbf{Glossaire :}\par

\begin{center}
\begin{tabularx}{\linewidth}{|l|l|X|}
\hline
\rowcolor{popblueL} \textbf{Mot anglais} & \textbf{Nature} & \textbf{Traduction} \\
\hline
elderly & adjectif & âgé(e) \\
\hline
generator & nom & groupe électrogène \\
\hline
certificate & nom & diplôme, certificat \\
\hline
to encourage & verbe & encourager \\
\hline
whole & adjectif & entier, tout entier \\
\hline
\end{tabularx}
\end{center}

\textbf{Questions de compréhension :}\par

\begin{enumerate}
\item \eng{What time does Mrs Ndi get up?} \emph{(À quelle heure Mme Ndi se lève-t-elle ?)}
\item \eng{Who does she look after?} \emph{(De qui s'occupe-t-elle ?)}
\item \eng{What happened to her son two years ago?} \emph{(Qu'est-il arrivé à son fils il y a deux ans ?)}
\item \eng{What advice did her parents give her?} \emph{(Quel conseil ses parents lui ont-ils donné ?)}
\end{enumerate}

\courssub{Repérage : des verbes accompagnés de particules}

Reprenons le texte et soulignons les verbes suivis d'un petit mot :

\begin{itemize}
\item \eng{gets \textbf{up}} — se lève
\item \eng{wakes the children \textbf{up}} — réveille les enfants
\item \eng{put \textbf{on} their uniforms} — enfiler leurs uniformes
\item \eng{turns \textbf{off} the generator} — éteint le groupe électrogène
\item \eng{looks \textbf{after} her mother} — s'occupe de sa mère
\item \eng{dropped \textbf{out} of school} — a abandonné l'école
\item \eng{gave \textbf{up} on him} — a cessé de croire en lui
\item \eng{came \textbf{back} home} — est rentré à la maison
\item \eng{grew \textbf{up} in a village} — a grandi dans un village
\item \eng{put \textbf{off} until tomorrow} — remettre au lendemain
\item \eng{get \textbf{on} with my work} — poursuivre mon travail
\end{itemize}

\begin{definition}[Phrasal verb — verbe à particule]
Un \cle{phrasal verb} (verbe à particule, ou verbe à postposition) est un verbe suivi d'un petit mot appelé \cle{particule} — un adverbe (\eng{up, out, off, back, on, away}) ou une préposition (\eng{after, for, at, into}) — dont le \textbf{sens global diffère du sens du verbe simple}. Le verbe et sa particule forment une seule unité de sens, comme un seul mot.
\end{definition}

\courssub{Question guidée : la particule change-t-elle le sens ?}

Comparons attentivement :

\begin{center}
\begin{tabularx}{\linewidth}{|l|X|X|}
\hline
\rowcolor{popblueL} \textbf{Verbe} & \textbf{Exemple anglais} & \textbf{Traduction} \\
\hline
to look & She is looking at the picture. & Elle regarde la photo. \\
\hline
to look after & She looks after her mother. & Elle s'occupe de sa mère. \\
\hline
to look for & I am looking for my keys. & Je cherche mes clés. \\
\hline
to look up & Look up the word in the dictionary. & Cherche le mot dans le dictionnaire. \\
\hline
to look forward to & I look forward to the holidays. & J'attends les vacances avec impatience. \\
\hline
\end{tabularx}
\end{center}

Le constat est clair : \textbf{la particule transforme le sens du verbe}. \eng{To look} signifie « regarder », mais \eng{to look after} signifie « s'occuper de » : on ne peut pas deviner ce sens à partir de « regarder ». C'est ce qu'on appelle un sens \cle{idiomatique} : le phrasal verb se mémorise comme un mot de vocabulaire à part entière, avec sa traduction.

\begin{popfigure}
\begin{tikzpicture}[mindmap, grow cyclic, every node/.style={concept}, root concept/.append style={concept color=popblue, font=\Large\bfseries}, level 1 concept/.append style={level distance=3.6cm, sibling angle=90, font=\small}, text=white]
\node[root concept]{LOOK}
  child[concept color=poporange] { node[align=center]{look after\\ s'occuper de} }
  child[concept color=popgreen] { node[align=center]{look for\\ chercher} }
  child[concept color=poppurple] { node[align=center]{look up\\ consulter} }
  child[concept color=poppink] { node[align=center]{look forward to\\ attendre avec\\ impatience} };
\end{tikzpicture}
\legende{Carte mentale : un même verbe, quatre sens différents selon la particule.}
\end{popfigure}

\courssub{Phrasal verb ou verbe + préposition ordinaire ?}

Attention : tout verbe suivi d'une préposition n'est pas un phrasal verb. Il faut distinguer deux cas.

\begin{propriete}[Le test du sens]
\begin{itemize}
\item Si le sens est \textbf{littéral} (le verbe garde son sens habituel et la préposition garde son sens habituel), c'est un simple \textbf{verbe + préposition} : \eng{She is looking \textbf{at} the picture.} « Elle regarde la photo. » — regarder + vers : sens normal.
\item Si le sens est \textbf{nouveau, imprévisible}, c'est un \textbf{phrasal verb} : \eng{She looks \textbf{after} the children.} « Elle s'occupe des enfants » — rien à voir avec « regarder après » !
\end{itemize}
\end{propriete}

\begin{exemplebox}[Sens littéral ou sens idiomatique ?]
\eng{He ran \textbf{across} the road.} = Il a traversé la route en courant. \emph{(littéral : courir + à travers)}

\eng{I ran \textbf{into} an old friend in Douala.} = Je suis tombé sur un vieil ami à Douala. \emph{(idiomatique : rencontrer par hasard)}

\eng{The plane took \textbf{off} at noon.} = L'avion a décollé à midi. \emph{(idiomatique)}

\eng{Take your feet \textbf{off} the table!} = Enlève tes pieds de la table ! \emph{(littéral)}
\end{exemplebox}

\begin{aretenir}[Le constat essentiel]
Le sens d'un phrasal verb est souvent \textbf{imprévisible} à partir du verbe seul :
\eng{to give up} = abandonner, renoncer (et non « donner en haut ») ;
\eng{to turn up} = arriver, apparaître ;
\eng{to put up with} = supporter, tolérer.
Chaque phrasal verb s'apprend donc comme une entrée de vocabulaire : verbe + particule + traduction + exemple.
\end{aretenir}

\courssub{Premiers exemples traduits}

\begin{exemplebox}[Phrasal verbs de la vie familiale et sociale]
\eng{She looks after her mother.} = Elle s'occupe de sa mère.

\eng{He gave up smoking.} = Il a arrêté de fumer.

\eng{They came back late.} = Ils sont rentrés tard.

\eng{The children get up early on school days.} = Les enfants se lèvent tôt les jours d'école.

\eng{Please turn off the lights before leaving.} = Éteins les lumières avant de partir, s'il te plaît.

\eng{I grew up in Yaoundé.} = J'ai grandi à Yaoundé.
\end{exemplebox}

\begin{savaistu}[Pourquoi tant de phrasal verbs en anglais ?]
Les anglophones utilisent les phrasal verbs \textbf{constamment à l'oral} et dans l'écrit informel. Le verbe « savant » équivalent est presque toujours d'origine latine ou française : \eng{put off} = \emph{postpone} (reporter), \eng{find out} = \emph{discover} (découvrir), \eng{give up} = \emph{abandon} (abandonner), \eng{go on} = \emph{continue} (continuer). C'est un héritage de l'histoire : après 1066, l'anglais populaire germanique a coexisté avec le français des Normands. Résultat : un texte formel (dissertation, lettre officielle) préférera souvent le verbe latin, tandis qu'une conversation préférera le phrasal verb. Connaître les deux registres est la marque d'un excellent anglophone !
\end{savaistu}

\begin{exoresolu}[Identifier les phrasal verbs]
\textbf{Énoncé.} Dans les phrases suivantes, souligne le phrasal verb et dis si son sens est littéral ou idiomatique, puis traduis la phrase.

1. \eng{My little brother put on his shoes and ran outside.}

2. \eng{We are looking at the photographs of the wedding.}

3. \eng{The meeting was put off because of the rain.}

\tcblower
\textbf{Solution détaillée.}

1. \eng{put \textbf{on}} : sens idiomatique (« enfiler, mettre » un vêtement). Traduction : « Mon petit frère a mis ses chaussures et a couru dehors. »

2. \eng{looking \textbf{at}} : sens littéral (regarder + vers) ; ce n'est donc pas un vrai phrasal verb, mais un verbe suivi d'une préposition ordinaire. Traduction : « Nous regardons les photos du mariage. »

3. \eng{put \textbf{off}} : sens idiomatique (« reporter »). Traduction : « La réunion a été reportée à cause de la pluie. »

\textbf{Méthode :} pour trancher, demande-toi toujours si le verbe garde son sens de base. Si oui, c'est une préposition ordinaire ; si non, c'est un phrasal verb à apprendre par cœur.
\end{exoresolu}

\begin{attention}[Pièges fréquents des francophones]
\begin{itemize}
\item \textbf{Ne traduis jamais mot à mot.} \eng{to look after} ne veut pas dire « regarder après » ; \eng{to give up} ne veut pas dire « donner vers le haut ».
\item \textbf{Place de la particule (pronoms).} Avec les phrasal verbs \textbf{séparables} (ex. \eng{wake up}, \eng{turn off}, \eng{put on}), l'objet nominal peut se placer avant ou après la particule : \eng{She wakes the children up} = \eng{She wakes up the children}. Mais si l'objet est un \textbf{pronom}, il \textbf{doit} se placer entre le verbe et la particule : \eng{She wakes \textbf{them} up} (et non *\eng{She wakes up them}). 
\emph{Attention :} \eng{look after} (s'occuper de) est \textbf{inséparable} : l'objet suit toujours la particule (\eng{look after the children}, \eng{look after them}). Nous verrons la classification complète (4 types) dans la partie suivante.
\item \textbf{Faux ami.} \eng{to assist} signifie « aider », pas « assister à » ; pour « assister à une réunion », on dit \eng{to attend a meeting}.
\item \textbf{N'invente pas de phrasal verbs.} On ne peut pas combiner n'importe quel verbe avec n'importe quelle particule : chaque combinaison doit être vérifiée dans un dictionnaire.
\end{itemize}
\end{attention}

\begin{exercice}[À toi de jouer]
Retrouve dans le texte de la famille Ndi cinq phrasal verbs, note leur particule et propose une traduction française pour chacun. Puis réponds à la question de Divine : pourquoi sa phrase \eng{I will look the children after} a-t-elle fait sourire Tom ?
\end{exercice}

Dans la prochaine partie, nous classerons ces phrasal verbs en quatre types grammaticaux précis, afin de savoir exactement où placer la particule et le complément — et de ne plus jamais faire sourire Tom !

\coursec{Les quatre types de phrasal verbs}

Dans la section précédente, nous avons observé, à travers le texte sur la vie quotidienne d'une famille camerounaise, que certains verbes anglais se composent d'un verbe suivi d'une petite particule (\eng{get up}, \eng{look after}, \eng{put on}). Nous avons aussi remarqué que la particule ne se comporte pas toujours de la même façon : parfois elle peut être séparée du verbe, parfois non. C'est précisément ce que nous allons systématiser maintenant. Il existe \cle{quatre types} de phrasal verbs, et les connaître, c'est savoir exactement où placer le complément.

\courssub{Type 1 : les phrasal verbs intransitifs (sans complément)}

\begin{definition}[Phrasal verb intransitif]
Un phrasal verb \cle{intransitif} est un verbe composé d'un \cle{verbe} et d'une \cle{particule} (adverbe), qui n'admet \textbf{aucun complément d'objet}. Le sens complet est porté par le verbe et sa particule : on ne peut rien ajouter après.
\end{definition}

\begin{exemplebox}[Observation]
\eng{I get up at five o'clock every morning.} « Je me lève à cinq heures tous les matins. »

\eng{The children grew up in Bamenda.} « Les enfants ont grandi à Bamenda. »

\eng{Come back before dark!} « Reviens avant la nuit ! »

\eng{Please sit down and listen.} « Asseyez-vous et écoutez, s'il vous plaît. »

\eng{The plane took off on time.} « L'avion a décollé à l'heure. »
\end{exemplebox}

\begin{propriete}[Règle du type 1]
Forme : \textbf{verbe + particule}, puis \textbf{rien} (ou éventuellement un complément circonstanciel introduit par une préposition).

Emploi : le sens est souvent différent du verbe simple : \eng{grow} = « pousser », mais \eng{grow up} = « grandir, devenir adulte » ; \eng{get up} = « se lever ».

Particules fréquentes : \eng{up, down, back, away, off, out}.
\end{propriete}

\begin{attention}[Piège francophone]
Ne cherchez pas à ajouter un complément direct après un phrasal verb intransitif. On ne dit jamais \eng{I get up my bed} pour « je me lève de mon lit » ; dites \eng{I get up} ou \eng{I get out of bed}. De même, \eng{grow up} ne prend jamais de complément : \eng{She grew up in Douala} (la préposition \eng{in} introduit un lieu, pas un objet).
\end{attention}

\courssub{Type 2 : les phrasal verbs transitifs séparables}

\begin{definition}[Phrasal verb séparable]
Un phrasal verb \cle{transitif séparable} admet un complément d'objet, et ce complément peut se placer \textbf{soit après la particule, soit entre le verbe et la particule}.
\end{definition}

\begin{exemplebox}[Les deux positions sont possibles avec un complément nominal]
\eng{Turn off the light.} = \eng{Turn the light off.} « Éteins la lumière. »

\eng{Put on your shoes.} = \eng{Put your shoes on.} « Mets tes chaussures. »

\eng{She picked up the phone.} = \eng{She picked the phone up.} « Elle a décroché le téléphone. »

\eng{Wake up the children.} = \eng{Wake the children up.} « Réveille les enfants. »
\end{exemplebox}

\begin{propriete}[La règle d'or : le pronom se place au milieu]
Avec un phrasal verb séparable, un complément \cle{pronom} (\eng{it, them, him, her, us}) se place \textbf{obligatoirement entre le verbe et la particule}.

\eng{She picked it up.} « Elle l'a décroché. » — jamais \eng{She picked up it.}

\eng{Turn it off, please.} « Éteins-le, s'il te plaît. » — jamais \eng{Turn off it.}

\eng{Put them on.} « Mets-les. » — jamais \eng{Put on them.}
\end{propriete}

\begin{attention}[Erreur fréquente des francophones]
En français, le pronom suit parfois le verbe à l'impératif (« éteins-le ») ou le précède (« je l'éteins »), mais il ne vient jamais \emph{au milieu}. Par calque, beaucoup d'élèves écrivent \eng{I turned off it} pour « je l'ai éteint ». C'est \textbf{faux} : dites \eng{I turned it off}. Retenez : \textbf{pronom = toujours au milieu} avec les verbes séparables.
\end{attention}

\courssub{Type 3 : les phrasal verbs transitifs inséparables}

\begin{definition}[Phrasal verb inséparable]
Un phrasal verb \cle{transitif inséparable} est formé d'un verbe et d'une particule (souvent une préposition : \eng{after, for, with, on, at}) ; le complément, nominal ou pronom, se place \textbf{toujours après la particule}. On ne peut jamais insérer quoi que ce soit entre le verbe et la particule.
\end{definition}

\begin{exemplebox}[Le complément suit toujours]
\eng{I look after my little sister.} « Je m'occupe de ma petite sœur. »

\eng{I look after her.} « Je m'occupe d'elle. » — jamais \eng{I look her after.}

\eng{They are looking for their keys.} « Ils cherchent leurs clés. »

\eng{She deals with customers every day.} « Elle s'occupe des clients tous les jours. »

\eng{You can rely on him.} « Tu peux compter sur lui. »
\end{exemplebox}

\begin{propriete}[Règle du type 3]
Forme : \textbf{verbe + particule + complément} (nom ou pronom), dans cet ordre, sans exception.

Astuce de reconnaissance : quand la particule est une vraie préposition (\eng{after, for, with, on, at, into}), le verbe est presque toujours inséparable, car la préposition « appelle » son complément, exactement comme en français « compter \emph{sur} quelqu'un ».
\end{propriete}

\courssub{Type 4 : les verbes à deux particules (phrasal-prepositional verbs)}

\begin{definition}[Phrasal-prepositional verb]
Certains verbes combinent \textbf{un verbe + une particule adverbiale + une préposition}. Ils sont toujours \cle{inséparables} : le complément vient après la préposition, jamais ailleurs.
\end{definition}

\begin{exemplebox}[Trois éléments soudés]
\eng{We look forward to the holidays.} « Nous attendons les vacances avec impatience. »

\eng{We look forward to them.} « Nous les attendons avec impatience. » — jamais \eng{We look them forward to.}

\eng{I can't put up with this noise.} « Je ne supporte pas ce bruit. »

\eng{She gets on well with her neighbours.} « Elle s'entend bien avec ses voisins. »

\eng{We have run out of sugar.} « Nous n'avons plus de sucre. » (littéralement : nous avons épuisé notre stock de sucre)
\end{exemplebox}

\begin{attention}[Attention à « to »]
Dans \eng{look forward to}, le mot \eng{to} est une \textbf{préposition}, pas une marque d'infinitif. Il est donc suivi d'un nom ou d'un verbe en \eng{-ing} : \eng{I look forward to seeing you.} « J'ai hâte de te voir. » — jamais \eng{I look forward to see you.}
\end{attention}

\courssub{Tableau récapitulatif des quatre types}

\begin{center}
\begin{tabularx}{\linewidth}{|l|X|X|X|}
\hline
\rowcolor{popblueL} Type & Structure & Exemple & Traduction \\
\hline
1. Intransitif & V + particule & \eng{The bus broke down.} & Le bus est tombé en panne. \\
\hline
2. Séparable & V + C + particule \emph{ou} V + particule + C & \eng{Turn the TV off / Turn off the TV / Turn it off.} & Éteins la télé / Éteins-la. \\
\hline
3. Inséparable & V + particule + C & \eng{She looks after the baby. / She looks after him.} & Elle garde le bébé. / Elle le garde. \\
\hline
4. Deux particules & V + part. + prép. + C & \eng{We ran out of fuel.} & Nous sommes tombés en panne d'essence. \\
\hline
\end{tabularx}
\end{center}

\begin{popfigure}
\begin{tikzpicture}[node distance=0.35cm]
\node[draw=popblue, fill=popblueL, rounded corners, minimum width=6.2cm, minimum height=1.5cm, align=center] (t1) {\textbf{Type 1 : intransitif}\\ V + particule\\ \eng{get up}};
\node[draw=popgreen, fill=popgreenL, rounded corners, minimum width=6.2cm, minimum height=1.5cm, align=center, right=0.6cm of t1] (t2) {\textbf{Type 2 : séparable}\\ V + (C) + particule + (C)\\ \eng{pick it up}};
\node[draw=poporange, fill=poporangeL, rounded corners, minimum width=6.2cm, minimum height=1.5cm, align=center, below=0.6cm of t1] (t3) {\textbf{Type 3 : inséparable}\\ V + particule + C\\ \eng{look after her}};
\node[draw=poppurple, fill=poppurpleL, rounded corners, minimum width=6.2cm, minimum height=1.5cm, align=center, below=0.6cm of t2] (t4) {\textbf{Type 4 : deux particules}\\ V + part. + prép. + C\\ \eng{look forward to it}};
\end{tikzpicture}
\legende{Les quatre structures de phrasal verbs : la place du complément (C) change tout.}
\end{popfigure}

\courssub{Le test pratique : séparable ou inséparable ?}

\begin{methode}[Le test du pronom « it »]
Pour savoir si un phrasal verb transitif est séparable, procédez ainsi :
\begin{enumerate}
\item Remplacez le complément nominal par le pronom \eng{it}.
\item Essayez de placer ce pronom \textbf{entre le verbe et la particule}.
\item Si la phrase est correcte (\eng{pick it up}), le verbe est \textbf{séparable} : avec un pronom, la séparation devient obligatoire.
\item Si la phrase est impossible (\eng{look it after} est faux ; il faut \eng{look after it}), le verbe est \textbf{inséparable} : le complément reste toujours après la particule.
\end{enumerate}
Exemples : \eng{turn on} $\rightarrow$ \eng{turn it on} (correct : séparable) ; \eng{deal with} $\rightarrow$ \eng{deal it with} (incorrect : inséparable, dites \eng{deal with it}).
\end{methode}

\begin{exoresolu}[Séparable ou inséparable ?]
Énoncé. Pour chaque phrase, réécrivez-la en remplaçant le complément souligné par un pronom, en le plaçant correctement : (a) \eng{Please wake up the baby.} (b) \eng{She is looking for her passport.} (c) \eng{He put on his uniform.}
\tcblower
Solution détaillée.

(a) \eng{wake up} : testons \eng{wake it up} — correct, donc séparable. Réponse : \eng{Please wake him up.} (le bébé = \eng{him} ou \eng{her}). « Réveille-le, s'il te plaît. »

(b) \eng{look for} : testons \eng{look it for} — impossible, donc inséparable. Le pronom reste après la particule : \eng{She is looking for it.} « Elle le cherche. »

(c) \eng{put on} : testons \eng{put it on} — correct, donc séparable. Réponse : \eng{He put it on.} « Il l'a mis. »
\end{exoresolu}

\begin{exercice}[À vous de jouer]
Classez ces phrasal verbs selon leur type (1, 2, 3 ou 4), puis traduisez chaque phrase : (a) \eng{The meeting started late, so we went back home.} (b) \eng{Can you turn down the radio?} (c) \eng{My grandmother looks after us during the holidays.} (d) \eng{I look forward to your letter.}
\end{exercice}

\begin{corrige}[Exercice : séparable ou inséparable ?]
(a) \eng{go back} : type 1, intransitif — « La réunion a commencé en retard, alors nous sommes rentrés à la maison. » (b) \eng{turn down} : type 2, séparable (\eng{turn it down}) — « Peux-tu baisser la radio ? » (c) \eng{look after} : type 3, inséparable (\eng{looks after us}) — « Ma grand-mère s'occupe de nous pendant les vacances. » (d) \eng{look forward to} : type 4, deux particules — « J'attends ta lettre avec impatience. »
\end{corrige}

\begin{aretenir}[L'essentiel de la section]
\begin{itemize}
\item Type 1 : verbe + particule, sans complément (\eng{get up}).
\item Type 2 : séparable ; avec un nom, deux positions possibles ; avec un pronom, séparation \textbf{obligatoire} (\eng{turn it off}).
\item Type 3 : inséparable ; le complément suit toujours (\eng{look after her}).
\item Type 4 : verbe + particule + préposition, inséparable (\eng{put up with it}).
\item En cas de doute, appliquez le test du pronom \eng{it}.
\end{itemize}
\end{aretenir}

\coursec{Phrasal verbs et locutions prépositionnelles : vie familiale et sociale}

\courssub{Les quatre types de phrasal verbs : rappel grammatical}

Avant d'aborder le vocabulaire du chapitre, il est indispensable de maîtriser la classification grammaticale des \cle{phrasal verbs}. Le type détermine la place de l'objet (nom ou pronom) et la possibilité de passer à la voix passive.

\begin{propriete}[Classification des phrasal verbs]
\begin{enumerate}
\item \textbf{Intransitifs} : pas d'objet. Ex. \eng{grow up, wake up, get up, go out, come back, fall out, make up, drop out, give up}.
\item \textbf{Transitifs séparables} : verbe + particule (adverbe). L'objet nom se place \textbf{avant ou après} la particule ; l'objet pronom \textbf{doit} se placer \textbf{entre} le verbe et la particule. Ex. \eng{bring up, put on, take off, tidy up, put off, find out, carry on (sth)}.
  \begin{itemize}
  \item \eng{She brought up the children.} / \eng{She brought the children up.} / \eng{She brought them up.}
  \item \textbf{Passif fréquent} : \eng{The children were brought up by their grandmother.}
  \end{itemize}
\item \textbf{Transitifs inséparables (verbes prépositionnels)} : verbe + préposition. L'objet (nom ou pronom) se place \textbf{toujours après} la préposition. Ex. \eng{look after, take after, look like, get on with, drop out of, catch up with, carry on with, put up with, find out (sth)}.
  \item \eng{She looks after the baby.} / \eng{She looks after him.} (jamais *\eng{looks him after})
  \item \textbf{Passif possible} : \eng{The baby is looked after by her.}
\item \textbf{Verbes à deux particules} : verbe + adverbe + préposition. Toujours inséparables. Ex. \eng{get on with, put up with, catch up with, carry on with, drop out of, look forward to}.
\end{enumerate}
\end{propriete}

\begin{attention}[Pronom objet = séparation obligatoire]
Avec un phrasal verb \textbf{séparable}, si l'objet est un pronom (\eng{me, you, him, her, it, us, them}), il \textbf{doit} s'intercaler : \eng{Turn it off} (jamais *\eng{Turn off it}). C'est une erreur majeure au Bac.
\end{attention}

\courssub{Observation : un texte pour découvrir le vocabulaire}

\begin{textebox}[A family in Bamenda]
Ali grew up in Bamenda, in the North-West Region of Cameroon. His parents brought up five children in a small house near the market. Every morning, Ali wakes up at six, gets up immediately, puts on his school uniform and helps his mother wash up before leaving for school. He takes after his father: they are both calm and hardworking. At school, Ali gets on well with his classmates and his teachers. Last year, his best friend dropped out of school because his family had no money, but Ali never gave up his studies. When he falls out with his younger brother, they always make up quickly, because their grandmother says that family is more important than pride. In the evening, Ali tidies up his room, does his homework and never puts off his duties. When he does not understand a lesson, he asks questions until he finds out the answer. His dream is to carry on his studies at the University of Buea and to come back to Bamenda as a doctor.
\end{textebox}

\textbf{Glossaire}

\begin{center}
\begin{tabularx}{\linewidth}{|l|X|l|}\hline
\rowcolor{popblueL} Mot anglais & Nature & Traduction \\ \hline
to grow up & verbe & grandir \\ \hline
to bring up & verbe & élever (des enfants) \\ \hline
to take after & verbe & ressembler à (un parent) \\ \hline
to get on with & verbe & s'entendre avec \\ \hline
to drop out & verbe & abandonner (ses études) \\ \hline
to give up & verbe & abandonner, renoncer \\ \hline
to fall out & verbe & se brouiller \\ \hline
to make up & verbe & se réconcilier \\ \hline
to tidy up & verbe & ranger \\ \hline
to put off & verbe & remettre à plus tard \\ \hline
to find out & verbe & découvrir \\ \hline
to carry on & verbe & continuer \\ \hline
\end{tabularx}
\end{center}

\textbf{Questions de compréhension}
\begin{enumerate}
\item Where did Ali grow up?
\item Who does Ali take after?
\item Why did his best friend drop out of school?
\item What does Ali do when he falls out with his brother?
\end{enumerate}

\begin{aretenir}[Observation grammaticale]
Dans ce texte, chaque verbe souligné par le sens est un \cle{phrasal verb} : verbe + particule. Remarque que le sens global \textbf{ne se devine pas} à partir du verbe seul : \eng{to make up} ne signifie pas « faire vers le haut » mais « se réconcilier ». Il faut donc apprendre chaque phrasal verb comme un \textbf{mot de vocabulaire à part entière}, avec sa traduction, son type (séparable/inséparable) et ses formes irrégulières.
\end{aretenir}

\courssub{Carte mentale : trois familles de sens}

\begin{popfigure}
\begin{tikzpicture}[
  mindmap,
  grow cyclic,
  every node/.style={concept, execute at begin node=\hskip0pt},
  root concept/.append style={concept color=popblue, font=\bfseries\large, text=white},
  level 1/.append style={level distance=3.5cm, sibling angle=120, concept color=popgreen!80!popblue, font=\bfseries, text=white},
  level 2/.append style={level distance=2.5cm, sibling angle=50, font=\small, text=white},
  text=white
]
\node[root concept]{PHRASAL VERBS}
  child[concept color=popgreen]{ node {FAMILLE}
    child { node {bring up} }
    child { node {grow up} }
    child { node {look after} }
    child { node {take after} }
    child { node {fall out} }
    child { node {make up} }
  }
  child[concept color=poporange]{ node {ROUTINE QUOTIDIENNE}
    child { node {wake up} }
    child { node {get up} }
    child { node {put on} }
    child { node {take off} }
    child { node {tidy up} }
    child { node {wash up} }
    child { node {go out / come back} }
  }
  child[concept color=poppurple]{ node {VIE SOCIALE \& SCOLAIRE}
    child { node {get on with} }
    child { node {drop out (of)} }
    child { node {catch up (with)} }
    child { node {give up} }
    child { node {carry on (with)} }
    child { node {find out} }
    child { node {put up with} }
  };
\end{tikzpicture}
\legende{Carte mentale des phrasal verbs essentiels du chapitre, classés en trois champs lexicaux.}
\end{popfigure}

\courssub{Le champ lexical de la famille}

\begin{definition}[Les six verbes de la vie familiale]
\begin{itemize}
\item \eng{to bring up} = élever (des enfants). Verbe transitif \textbf{séparable} : \eng{She brought up six children} ou \eng{She brought six children up}. Au passif, le participe passé est très fréquent : \eng{He was brought up by his grandmother.}
\item \eng{to grow up} = grandir. Verbe \textbf{intransitif} : on ne peut pas « grow up quelqu'un ». \eng{I grew up in Douala.}
\item \eng{to look after} = s'occuper de, prendre soin de. Transitif \textbf{inséparable} (verbe prépositionnel) : \eng{She looks after her little sister.}
\item \eng{to take after} = ressembler à (un parent, trait de caractère hérité). Transitif \textbf{inséparable} : \eng{Ngono takes after her father.}
\item \eng{to fall out (with somebody)} = se brouiller (avec quelqu'un). Intransitif, suivi de la préposition \eng{with} pour introduire la personne.
\item \eng{to make up (with somebody)} = se réconcilier (avec quelqu'un). C'est le contraire de \eng{fall out}. Intransitif.
\end{itemize}
\end{definition}

\begin{exemplebox}[Exemples traduits — la famille]
\eng{My grandmother brought up six children.} « Ma grand-mère a élevé six enfants. »

\eng{Ngono takes after her father.} « Ngono ressemble à son père. »

\eng{They fell out but made up quickly.} « Ils se sont brouillés mais se sont vite réconciliés. »

\eng{My aunt looks after the twins while their mother is at work.} « Ma tante s'occupe des jumeaux pendant que leur mère travaille. »

\eng{Children who grow up in a village learn many skills.} « Les enfants qui grandissent dans un village apprennent de nombreux savoir-faire. »
\end{exemplebox}

\begin{attention}[Ne confonds pas \eng{take after}, \eng{look like} et \eng{look after}]
Ces trois verbes sont des pièges classiques pour les francophones :
\begin{itemize}
\item \eng{to take after somebody} = ressembler à un \textbf{parent} (caractère ou physique hérité) : \eng{He takes after his mother.} « Il tient de sa mère. »
\item \eng{to look like somebody} = ressembler \textbf{physiquement} à quelqu'un en général : \eng{She looks like her teacher.} « Elle ressemble à son professeur. »
\item \eng{to look after somebody} = \textbf{s'occuper} de quelqu'un : \eng{I look after my grandfather.} « Je m'occupe de mon grand-père. »
\end{itemize}
Une seule particule change tout le sens : \eng{look after} $\neq$ \eng{look like}. Au Bac, cette confusion coûte des points en compréhension comme en expression écrite.
\end{attention}

\courssub{Le champ lexical de la vie quotidienne}

\begin{definition}[Les verbes de la routine quotidienne]
\begin{itemize}
\item \eng{to wake up} = se réveiller (intransitif) ; \eng{to wake somebody up} = réveiller quelqu'un (transitif \textbf{séparable}).
\item \eng{to get up} = se lever, sortir du lit (intransitif).
\item \eng{to put on} = mettre (un vêtement) — transitif \textbf{séparable} : \eng{Put on your shoes} ou \eng{Put your shoes on}.
\item \eng{to take off} = enlever (un vêtement) — transitif \textbf{séparable}. Contraire de \eng{put on}. Sens secondaire : décoller (avion, intransitif).
\item \eng{to tidy up} = ranger — peut être intransitif (\eng{I must tidy up}) ou transitif séparable (\eng{Tidy up your room} / \eng{Tidy your room up}).
\item \eng{to wash up} = faire la vaisselle (intransitif en anglais britannique, standard au Cameroun).
\item \eng{to go out} = sortir (se divertir) ; \eng{to come back} = revenir.
\end{itemize}
\end{definition}

\begin{exemplebox}[Exemples traduits — la routine]
\eng{I wake up at five thirty and get up at six.} « Je me réveille à cinq heures trente et je me lève à six heures. »

\eng{Put on your uniform, we are late!} « Mets ton uniforme, nous sommes en retard ! »

\eng{Take off your shoes before entering the house.} « Enlève tes chaussures avant d'entrer dans la maison. »

\eng{After dinner, the children wash up and tidy up the kitchen.} « Après le dîner, les enfants font la vaisselle et rangent la cuisine. »

\eng{On Saturdays, we go out with friends and come back late.} « Le samedi, nous sortons avec des amis et nous rentrons tard. »
\end{exemplebox}

\courssub{Le champ lexical social et scolaire}

\begin{definition}[Les verbes de la vie sociale et scolaire]
\begin{itemize}
\item \eng{to get on with somebody} = s'entendre avec quelqu'un (verbe à \textbf{deux particules}, inséparable). Équivalent américain : \eng{get along with}.
\item \eng{to drop out (of school)} = abandonner ses études, décrocher.
\item \eng{to catch up (with)} = rattraper (son retard, quelqu'un).
\item \eng{to give up} = abandonner, renoncer : \eng{give up smoking} « arrêter de fumer ». Transitif séparable ou intransitif.
\item \eng{to carry on} = continuer. Souvent \eng{carry on with sth} (inséparable) ou \eng{carry on doing sth}.
\item \eng{to find out} = découvrir, apprendre (par recherche) : \eng{I found out the truth.} Transitif séparable.
\item \eng{to put up with} = tolérer, supporter (verbe à deux particules, inséparable) : \eng{I can't put up with this noise.}
\end{itemize}
\end{definition}

\begin{exemplebox}[Exemples traduits — vie sociale et scolaire]
\eng{Ali gets on well with his neighbours.} « Ali s'entend bien avec ses voisins. »

\eng{Many students drop out of school because of poverty.} « Beaucoup d'élèves abandonnent l'école à cause de la pauvreté. »

\eng{After two weeks of illness, she worked hard to catch up with the class.} « Après deux semaines de maladie, elle a travaillé dur pour rattraper la classe. »

\eng{Never give up your dreams.} « N'abandonne jamais tes rêves. »

\eng{The teacher told us to carry on with the exercise.} « Le professeur nous a dit de continuer l'exercice. »

\eng{We must find out why the boy is always absent.} « Nous devons découvrir pourquoi le garçon est toujours absent. »

\eng{She put up with his bad behaviour for years.} « Elle a supporté son mauvais comportement pendant des années. »
\end{exemplebox}

\courssub{Les prepositional phrases (locutions prépositionnelles)}

\begin{definition}[Prepositional phrases]
Une \cle{prepositional phrase} (ou locution prépositionnelle) est un groupe de mots qui fonctionne comme une préposition unique (avant un nom/pronom/gérondif). Elle ne contient \textbf{pas de verbe}. Elles sont indispensables pour structurer une composition (connecteurs logiques).
\end{definition}

\begin{center}
\begin{tabularx}{\linewidth}{|l|X|l|}\hline
\rowcolor{popblueL} Locution & Traduction & Exemple \\ \hline
in front of & devant & The car is \eng{in front of} the house. \\ \hline
behind & derrière & (préposition simple, mais souvent groupée) \\ \hline
next to / beside & à côté de & She sits \eng{next to} me. \\ \hline
because of / due to & à cause de / en raison de & \eng{Because of} the rain, we stayed home. \\ \hline
in spite of / despite & malgré & \eng{In spite of} his age, he runs fast. \\ \hline
instead of & au lieu de & \eng{Instead of} complaining, act. \\ \hline
according to & selon & \eng{According to} the news, it will rain. \\ \hline
in addition to & en plus de & \eng{In addition to} English, she speaks German. \\ \hline
on behalf of & au nom de & I speak \eng{on behalf of} the class. \\ \hline
with regard to / concerning & concernant & \eng{With regard to} your application... \\ \hline
\end{tabularx}
\end{center}

\begin{attention}[Prepositional phrase $\neq$ Phrasal verb]
\begin{itemize}
\item \textbf{Phrasal verb} = Verbe + Particule(s). Fonction : verbe. \eng{He \textbf{gave up} smoking.}
\item \textbf{Prepositional phrase} = Préposition complexe (2-3 mots). Fonction : complément circonstanciel. \eng{\textbf{In spite of} the rain, he came.}
\end{itemize}
Ne confonds pas \eng{drop out of} (phrasal verb à 3 mots, verbe) et \eng{because of} (locution prépositionnelle, pas de verbe).
\end{attention}

\courssub{Tableau récapitulatif anglais / français}

\begin{center}
\begin{tabularx}{\linewidth}{|l|X|X|}\hline
\rowcolor{popblueL} Phrasal verb & Traduction & Exemple en contexte \\ \hline
bring up & élever & My grandmother brought up six children. \\ \hline
grow up & grandir & Ali grew up in Bamenda. \\ \hline
look after & s'occuper de & She looks after her little sister. \\ \hline
take after & ressembler à (parent) & Ngono takes after her father. \\ \hline
fall out & se brouiller & The brothers fell out over money. \\ \hline
make up & se réconcilier & They made up after the quarrel. \\ \hline
get up & se lever & I get up at six every morning. \\ \hline
wake up & se réveiller & The baby wakes up at night. \\ \hline
put on & mettre (vêtement) & Put on your shoes. \\ \hline
take off & enlever ; décoller & Take off your hat. / The plane took off. \\ \hline
tidy up & ranger & Tidy up your room, please. \\ \hline
wash up & faire la vaisselle & We wash up after lunch. \\ \hline
go out & sortir & They go out on Saturdays. \\ \hline
come back & revenir & He came back from Buea. \\ \hline
get on with & s'entendre avec & I get on well with my classmates. \\ \hline
drop out (of) & abandonner (études) & His friend dropped out of school. \\ \hline
catch up (with) & rattraper & She caught up with the class. \\ \hline
give up & abandonner & Never give up. \\ \hline
carry on (with) & continuer & Carry on with your work. \\ \hline
find out & découvrir & I found out the truth. \\ \hline
put up with & tolérer & I can't put up with this noise. \\ \hline
\end{tabularx}
\end{center}

\courssub{Synonymes « savants » : un atout pour le Bac}

À l'examen, les exercices de transformation demandent souvent de remplacer un verbe simple (soutenu) par un phrasal verb (courant), ou l'inverse. Voici les paires à connaître par cœur :

\begin{center}
\begin{tabularx}{\linewidth}{|X|X|X|}\hline
\rowcolor{popgreenL} Verbe soutenu & Phrasal verb & Exemple \\ \hline
to postpone & to put off & They put off the meeting. \\ \hline
to discover & to find out & We found out the truth. \\ \hline
to continue & to carry on & Carry on reading. \\ \hline
to abandon & to give up & He gave up smoking. \\ \hline
to tolerate & to put up with & I can't put up with noise. \\ \hline
to raise (children) & to bring up & She brought up five children. \\ \hline
to resemble (parent) & to take after & He takes after his uncle. \\ \hline
to quarrel & to fall out & They fell out yesterday. \\ \hline
to reconcile & to make up & They made up later. \\ \hline
\end{tabularx}
\end{center}

\begin{savaistu}[Les transformations au Bac]
Au baccalauréat camerounais, les exercices de \emph{sentence transformation} adorent les phrasal verbs : on te donne une phrase avec un verbe simple et tu dois la réécrire avec le phrasal verb équivalent sans changer le sens. Exemple : \eng{The students continued their work.} $\rightarrow$ \eng{The students carried on with their work.} Autre cas classique : \eng{I cannot tolerate his behaviour.} $\rightarrow$ \eng{I cannot put up with his behaviour.} Apprendre les paires du tableau ci-dessus, c'est gagner des points presque automatiquement.
\end{savaistu}

\courssub{Un même verbe, plusieurs sens : le cas de \eng{take}}

La particule change radicalement le sens du verbe de base. Le verbe \eng{take} en est le meilleur exemple :

\begin{center}
\begin{tabular}{|l|l|l|}\hline
\rowcolor{poppurpleL} Verbe & Sens & Exemple \\ \hline
take off & enlever & Take off your jacket. \\ \hline
take off & décoller (avion) & The plane took off at noon. \\ \hline
take up & se mettre à (activité) & She took up football at ten. \\ \hline
take after & ressembler à & He takes after his father. \\ \hline
take over & reprendre (entreprise) & His son took over the shop. \\ \hline
\end{tabular}
\end{center}

\begin{attention}[Un phrasal verb = un mot nouveau]
Ne traduis jamais un phrasal verb mot à mot : \eng{The plane took off} ne veut pas dire « l'avion a pris dehors » mais « l'avion a décollé ». De même, \eng{take up a sport} signifie « se mettre à un sport », et non « prendre un sport vers le haut ». Apprends chaque combinaison verbe + particule comme une entrée de dictionnaire indépendante, avec son exemple.
\end{attention}

\courssub{Entraînement guidé}

\begin{exoresolu}[Transformation : du verbe simple au phrasal verb]
\textbf{Énoncé.} Réécris chaque phrase en remplaçant le verbe souligné par le phrasal verb équivalent, sans changer le sens.
\begin{enumerate}
\item My aunt \textbf{raised} four children alone.
\item The students \textbf{continued} their revision after the break.
\item She \textbf{discovered} that her friend had lied.
\item He \textbf{abandoned} his studies in 2022.
\item They \textbf{tolerate} the noise from the bar.
\end{enumerate}
\tcblower
\textbf{Solution détaillée.}
\begin{enumerate}
\item \eng{My aunt brought up four children alone.} — \eng{raise} (élever) = \eng{bring up} ; attention au prétérit irrégulier : \eng{bring} $\rightarrow$ \eng{brought}.
\item \eng{The students carried on with their revision after the break.} — \eng{continue} = \eng{carry on (with)} ; prétérit régulier : \eng{carried}.
\item \eng{She found out that her friend had lied.} — \eng{discover} = \eng{find out} ; prétérit irrégulier : \eng{find} $\rightarrow$ \eng{found}.
\item \eng{He gave up his studies in 2022.} — \eng{abandon} = \eng{give up} ; prétérit irrégulier : \eng{give} $\rightarrow$ \eng{gave}.
\item \eng{They put up with the noise from the bar.} — \eng{tolerate} = \eng{put up with} (inséparable, deux particules) ; prétérit : \eng{put} (invariable).
\end{enumerate}
\textbf{Méthode :} repère le verbe simple, choisis le phrasal verb équivalent, puis conjugue-le au même temps que le verbe d'origine en faisant attention aux formes irrégulières et à la structure (séparable / inséparable).
\end{exoresolu}

\begin{exercice}[À toi de jouer]
Complète les phrases avec le phrasal verb qui convient, conjugué au bon temps : \eng{bring up, fall out, make up, look after, get on with, put off, take after, grow up}.
\begin{enumerate}
\item Ekane .................. in a small village near Ebolowa.
\item The two sisters .................. last week, but they .................. the next day.
\item Who .................. your baby brother when your mother is at the market?
\item My parents .................. the wedding because of the rain. (remettre à plus tard)
\item Little Andela .................. her grandmother: they both love singing.
\item Do you .................. well .................. your new classmates?
\end{enumerate}
\end{exercice}

\begin{corrige}[Exercice : À toi de jouer]
\begin{enumerate}
\item \eng{grew up} (prétérit irrégulier de \eng{grow up} : grow, grew, grown).
\item \eng{fell out} ; \eng{made up} (prétérits irréguliers : fall, fell, fallen ; make, made, made).
\item \eng{looks after} (présent simple, troisième personne : \eng{-s} obligatoire ; inséparable).
\item \eng{put off} (prétérit : \eng{put} est invariable : put, put, put ; séparable mais objet \eng{the wedding} après).
\item \eng{takes after} (inséparable : la particule \eng{after} reste collée au verbe ; présent 3e pers. sg).
\item \eng{get on ... with} (verbe à deux particules : \eng{get on with somebody} ; présent simple).
\end{enumerate}
\end{corrige}

\begin{aretenir}[L'essentiel de la leçon]
\begin{itemize}
\item \textbf{4 types de phrasal verbs} : intransitifs, transitifs séparables (objet pronom $\rightarrow$ intercalation), transitifs inséparables (verbes prépositionnels), verbes à deux particules (toujours inséparables).
\item \textbf{Trois familles lexicales} : \textbf{Famille} (\eng{bring up, grow up, look after, take after, fall out, make up}), \textbf{Routine} (\eng{wake up, get up, put on, take off, tidy up, wash up, go out, come back}), \textbf{Vie sociale/scolaire} (\eng{get on with, drop out, catch up, give up, carry on, find out, put up with}).
\item \textbf{Prepositional phrases} : locutions prépositionnelles sans verbe (\eng{because of, in spite of, instead of, on behalf of, in addition to...}) pour lier les idées dans une composition.
\item \textbf{Transformations Bac} : maîtriser les paires \eng{postpone/put off}, \eng{discover/find out}, \eng{continue/carry on}, \eng{abandon/give up}, \eng{tolerate/put up with}, \eng{raise/bring up}, \eng{resemble/take after}.
\item \textbf{Pièges} : \eng{take after} (héritage) $\neq$ \eng{look like} (physique) $\neq$ \eng{look after} (soins) ; séparation obligatoire avec pronom (\eng{turn it off}) ; passif des séparables (\eng{was brought up}).
\end{itemize}
\end{aretenir}

\coursec{Les prepositional phrases}

Dans la section précédente, nous avons vu que les \emph{phrasal verbs} associent un verbe à une particule (\eng{look after}, \eng{get on with}, \eng{put up with}). Nous allons maintenant étudier une famille voisine, encore plus fréquente dans la langue : les \cle{prepositional phrases}, c'est-à-dire les groupes de mots construits autour d'une préposition. Elles sont partout : \eng{in the morning}, \eng{at school}, \eng{with my friends}, \eng{because of the rain}. Bien les maîtriser, c'est gagner en précision et en naturel, à l'oral comme à l'écrit.

\courssub{Observation et découverte}

\begin{textebox}[A letter from Buea]
Dear Auntie,

I am writing from Buea, where I have been living with my uncle's family since September. Life here is different from life in Douala. In the morning, I wake up at 6 a.m. and I walk to school with my cousins. We are never late, because we leave home on time.

My uncle is very proud of his children. He is responsible for the whole family, and everybody depends on him. Last week, my cousin Manka succeeded in passing her entrance examination, so the family organised a small party in front of the house. In spite of the rain, many neighbours came, and we danced until midnight. On behalf of the family, my uncle thanked everybody for coming.

I am very fond of this town, but sometimes I am afraid of the cold weather on the mountain. I am also interested in the local history, so I often listen to the old people who tell stories about the past. I apologised for not writing earlier, but I was busy with my studies.

Please, give my love to everybody at home. I am looking forward to seeing you during the holidays.

Your loving niece,
Clarisse
\end{textebox}

\textbf{Glossaire}

\begin{center}
\begin{tabularx}{\linewidth}{|l|l|X|}
\hline
\rowcolor{popblueL} Mot anglais & Nature & Traduction française \\
\hline
to wake up & verbe & se réveiller \\
\hline
on time & locution & à l'heure \\
\hline
proud of & adjectif + préposition & fier de \\
\hline
entrance examination & nom & examen d'entrée \\
\hline
in spite of & préposition complexe & malgré \\
\hline
on behalf of & préposition complexe & au nom de, de la part de \\
\hline
fond of & adjectif + préposition & amateur de, attaché à \\
\hline
to look forward to & verbe & attendre avec impatience \\
\hline
\end{tabularx}
\end{center}

\textbf{Questions de compréhension}

\begin{enumerate}
\item Where does Clarisse live now, and with whom?
\item Why is the family proud of Manka?
\item What happened in spite of the rain?
\item What is Clarisse interested in?
\end{enumerate}

\textbf{Observation.} Soulignons dans la lettre tous les groupes introduits par une préposition :
\begin{itemize}
\item \eng{from Buea}, \eng{with my uncle's family}, \eng{since September}, \eng{in the morning}, \eng{at 6 a.m.}, \eng{to school}, \eng{with my cousins}, \eng{on time} (lieu, temps, manière) ;
\item \eng{proud of his children}, \eng{responsible for the whole family}, \eng{depend on him} (complément d'adjectif / verbe prépositionnel) ;
\item \eng{in front of the house}, \eng{in spite of the rain}, \eng{on behalf of the family} (prépositions complexes) ;
\item \eng{fond of this town}, \eng{afraid of the cold weather}, \eng{interested in the local history}, \eng{apologised for not writing}, \eng{busy with my studies}, \eng{looking forward to seeing you} (adjectif / verbe + préposition + gérondif).
\end{itemize}
Chacun de ces groupes est une \cle{prepositional phrase}.

\courssub{Définition et structure}

\begin{definition}[Prepositional phrase]
Une \cle{prepositional phrase} (groupe prépositionnel) est un groupe de mots introduit par une \textbf{préposition} (\eng{in}, \eng{on}, \eng{at}, \eng{with}, \eng{for}, \eng{of}, \eng{to}, \eng{from}, \eng{because of}...) et suivi d'un \textbf{nom}, d'un \textbf{pronom} ou d'un \textbf{gérondif} : \eng{at school}, \eng{with them}, \eng{good at cooking} (\eng{at} + V-ing).
\end{definition}

\begin{popfigure}
\begin{center}
\begin{tikzpicture}[
  box/.style={draw=popblue, thick, rounded corners, fill=popblueL, align=center, inner sep=6pt},
  part/.style={draw=poporange, thick, rounded corners, fill=poporangeL, align=center, inner sep=5pt},
  lab/.style={align=center, font=\small\itshape, text=popmuted}
]
\node[box, font=\bfseries] (pp) {PREPOSITIONAL PHRASE};
\node[part, below left=1.1cm and 0.4cm of pp, align=center] (prep){PREPOSITION\\ \small in / on / at / with / for / of};
\node[part, below right=1.1cm and 0.4cm of pp, align=center] (np){NOUN PHRASE\\ \small the village / my parents / school / them / cooking};
\draw[-{Stealth}, thick, popblue] (pp) -- (prep);
\draw[-{Stealth}, thick, popblue] (pp) -- (np);
\node[lab, below=0.15cm of prep] {jamais seule};
\node[lab, below=0.15cm of np] {nom, pronom ou V-ing};
\end{tikzpicture}
\end{center}
\legende{Structure de la prepositional phrase : une préposition suivie d'un groupe nominal.}
\end{popfigure}

\begin{aretenir}[La règle d'or]
Une préposition anglaise n'est \textbf{jamais seule} : elle est toujours suivie de son complément. On ne peut pas dire \eng{I am interested in.} ✗ ; il faut compléter : \eng{I am interested in music.} (« Je m'intéresse à la musique. »)
\end{aretenir}

\courssub{Les fonctions des prepositional phrases}

Les groupes prépositionnels jouent plusieurs rôles dans la phrase. Les cinq fonctions essentielles à connaître en Terminale sont les suivantes.

\textbf{1. Complément de lieu} — il répond à la question \emph{où ?}

\eng{My grandmother lives in the village.} « Ma grand-mère vit au village. »

\eng{There is a mango tree in front of our house.} « Il y a un manguier devant notre maison. »

\textbf{2. Complément de temps} — il répond à la question \emph{quand ?}

\eng{We visit our grandparents on Sundays.} « Nous rendons visite à nos grands-parents le dimanche. »

\eng{The meeting starts at 5 p.m.} « La réunion commence à 17 heures. »

\eng{My sister was born in 2008.} « Ma sœur est née en 2008. »

\textbf{3. Complément de manière ou de moyen} — il répond à la question \emph{comment ?}

\eng{She speaks to the children with care.} « Elle parle aux enfants avec soin. »

\eng{My father goes to work by bus.} « Mon père va au travail en bus. »

\textbf{4. Complément d'adjectif} — il complète un adjectif.

\eng{We are proud of our school.} « Nous sommes fiers de notre école. »

\eng{She is interested in politics.} « Elle s'intéresse à la politique. »

\eng{My brother is very good at football.} « Mon frère est très bon au football. »

\textbf{5. Complément d'objet prépositionnel (verbe + préposition)} — il complète un verbe qui exige une préposition fixe.

\eng{Children depend on their parents.} « Les enfants dépendent de leurs parents. »

\eng{She succeeded in passing the exam.} « Elle a réussi à l'examen. »

\eng{This book belongs to me.} « Ce livre m'appartient. »

\courssub{Constructions figées : verbe + préposition et adjectif + préposition}

En anglais, de très nombreux verbes et adjectifs se construisent avec une préposition \textbf{fixe}, qu'il faut apprendre avec eux, comme une partie du mot. C'est un des points où le français et l'anglais diffèrent le plus.

\begin{center}
\begin{tabularx}{\linewidth}{|X|X|X|}
\hline
\rowcolor{popblueL} Construction anglaise & Exemple anglais & Traduction française \\
\hline
to listen \textbf{to} & I listen to the radio. & J'écoute la radio. \\
\hline
to wait \textbf{for} & We are waiting for the bus. & Nous attendons le bus. \\
\hline
to belong \textbf{to} & This land belongs to the family. & Cette terre appartient à la famille. \\
\hline
to depend \textbf{on} & Children depend on their parents. & Les enfants dépendent de leurs parents. \\
\hline
to apologise \textbf{for} & He apologised for the delay. & Il s'est excusé du retard. \\
\hline
to succeed \textbf{in} & She succeeded in her exam. & Elle a réussi à son examen. \\
\hline
to insist \textbf{on} & They insisted on paying. & Ils ont insisté pour payer. \\
\hline
to look \textbf{at} & Look at the picture. & Regarde l'image. \\
\hline
to think \textbf{of / about} & I am thinking of my mother. & Je pense à ma mère. \\
\hline
to pay \textbf{for} & Who paid for the meal? & Qui a payé le repas ? \\
\hline
\end{tabularx}
\end{center}

\begin{center}
\begin{tabularx}{\linewidth}{|X|X|X|}
\hline
\rowcolor{popgreenL} Adjectif + préposition & Exemple anglais & Traduction française \\
\hline
afraid \textbf{of} & The child is afraid of dogs. & L'enfant a peur des chiens. \\
\hline
fond \textbf{of} & She is fond of music. & Elle aime beaucoup la musique. \\
\hline
angry \textbf{with} & Father is angry with his son. & Le père est en colère contre son fils. \\
\hline
responsible \textbf{for} & He is responsible for the shop. & Il est responsable de la boutique. \\
\hline
different \textbf{from} & Buea is different from Douala. & Buea est différente de Douala. \\
\hline
interested \textbf{in} & I am interested in history. & Je m'intéresse à l'histoire. \\
\hline
good \textbf{at} & She is good at cooking. & Elle est bonne en cuisine. \\
\hline
proud \textbf{of} & We are proud of you. & Nous sommes fiers de toi. \\
\hline
married \textbf{to} & My aunt is married to a teacher. & Ma tante est mariée à un enseignant. \\
\hline
full \textbf{of} & The market is full of people. & Le marché est plein de monde. \\
\hline
\end{tabularx}
\end{center}

\begin{exemplebox}[Exemples traduits]
\eng{We depend on our parents.} « Nous dépendons de nos parents. »

\eng{She succeeded in passing the exam.} « Elle a réussi à l'examen. »

\eng{He apologised for being late.} « Il s'est excusé d'être en retard. »

\eng{This house belongs to my grandfather.} « Cette maison appartient à mon grand-père. »

\eng{The children are fond of their grandmother.} « Les enfants adorent leur grand-mère. »
\end{exemplebox}

\begin{propriete}[Préposition + gérondif]
Après une préposition, le verbe se met \textbf{toujours au gérondif} (-ing), jamais à l'infinitif.

\eng{He insisted on paying.} « Il a insisté pour payer. »

\eng{Thank you for coming.} « Merci d'être venu. »

\eng{She is good at singing.} « Elle chante bien. »

\eng{He left without saying goodbye.} « Il est parti sans dire au revoir. »
\end{propriete}

\begin{attention}[Piège du francophone : l'infinitif après une préposition]
En français, on dit « merci de venir », « il est parti sans payer » avec un infinitif. En anglais, c'est impossible : après une préposition, le verbe prend \textbf{-ing}. Ne dites jamais \eng{Thank you for to come} ✗ ni \eng{Thank you for come} ✗, mais bien \eng{Thank you for coming} ✓.
\end{attention}

\courssub{Les prépositions complexes}

Certaines prépositions sont formées de deux ou trois mots, mais fonctionnent comme une seule préposition. Elles sont très fréquentes à l'écrit et dans les discours formels.

\begin{center}
\begin{tabularx}{\linewidth}{|X|X|X|}
\hline
\rowcolor{popgoldL} Préposition complexe & Traduction & Exemple anglais \\
\hline
in front of & devant & The taxi is in front of the school. \\
\hline
in spite of & malgré & In spite of the rain, we played football. \\
\hline
because of & à cause de & The match was cancelled because of the storm. \\
\hline
instead of & au lieu de & She took a bendskin instead of a taxi. \\
\hline
on behalf of & au nom de & He spoke on behalf of the family. \\
\hline
according to & selon & According to my uncle, the harvest is good. \\
\hline
thanks to & grâce à & Thanks to her efforts, she passed. \\
\hline
next to & à côté de & The church is next to the market. \\
\hline
\end{tabularx}
\end{center}

\begin{exemplebox}[Exemples traduits]
\eng{In spite of his age, my grandfather still works on the farm.} « Malgré son âge, mon grand-père travaille encore à la ferme. »

\eng{Instead of watching television, the children listened to their grandmother's stories.} « Au lieu de regarder la télévision, les enfants ont écouté les histoires de leur grand-mère. »

\eng{On behalf of the students, I thank our teachers.} « Au nom des élèves, je remercie nos professeurs. »
\end{exemplebox}

\begin{attention}[Ne pas confondre : préposition vs conjonction]
\eng{Because of} + nom (ou gérondif), mais \eng{because} + phrase complète (sujet + verbe) : \eng{We stayed at home because of the rain.} / \eng{We stayed at home because it was raining.} De même, \eng{in spite of} / \eng{despite} + nom : \eng{In spite of the rain, we went out.} ; pour une phrase complète, utilisez \eng{although} / \eng{though} : \eng{Although it was raining, we went out.}
\end{attention}

\courssub{Place du groupe prépositionnel dans la phrase}

\begin{propriete}[Position de la prepositional phrase]
Le groupe prépositionnel se place généralement \textbf{après le verbe} ou \textbf{après le complément direct}. Il ne s'intercale \textbf{jamais entre le verbe et son complément direct}.

\eng{She cooks the food with care.} ✓ « Elle prépare la nourriture avec soin. »

\eng{She cooks with care the food.} ✗ (ordre impossible)

Les compléments de temps et de lieu peuvent aussi se placer en tête de phrase, suivis d'une virgule : \eng{In the morning, I walk to school.} « Le matin, je vais à l'école à pied. »
\end{propriete}

\begin{attention}[Erreur fréquente : calque du français sur les prépositions]
Le français « écouter », « attendre », « regarder », « payer » se construisent sans préposition ; l'anglais en exige une. Mémorisez ces paires :

to listen the radio ✗ → to listen \textbf{to} the radio ✓

to wait the bus ✗ → to wait \textbf{for} the bus ✓

to look the picture ✗ → to look \textbf{at} the picture ✓

to pay the meal ✗ → to pay \textbf{for} the meal ✓

À l'inverse, certains verbes anglais n'ont pas de préposition alors que le français en a une : \eng{to enter the room} (« entrer dans la pièce », sans \eng{in}), \eng{to marry somebody} (« épouser quelqu'un », sans \eng{to} ni \eng{with}), \eng{to answer the question} (« répondre à la question », sans \eng{to}), \eng{to reach the village} (« arriver au village », sans \eng{at}).
\end{attention}

\courssub{Entraînement guidé}

\begin{exoresolu}[Complétez avec la préposition qui convient]
1. My little sister is afraid ... dogs. \quad 2. We are waiting ... the headmaster. \quad 3. He apologised ... being rude. \quad 4. This bag belongs ... my mother. \quad 5. She insisted ... cooking the meal herself. \quad 6. ... spite of the noise, the baby is sleeping.
\tcblower
1. \textbf{of} — \eng{My little sister is afraid of dogs.} « Ma petite sœur a peur des chiens. » (adjectif \eng{afraid} + \eng{of}).

2. \textbf{for} — \eng{We are waiting for the headmaster.} « Nous attendons le proviseur. » (verbe \eng{to wait for}, sans équivalent français de la préposition).

3. \textbf{for} — \eng{He apologised for being rude.} « Il s'est excusé d'avoir été impoli. » (\eng{apologise for} + gérondif \eng{being}).

4. \textbf{to} — \eng{This bag belongs to my mother.} « Ce sac appartient à ma mère. »

5. \textbf{on} — \eng{She insisted on cooking the meal herself.} « Elle a insisté pour préparer le repas elle-même. » (\eng{insist on} + gérondif).

6. \textbf{In} — \eng{In spite of the noise, the baby is sleeping.} « Malgré le bruit, le bébé dort. » (préposition complexe \eng{in spite of}).
\end{exoresolu}

\begin{exercice}[À vous de jouer]
A. Traduisez en anglais : 1. Nous écoutons les informations à la radio. 2. Merci d'être venus à la fête. 3. Malgré la pluie, le mariage a eu lieu devant la mairie. 4. Mon frère est doué en mathématiques. 5. Au nom de toute la famille, je vous remercie.

B. Corrigez les erreurs : 1. \eng{She is interested on music.} 2. \eng{They depend of their uncle.} 3. \eng{Thank you for help me.} 4. \eng{He entered in the room.} 5. \eng{I am angry to my brother.}
\end{exercice}

\begin{corrige}[Exercice : À vous de jouer]
A. 1. \eng{We listen to the news on the radio.} 2. \eng{Thank you for coming to the party.} 3. \eng{In spite of the rain, the wedding took place in front of the town hall.} 4. \eng{My brother is good at mathematics.} 5. \eng{On behalf of the whole family, I thank you.}

B. 1. \eng{She is interested in music.} (\eng{interested in}). 2. \eng{They depend on their uncle.} (\eng{depend on}). 3. \eng{Thank you for helping me.} (gérondif après \eng{for}). 4. \eng{He entered the room.} (pas de préposition après \eng{to enter}). 5. \eng{I am angry with my brother.} (\eng{angry with} quelqu'un).
\end{corrige}

\begin{savaistu}[Pourquoi tant de prépositions en anglais ?]
L'anglais exprime par des prépositions ce que le français exprime souvent par le verbe lui-même : \eng{to look for} = « chercher », \eng{to look after} = « garder », \eng{to put up with} = « supporter ». C'est pourquoi les élèves anglophones apprennent les prépositions « collées » aux verbes, exactement comme vous apprenez le genre d'un nom français avec son article. Bonne habitude : notez toujours un verbe nouveau \textbf{avec sa préposition} — \eng{to depend on}, \eng{to succeed in} — et non le verbe seul.
\end{savaistu}

\begin{aretenir}[L'essentiel de la section]
\begin{itemize}
\item Une prepositional phrase = préposition + nom, pronom ou gérondif : \eng{at school}, \eng{with them}, \eng{before leaving}.
\item Fonctions : lieu (\eng{in Yaoundé}), temps (\eng{on Mondays}, \eng{at 5 p.m.}), manière (\eng{by bus}, \eng{with care}), complément d'adjectif (\eng{proud of}, \eng{good at}), complément de verbe prépositionnel (\eng{depend on}, \eng{listen to}).
\item Après une préposition, le verbe est toujours en -ing : \eng{Thank you for coming.}
\item Apprenez les constructions figées avec leur préposition : \eng{listen to}, \eng{wait for}, \eng{belong to}, \eng{depend on}, \eng{apologise for}, \eng{succeed in}, \eng{insist on} ; \eng{afraid of}, \eng{fond of}, \eng{angry with}, \eng{responsible for}, \eng{different from}.
\item Prépositions complexes : \eng{in front of}, \eng{in spite of}, \eng{because of}, \eng{instead of}, \eng{on behalf of}.
\item Le groupe prépositionnel ne se place jamais entre le verbe et son complément direct.
\end{itemize}
\end{aretenir}

\coursec{Comparaison français/anglais et pièges}

Dans la section précédente, nous avons étudié les \cle{prepositional phrases} et constaté que l'anglais associe à certains verbes, adjectifs et noms une préposition fixe (\eng{depend on}, \eng{fond of}, \eng{in charge of}). Nous allons maintenant confronter directement les deux langues. C'est ici que se joue l'essentiel de la réussite d'un élève francophone : le français et l'anglais ne découpent pas la réalité de la même façon, et la traduction mot à mot est le premier ennemi du candidat au baccalauréat.

\courssub{Rappel : les quatre types de phrasal verbs}

Avant d'aborder les pièges de traduction, il est indispensable de maîtriser la classification qui conditionne la place du pronom (voir Piège 2).

\begin{definition}[Les quatre types de phrasal verbs]
\begin{enumerate}
  \item \textbf{Type 1 — Intransitifs} : pas d'objet. Ex. \eng{The car broke down.} « La voiture est tombée en panne. »
  \item \textbf{Type 2 — Transitifs séparables} : verbe + particule + objet ; l'objet \textbf{nom} peut aller avant ou après la particule, l'objet \textbf{pronom} va \textbf{obligatoirement} entre les deux. Ex. \eng{Turn off the light / Turn the light off / Turn it off.}
  \item \textbf{Type 3 — Transitifs inséparables} : verbe + préposition + objet ; l'objet (nom ou pronom) suit toujours la préposition. Ex. \eng{Look after the children / Look after them.} (jamais \eng{look them after}).
  \item \textbf{Type 4 — À trois mots (verbe + particule + préposition)} : toujours inséparables. Ex. \eng{put up with}, \eng{look forward to}, \eng{get on with}. L'objet suit l'ensemble : \eng{put up with it.}
\end{enumerate}
\end{definition}

\courssub{Le français n'a pas de phrasal verbs}

\begin{definition}[Une différence fondamentale]
Le français exprime généralement une action par \textbf{un seul verbe} (\emph{éteindre}, \emph{élever}, \emph{retarder}), tandis que l'anglais utilise souvent un \textbf{verbe + une particule} (\eng{turn off}, \eng{bring up}, \eng{put off}). Le phrasal verb anglais forme un \cle{bloc de sens} : la particule change complètement le sens du verbe de base.
\end{definition}

Observez ces paires avant de lire la règle :

\begin{exemplebox}[Un verbe français = verbe + particule en anglais]
\begin{itemize}
  \item \eng{Please turn off the lights before leaving.} « Éteins les lumières avant de partir. »
  \item \eng{She brought up three children on her own.} « Elle a élevé trois enfants toute seule. »
  \item \eng{The match was put off because of the rain.} « Le match a été reporté à cause de la pluie. »
  \item \eng{They made up after their quarrel.} « Ils se sont réconciliés après leur dispute. »
  \item \eng{I can't put up with his behaviour any longer.} « Je ne supporte plus son comportement. »
\end{itemize}
\end{exemplebox}

\begin{propriete}[Règle de correspondance français–anglais]
Il n'existe \textbf{aucune règle de traduction mécanique} entre un verbe français et un phrasal verb anglais. La particule (\eng{up}, \eng{off}, \eng{out}, \eng{on}...) ne se traduit pas : elle fait partie du mot. \eng{Turn} seul signifie « tourner », mais \eng{turn off} signifie « éteindre » et \eng{turn up} signifie « arriver » ou « monter (le son) ». Chaque phrasal verb doit donc être mémorisé comme \textbf{un vocabulaire entier}, exactement comme on apprend un mot nouveau.
\end{propriete}

\begin{popfigure}
\begin{center}
\begin{tikzpicture}[
  fr/.style={rectangle, rounded corners, draw=popblue, fill=popblueL, align=center, minimum width=4.6cm, minimum height=0.9cm, font=\small},
  en/.style={rectangle, rounded corners, draw=poporange, fill=poporangeL, align=center, minimum width=4.6cm, minimum height=0.9cm, font=\small},
  fleche/.style={-{Stealth[length=2.5mm]}, thick, popdark}
]
\node[fr] (f1) at (0,0) {éteindre la lumière};
\node[en] (e1) at (7.2,0) {turn the light off};
\node[fr] (f2) at (0,-1.3) {se réconcilier};
\node[en] (e2) at (7.2,-1.3) {make up};
\node[fr] (f3) at (0,-2.6) {assister à (la réunion)};
\node[en, align=center] (e3) at (7.2,-2.6){attend (the meeting)\\ \emph{sans « to » !}};
\node[fr] (f4) at (0,-3.9) {élever des enfants};
\node[en] (e4) at (7.2,-3.9) {bring up children};
\node[fr] (f5) at (0,-5.2) {remettre à plus tard};
\node[en] (e5) at (7.2,-5.2) {put off};
\draw[fleche] (f1) -- (e1);
\draw[fleche] (f2) -- (e2);
\draw[fleche] (f3) -- (e3);
\draw[fleche] (f4) -- (e4);
\draw[fleche] (f5) -- (e5);
\node[font=\bfseries\small, popblue] at (0,0.9) {Le français dit...};
\node[font=\bfseries\small, poporange] at (7.2,0.9) {L'anglais dit...};
\end{tikzpicture}
\end{center}
\legende{Un seul mot en français correspond souvent à un bloc verbe + particule en anglais.}
\end{popfigure}

\courssub{Piège 1 : traduire mot à mot le verbe seul}

L'erreur la plus fréquente consiste à traduire le verbe français par le verbe anglais de base, en oubliant la particule qui porte le sens.

\begin{attention}[To look $\neq$ s'occuper de]
\eng{Look} seul signifie « regarder ». Pour « s'occuper de », il faut \eng{look after} ; pour « chercher », \eng{look for} ; pour « admirer / respecter », \eng{look up to}.
\begin{itemize}
  \item \eng{She looks after her little brother.} « Elle s'occupe de son petit frère. » (\textbf{jamais} \eng{she looks her brother})
  \item \eng{I am looking for my keys.} « Je cherche mes clés. » (\textbf{jamais} \eng{I am looking my keys})
\end{itemize}
\end{attention}

\begin{attention}[To attend = assister à, sans préposition]
Attention au faux mouvement inverse : ici, c'est le français qui a une préposition (« assister \emph{à} ») et l'anglais qui n'en a pas !
\begin{itemize}
  \item \eng{I attended the meeting yesterday.} « J'ai assisté à la réunion hier. »
  \item Ne dites \textbf{jamais} \eng{I attended to the meeting}. \eng{Attend to} existe, mais signifie « s'occuper de, régler » : \eng{The nurse attended to the patient.} « L'infirmière s'est occupée du patient. »
\end{itemize}
\end{attention}

\courssub{Piège 2 : le placement du pronom}

\begin{propriete}[Le pronom se place au milieu]
Avec les phrasal verbs \textbf{séparables} (transitifs, Type 2), un \textbf{nom} peut se placer avant ou après la particule, mais un \textbf{pronom} (\eng{it}, \eng{them}, \eng{her}...) se place \textbf{obligatoirement entre le verbe et la particule}.
\begin{itemize}
  \item \eng{Turn off the light.} = \eng{Turn the light off.} (les deux sont corrects)
  \item \eng{Turn it off.} « Éteins-la. » — mais \textbf{jamais} \eng{turn off it}.
  \item \eng{I picked it up.} « Je l'ai ramassé. » — \textbf{jamais} \eng{I picked up it}.
\end{itemize}
En revanche, avec les verbes \textbf{inséparables} (Type 3 et Type 4 : \eng{look after}, \eng{look for}, \eng{run into}, \eng{put up with}, \eng{look forward to}), le pronom suit toujours la préposition : \eng{I am looking for them.} « Je les cherche. » ; \eng{I look forward to it.} « J'ai hâte de cela. »
\end{propriete}

\courssub{Piège 3 : les prépositions calquées du français}

C'est le piège classique des prepositional phrases : la préposition anglaise n'est presque jamais la traduction littérale de la préposition française.

\begin{exemplebox}[Exemples traduits à mémoriser]
\begin{itemize}
  \item \eng{She is married to a teacher.} « Elle est mariée à un professeur. » (\textbf{pas} \eng{married with})
  \item \eng{It depends on the weather.} « Cela dépend du temps. » (\textbf{pas} \eng{depends of})
  \item \eng{He is fond of music.} « Il aime la musique. »
  \item \eng{Are you interested in football?} « Est-ce que tu t'intéresses au football ? » (\textbf{pas} \eng{interested by})
  \item \eng{He apologized for being late.} « Il s'est excusé d'être en retard. »
\end{itemize}
\end{exemplebox}

\begin{center}
\begin{tabularx}{\linewidth}{|X|X|X|}
\hline
\rowcolor{popblueL} Français & English & Erreur fréquente \\
\hline
marié à & married \textbf{to} & married with \\
\hline
dépendre de & depend \textbf{on} & depend of \\
\hline
s'intéresser à & be interested \textbf{in} & interested by \\
\hline
rêver de & dream \textbf{of / about} & dream to \\
\hline
réussir à & succeed \textbf{in} & succeed to \\
\hline
appartenir à & belong \textbf{to} & belong at \\
\hline
s'excuser de & apologize \textbf{for} & apologize of \\
\hline
écouter (la radio) & listen \textbf{to} (the radio) & listen the radio \\
\hline
attendre (le bus) & wait \textbf{for} (the bus) & wait the bus \\
\hline
entrer dans (la salle) & enter (the room), \emph{sans prép.} & enter in \\
\hline
\end{tabularx}
\end{center}

\courssub{Piège 4 : l'infinitif après une préposition}

\begin{propriete}[Préposition + verbe en -ing]
En anglais, \textbf{tout verbe placé après une préposition} se met à la forme en \cle{-ing} (gérondif), jamais à l'infinitif avec \eng{to}. Le français, lui, utilise l'infinitif après une préposition (« merci d'\emph{être} venu »).
\begin{itemize}
  \item \eng{Thank you for coming.} « Merci d'être venu. » (\textbf{pas} \eng{for to come})
  \item \eng{He left without saying goodbye.} « Il est parti sans dire au revoir. »
  \item \eng{She is good at cooking.} « Elle est bonne en cuisine. »
  \item \eng{They insisted on paying the bill.} « Ils ont insisté pour payer l'addition. »
\end{itemize}
\textbf{Exception apparente :} attention à \eng{to} ! Dans \eng{look forward to}, \eng{to} est une préposition : \eng{I look forward to seeing you.} « J'ai hâte de te voir. » (\textbf{pas} \eng{to see}). De même : \eng{be used to doing}, \eng{get used to doing}, \eng{object to doing}.
\end{propriete}

\courssub{Piège 5 : particule adverbiale ou préposition ?}

Le même mot (\eng{up}, \eng{off}, \eng{in}) peut être soit une \cle{particule adverbiale} (qui fait partie du phrasal verb), soit une \cle{préposition} (qui introduit un complément de lieu ou de temps).

\begin{propriete}[Le test de distinction]
\begin{itemize}
  \item \eng{I get up early.} « Je me lève tôt. » → \eng{up} est une \textbf{particule} : \eng{get up} est un verbe à part entière (« se lever »), suivi ici d'un complément de temps.
  \item \eng{I go up the stairs.} « Je monte l'escalier. » → \eng{up} est une \textbf{préposition} : elle introduit le nom \eng{the stairs} (complément de lieu).
\end{itemize}
\textbf{Test pratique :} si le mot est suivi d'un nom qui lui appartient (\eng{up the stairs}), c'est une préposition ; s'il modifie le sens du verbe (\eng{get up} = se lever, et non « obtenir vers le haut »), c'est une particule.
\end{propriete}

\courssub{Stratégie de mémorisation}

\begin{methode}[Apprendre les blocs indissociables]
\begin{enumerate}
  \item N'apprends jamais un verbe seul : mémorise le \textbf{verbe + sa préposition} comme un seul mot : \eng{depend ON}, \eng{succeed IN}, \eng{insist ON}, \eng{apologize FOR}, \eng{belong TO}.
  \item Pour chaque phrasal verb, note \textbf{une traduction française et une phrase complète} en contexte (ex. \eng{bring up} = élever → \eng{My grandmother brought up six children in Bafoussam.}).
  \item Vérifie si le verbe est \textbf{séparable ou inséparable} (Type 2 vs Type 3/4) : cela décide de la place du pronom.
  \item Interroge-toi régulièrement dans les deux sens : français → anglais et anglais → français.
  \item En rédaction, relis-toi en traquant les prépositions : c'est là que les correcteurs repèrent les calques du français.
\end{enumerate}
\end{methode}

\begin{exoresolu}[Corrige les erreurs des francophones]
Énoncé : chaque phrase contient une erreur typique d'un élève francophone. Corrige-la et justifie.
\begin{enumerate}
  \item \eng{She is married with a doctor from Douala.}
  \item \eng{I am looking my pen everywhere.}
  \item \eng{He picked up it and left the room.}
  \item \eng{Thank you for to help me with my homework.}
  \item \eng{We attended to the wedding last Saturday.}
\end{enumerate}
\tcblower
Solution détaillée :
\begin{enumerate}
  \item \eng{She is married \textbf{to} a doctor from Douala.} — « marié à » se dit \eng{married to}, jamais \eng{with}.
  \item \eng{I am looking \textbf{for} my pen everywhere.} — « chercher » = \eng{look for} ; \eng{look} seul signifie « regarder ».
  \item \eng{He picked \textbf{it} up and left the room.} — \eng{pick up} est séparable (Type 2) : le pronom se place entre le verbe et la particule.
  \item \eng{Thank you for \textbf{helping} me with my homework.} — après une préposition (\eng{for}), le verbe prend la forme en \eng{-ing}.
  \item \eng{We attended the wedding last Saturday.} — « assister à » = \eng{attend} sans préposition ; \eng{attend to} signifie « s'occuper de ».
\end{enumerate}
\end{exoresolu}

\begin{exercice}[À toi de jouer]
Traduis en anglais en faisant attention aux pièges étudiés :
\begin{enumerate}
  \item Elle s'occupe de ses grands-parents le week-end.
  \item Cela dépend de ta décision.
  \item Éteins-le avant de sortir. (le = le téléviseur, \eng{it})
  \item Il s'est excusé d'être arrivé en retard.
  \item J'ai assisté à la cérémonie à Yaoundé.
\end{enumerate}
\end{exercice}

\begin{corrige}[Exercice « À toi de jouer »]
\begin{enumerate}
  \item \eng{She looks after her grandparents at the weekend.}
  \item \eng{It depends on your decision.}
  \item \eng{Turn it off before going out.}
  \item \eng{He apologized for arriving late.}
  \item \eng{I attended the ceremony in Yaoundé.}
\end{enumerate}
\end{corrige}

\begin{aretenir}[L'essentiel de la section]
\begin{itemize}
  \item Le français n'a pas de phrasal verbs : un verbe français = souvent verbe + particule en anglais.
  \item Jamais de traduction mot à mot : \eng{look after} = s'occuper de, \eng{look for} = chercher, \eng{attend} = assister à (sans préposition).
  \item Pronom avec un verbe séparable (Type 2) : \eng{pick it up}, jamais \eng{pick up it}.
  \item Verbes inséparables (Type 3 et 4) : pronom après la préposition (\eng{look after them}, \eng{put up with it}).
  \item Prépositions fixes à mémoriser en bloc : \eng{married to}, \eng{depend on}, \eng{interested in}, \eng{fond of}.
  \item Après une préposition, le verbe est en \eng{-ing} : \eng{thank you for coming}.
  \item Particule (\eng{get up early}) $\neq$ préposition (\eng{go up the stairs}).
\end{itemize}
\end{aretenir}

\coursec{Entraînement guidé et réinvestissement}

Après avoir analysé les différences structurelles entre le français et l'anglais — notamment l'ordre des mots avec les pronoms, les \cle{faux amis} (comme \eng{look for} $\neq$ \emph{regarder} mais \emph{chercher}) et la non-séparabilité des \cle{verbes prépositionnels} — nous passons maintenant à la mise en œuvre active. Cette section propose un parcours progressif : de la reconnaissance en contexte à la production autonome, en passant par les transformations types du Baccalauréat. L'objectif est d'automatiser les bons réflexes pour que l'utilisation des \cle{phrasal verbs} et des \cle{locutions prépositionnelles} devienne naturelle à l'oral comme à l'écrit.

\courssub{Exercice 1 : Identification et classification dans un texte authentique}

\begin{textebox}[Extrait de journal intime — A Week in Molyko]
Last Monday, I \textbf{woke up} late because my alarm didn't \textbf{go off}. I had to \textbf{hurry up} to catch the bendskin to school. In the traffic jam, the driver \textbf{broke down} near the Total station. Everyone \textbf{got off} and waited. Luckily, a mechanic \textbf{came along} and \textbf{fixed it up} in ten minutes. 

At home, my younger sister \textbf{looks after} our grandmother while Mum is at the farm. She \textbf{takes after} her in patience. Yesterday, Granny \textbf{ran out of} her medicine, so I \textbf{went out} to \textbf{look for} the pharmacy on Molyko main street. I \textbf{ran into} an old friend, Peter, who \textbf{gave up} his studies last year to \textbf{set up} a small business. He \textbf{puts up with} difficult customers but \textbf{looks forward to} expanding his shop. 

Tonight, I must \textbf{hand in} my assignment. I \textbf{put it off} too long. I will \textbf{stay up} to finish it, but I hope I don't \textbf{fall asleep} at my desk!
\end{textebox}

\begin{exercice}[Identification et classification]
Lisez le texte ci-dessus. 
\begin{enumerate}
    \item Surlignez tous les \cle{phrasal verbs} et \cle{verbes prépositionnels} (verbes à trois particules inclus).
    \item Classez-les dans le tableau ci-dessous selon les quatre types étudiés :
    \begin{itemize}
        \item Type 1 : Intransitifs (pas de COD)
        \item Type 2 : Transitifs inséparables / Verbes prépositionnels (verbe + préposition + COD)
        \item Type 3 : Transitifs séparables / Phrasaux purs (verbe + particule + COD)
        \item Type 4 : Verbes à trois particules / Phrasaux-prépositionnels (verbe + particule + préposition + COD)
    \end{itemize}
    \item Pour chaque verbe de type 3, soulignez la particule et indiquez si le COD est un nom ou un pronom dans le texte.
\end{enumerate}
\end{exercice}

\begin{corrige}[Exercice 1]
\begin{center}
\begin{tabularx}{\linewidth}{|X|X|X|X|}
\hline
\rowcolor{popblueL} \textbf{Type 1 : Intransitifs} & \textbf{Type 2 : Verbes prépositionnels (Inséparables)} & \textbf{Type 3 : Phrasaux séparables} & \textbf{Type 4 : Verbes à 3 particules (Inséparables)} \\
\hline
woke up & looks after & fixed \textbf{it} up & ran out of \\
go off & takes after & handed \textbf{it} in & put up with \\
broke down & ran into & put \textbf{it} off & looks forward to \\
got off & look for & set up (a business) & \\
came along & & gave up (his studies) & \\
went out & & & \\
stay up & & & \\
fall asleep & & & \\
\hline
\end{tabularx}
\end{center}

\textbf{Observations clés :}
\begin{itemize}
    \item \eng{woke up, go off, broke down, got off, came along, went out, stay up, fall asleep} : Type 1 (intransitifs, pas de COD).
    \item \eng{looks after, takes after, ran into, look for} : Type 2 (verbes prépositionnels). Le COD (nom ou pronom) suit \textbf{toujours} la préposition. Ex: \eng{looks after \textbf{her}}, \eng{look for \textbf{it}}.
    \item \eng{fixed \textbf{it} up, handed \textbf{it} in, put \textbf{it} off} : Type 3 (séparables). Le pronom \textbf{\eng{it}} est \textbf{obligatoirement} placé entre le verbe et la particule. Avec un COD nominal (\eng{set up a business, gave up his studies}), le nom se place \textbf{après} la particule (usage standard) ou entre les deux (\eng{set a business up} — possible mais moins courant pour les noms longs).
    \item \eng{ran out of, put up with, looks forward to} : Type 4 (trois particules). Structure Verbe + Particule + Préposition + COD. Inséparables. Ex: \eng{ran out of \textbf{it}}, \eng{put up with \textbf{them}}.
\end{itemize}
\end{corrige}

\courssub{Exercice 2 : Choix de la particule appropriée (Gap-fill)}

\begin{exercice}[Complétez avec la particule ou préposition correcte]
Complétez les phrases suivantes avec un mot parmi la liste : \eng{up, off, after, for, with, out, down, in, on, away, over}.
\begin{enumerate}
    \item The firemen managed to put \_\_\_\_\_ the fire before it spread to the market.
    \item My father wants me to take \_\_\_\_\_ the family cocoa farm when he retires.
    \item Don't let the baby play \_\_\_\_\_ the knives; put them \_\_\_\_\_ immediately.
    \item She takes \_\_\_\_\_ her mother; they have the same smile and the same temper.
    \item We have run \_\_\_\_\_ sugar. Could you go to the \emph{boutique} and buy some?
    \item The teacher asked the students to hand \_\_\_\_\_ their scripts at the end of the exam.
    \item I can't put \_\_\_\_\_ \_\_\_\_\_ this noise any longer; I'm going to the library.
    \item Look \_\_\_\_\_! That \emph{bendskin} is driving on the wrong side of the road.
    \item The meeting was called \_\_\_\_\_ because the Principal was sick.
    \item If you don't know the word, look it \_\_\_\_\_ in the dictionary.
\end{enumerate}
\end{exercice}

\begin{corrige}[Exercice 2]
\begin{enumerate}
    \item \textbf{out} (\eng{put out} = éteindre — Type 3, séparable).
    \item \textbf{over} (\eng{take over} = reprendre, succéder — Type 2, inséparable).
    \item \textbf{with} / \textbf{away} (\eng{play with} Type 2 ; \eng{put away} = ranger — Type 3).
    \item \textbf{after} (\eng{take after} = ressembler — Type 2, inséparable).
    \item \textbf{out of} (\eng{run out of} = manquer de — Type 4, deux mots).
    \item \textbf{in} (\eng{hand in} = rendre, soumettre — Type 3).
    \item \textbf{up with} (\eng{put up with} = supporter — Type 4, deux mots).
    \item \textbf{out} (\eng{Look out!} = Attention ! — Type 1, intransitif, impératif).
    \item \textbf{off} (\eng{call off} = annuler — Type 3).
    \item \textbf{up} (\eng{look up} = chercher (un mot) — Type 3, pronom \eng{it} inséré : \eng{look \textbf{it} up}).
\end{enumerate}
\end{corrige}

\courssub{Exercice 3 : Placement du pronom complément (Transformation \eng{it} / \eng{them})}

\begin{methode}[Méthode : Transformer un COD nom en pronom]
\begin{enumerate}
    \item \textbf{Identifiez} le phrasal verb et son type (séparable ou inséparable).
    \item \textbf{Remplacez} le groupe nominal (COD) par le pronom approprié (\eng{it} singulier, \eng{them} pluriel, \eng{him/her} pour personnes).
    \item \textbf{Placez} le pronom :
    \begin{itemize}
        \item \textbf{Entre} le verbe et la particule si le verbe est de \textbf{Type 3 (séparable)}.
        \item \textbf{Après} la préposition (ou la dernière préposition) si le verbe est de \textbf{Type 2 ou Type 4 (inséparables)}.
    \end{itemize}
    \item \textbf{Vérifiez} : ne séparez jamais un \cle{verbe prépositionnel} (Type 2) ni un \cle{verbe à trois particules} (Type 4).
\end{enumerate}
\end{methode}

\begin{popfigure}
\begin{tikzpicture}[
    node distance=1.3cm and 2.5cm,
    startstop/.style={rectangle, rounded corners, minimum width=3.2cm, minimum height=1cm, text centered, draw=popblue, fill=popblueL, font=\small},
    decision/.style={diamond, aspect=2, minimum width=3.5cm, minimum height=1.2cm, text centered, draw=poporange, fill=poporangeL, font=\small},
    process/.style={rectangle, minimum width=4cm, minimum height=1cm, text centered, draw=popgreen, fill=popgreenL, font=\small},
    arrow/.style={thick,->,>=Stealth, draw=popdark}
]
\node (start) [startstop, align=center]{1. Identifier\\ Verbe + Particule/Prép.};
\node (type) [decision, below of=start, align=center]{Est-ce un Type 3 ?\\(Séparable)};
\node (sep) [process, left of=type, xshift=-3cm, align=center]{OUI : Type 3\\ Pronom \textbf{entre} V et Part\\ \eng{turn \textbf{it} off}};
\node (insep) [process, right of=type, xshift=3cm, align=center]{NON : Type 2 ou 4\\ Pronom \textbf{après} la préposition\\ \eng{look for \textbf{it}}};
\node (end) [startstop, below of=type, yshift=-1.8cm] {Phrase correcte};

\draw [arrow] (start) -- (type);
\draw [arrow] (type) -- node[above, font=\tiny, sloped] {Oui} (sep);
\draw [arrow] (type) -- node[above, font=\tiny, sloped] {Non} (insep);
\draw [arrow] (sep) |- (end);
\draw [arrow] (insep) |- (end);
\end{tikzpicture}
\legende{Organigramme de décision pour le placement du pronom COD}
\end{popfigure}

\begin{exercice}[Transformation pronominale]
Remplacez le COD (en gras) par le pronom \eng{it}, \eng{them}, \eng{him} ou \eng{her} et réécrivez la phrase correctement.
\begin{enumerate}
    \item She switched \textbf{the light} off.
    \item They are looking for \textbf{their keys}.
    \item He gave up \textbf{smoking}.
    \item We must fill in \textbf{the form}.
    \item She looks after \textbf{her little brother}.
    \item The students handed in \textbf{their assignments}.
    \item I ran into \textbf{my old teacher} at the market.
    \item They put off \textbf{the meeting}.
    \item We look forward to \textbf{the holidays}.
    \item He takes after \textbf{his grandfather}.
\end{enumerate}
\end{exercice}

\begin{corrige}[Exercice 3]
\begin{enumerate}
    \item She switched \textbf{it} off. (Type 3 : séparable)
    \item They are looking for \textbf{them}. (Type 2 : inséparable, \eng{for} préposition)
    \item He gave \textbf{it} up. (Type 3 : séparable, \eng{smoking} = \eng{it})
    \item We must fill \textbf{it} in. (Type 3 : séparable)
    \item She looks after \textbf{him}. (Type 2 : inséparable, \eng{after} préposition $\rightarrow$ pronom après. \textit{NB : personne $\rightarrow$ \eng{him}/\eng{her}})
    \item The students handed \textbf{them} in. (Type 3 : séparable, pluriel $\rightarrow$ \eng{them})
    \item I ran into \textbf{him} at the market. (Type 2 : inséparable, \eng{into} préposition)
    \item They put \textbf{it} off. (Type 3 : séparable)
    \item We look forward to \textbf{them}. (Type 4 : inséparable, \eng{to} préposition)
    \item He takes after \textbf{him}. (Type 2 : inséparable, \eng{after} préposition)
\end{enumerate}
\end{corrige}

\courssub{Exercice 4 : Transformation type Baccalauréat (Remplacement lexical)}

\begin{exoresolu}[Modèle de transformation type Bac]
\textbf{Consigne :} Remplacez le verbe souligné par un \cle{phrasal verb} ou une \cle{locution verbale prépositionnelle} de sens équivalent. \textbf{Ne changez ni le temps, ni le sens, ni la voix de la phrase.}

\begin{enumerate}
    \item \textbf{Énoncé :} The government \uline{increased} fuel prices last week.
    \item \textbf{Analyse :} Temps = \emph{Past Simple}. Verbe = \eng{increase} (transitif). Sujet = \eng{The government}. COD = \eng{fuel prices}.
    \item \textbf{Choix :} \eng{put up} (augmenter, Type 3 séparable) ou \eng{hike up} (familier). \eng{Raise} n'est pas un phrasal verb.
    \item \textbf{Application règle Type 3 :} COD nom \eng{fuel prices} $\rightarrow$ placé \textbf{après} la particule (forme standard) ou entre les deux. Forme standard : \eng{put up fuel prices}.
    \item \textbf{Solution :} The government \textbf{put up} fuel prices last week. (ou \eng{put fuel prices up}).
\end{enumerate}
\end{exoresolu}

\begin{attention}[Piège majeur au Baccalauréat]
Dans les items de \textbf{transformation de phrase} (Sentence Transformation) :
\begin{itemize}
    \item \textbf{Ne changez JAMAIS le temps du verbe.} Si l'original est au \emph{Present Perfect}, le phrasal verb doit être au \emph{Present Perfect}.
    \item \textbf{Ne changez JAMAIS la voix.} Actif $\rightarrow$ Actif ; Passif $\rightarrow$ Passif.
    \item \textbf{Respectez la séparabilité.} Si le COD original est un pronom (\eng{it, them}), le phrasal verb choisi \textbf{doit être séparable} (Type 3). Ex: \eng{He abandoned it} $\rightarrow$ \eng{He gave it up} (correct). \eng{He gave up it} (incorrect, 0 point).
    \item \textbf{Évitez les verbes à trois particules (Type 4)} si le COD est un pronom, sauf si la structure le permet (\eng{put up with it} est correct car inséparable).
    \item \textbf{Attention au gérondif :} Après \eng{look forward to}, \eng{be used to}, \eng{get used to}, \eng{object to} $\rightarrow$ \textsc{v-ing} (\eng{seeing}, \eng{doing}).
\end{itemize}
\end{attention}

\begin{exercice}[Transformation — Niveau Bac]
Transformez les phrases en remplaçant le verbe souligné par un phrasal verb approprié. Conservez le temps et la structure.
\begin{enumerate}
    \item The children \uline{resemble} their parents in many ways.
    \item The Principal \uline{cancelled} the sports day due to rain.
    \item We must \uline{tolerate} his bad behaviour for now.
    \item She \uline{investigated} the cause of the accident.
    \item They \uline{abandoned} the project for lack of funds.
    \item I \uline{meet} my friends every Saturday at the \emph{buvette}.
    \item The fire \uline{destroyed} the whole warehouse.
    \item You should \uline{review} your lessons before the exam.
    \item He \uline{recovered from} malaria quickly.
    \item The bomb \uline{exploded} in the market.
\end{enumerate}
\end{exercice}

\begin{corrige}[Exercice 4]
\begin{enumerate}
    \item The children \textbf{take after} their parents... (Type 2, inséparable, \eng{resemble} $\rightarrow$ \eng{take after})
    \item The Principal \textbf{called off} the sports day... (Type 3, séparable, \eng{cancel} $\rightarrow$ \eng{call off})
    \item We must \textbf{put up with} his bad behaviour... (Type 4, inséparable, \eng{tolerate} $\rightarrow$ \eng{put up with})
    \item She \textbf{looked into} the cause... (Type 2, inséparable, \eng{investigate} $\rightarrow$ \eng{look into})
    \item They \textbf{gave up} the project... (Type 3, séparable, \eng{abandon} $\rightarrow$ \eng{give up})
    \item I \textbf{meet up with} my friends... (Type 4, inséparable, \eng{meet} $\rightarrow$ \eng{meet up with})
    \item The fire \textbf{burnt down} the whole warehouse. (Type 3, séparable, \eng{destroy} $\rightarrow$ \eng{burn down})
    \item You should \textbf{go over} your lessons... (Type 2, inséparable, \eng{review} $\rightarrow$ \eng{go over})
    \item He \textbf{got over} malaria quickly. (Type 2, inséparable, \eng{recover from} $\rightarrow$ \eng{get over})
    \item The bomb \textbf{went off} in the market. (Type 1, intransitif, \eng{explode} $\rightarrow$ \eng{go off})
\end{enumerate}
\end{corrige}

\courssub{Exercice 5 : Thème ciblé (Français $\rightarrow$ Anglais)}

\begin{exercice}[Traduction — Pièges francophones]
Traduisez en anglais en utilisant obligatoirement un \cle{phrasal verb} ou une \cle{locution verbale prépositionnelle} souligné(e) dans la correction.
\begin{enumerate}
    \item Mon grand-père s'occupe de mes cousins pendant les vacances. (\textit{indice : look after})
    \item Elle a annulé son voyage à Douala à la dernière minute. (\textit{indice : call off})
    \item Nous n'avons plus d'essence ; il faut s'arrêter à la station. (\textit{indice : run out of})
    \item Il supporte difficilement les critiques de son patron. (\textit{indice : put up with})
    \item J'ai hâte de revoir mes camarades de classe. (\textit{indice : look forward to})
    \item Le professeur a distribué les copies au début du cours. (\textit{indice : hand out / give out})
    \item N'éteins pas la lumière, je lis encore. (\textit{indice : turn off} — attention pronom !)
    \item Il a arrêté de fumer l'année dernière. (\textit{indice : give up})
    \item Cette couleur ne va pas avec ton pagne. (\textit{indice : go with})
    \item Le voleur a réussi à s'enfuir par la fenêtre. (\textit{indice : get away})
\end{enumerate}
\end{exercice}

\begin{corrige}[Exercice 5]
\begin{enumerate}
    \item My grandfather \textbf{looks after} my cousins during the holidays. (Type 2 : \eng{after} préposition $\rightarrow$ \eng{look after \textbf{them}} si pronom).
    \item She \textbf{called off} her trip to Douala at the last minute. (Type 3 : \eng{called \textbf{it} off} si pronom).
    \item We have \textbf{run out of} petrol; we must stop at the station. (Type 4 : \eng{run out of \textbf{it}}).
    \item He hardly \textbf{puts up with} his boss's criticism. (Type 4 : \eng{puts up with \textbf{it}}).
    \item I \textbf{look forward to} seeing my classmates again. (Type 4 : \eng{look forward to \textbf{it / doing it}}. \textbf{Piège :} \eng{to} est préposition $\rightarrow$ \textsc{v-ing}).
    \item The teacher \textbf{handed out} the scripts at the beginning of the lesson. (Type 3 : \eng{handed \textbf{them} out}).
    \item Don't \textbf{turn it off}, I'm still reading. (\textbf{Piège majeur :} Type 3 séparable + pronom \eng{it} $\rightarrow$ \textbf{obligatoirement} \eng{turn \textbf{it} off}. \eng{Turn off it} = \textbf{Erreur éliminatoire}).
    \item He \textbf{gave up} smoking last year. (Type 3 : \eng{gave \textbf{it} up}).
    \item This colour doesn't \textbf{go with} your wrapper. (Type 2 : \eng{go with \textbf{it}}).
    \item The thief managed to \textbf{get away} through the window. (Type 1 : intransitif).
\end{enumerate}
\end{corrige}

\courssub{Production finale : Ma routine familiale}

\begin{experience}[Production écrite guidée — 6 à 8 lignes]
\textbf{Consigne :} Rédigez un court paragraphe (environ 80--100 mots) décrivant une journée typique dans votre famille (routine, tâches, interactions). 
\textbf{Contraintes obligatoires :}
\begin{itemize}
    \item Utilisez \textbf{au moins 5 phrasal verbs} différents (variez les types : au moins un Type 1, un Type 3, un Type 4).
    \item Utilisez \textbf{au moins 3 locutions prépositionnelles} (ex: \eng{in front of, because of, according to, in spite of, instead of, on time}).
    \item Respectez la chronologie (marqueurs : \eng{First, Then, After that, Finally}).
    \item Vérifiez le placement des pronoms COD (\eng{it, them, him, her}) avec les verbes séparables (Type 3).
\end{itemize}

\end{experience}

\begin{exemplebox}[Modèle de production attendue]
\textbf{A Typical Day in My Family}

“\eng{First, I \textbf{wake up} early at 5:30 a.m. and \textbf{tidy up} my room before breakfast. My mother \textbf{looks after} the baby while my father \textbf{goes out} to buy bread. \textbf{Because of} the traffic, he sometimes \textbf{comes back} late. \textbf{In spite of} the noise, I \textbf{get on with} my homework. At noon, we \textbf{sit down} to eat together. In the evening, my elder brother \textbf{helps me out} with maths. I \textbf{look forward to} the weekend when we \textbf{go out} as a family.}”
\par
\textbf{Analyse :} \cle{Phrasal verbs} : \eng{wake up} (T1), \eng{tidy up} (T3), \eng{looks after} (T2), \eng{goes out} (T1), \eng{comes back} (T1), \eng{get on with} (T4), \eng{sit down} (T1), \eng{helps out} (T3), \eng{look forward to} (T4), \eng{go out} (T1). \cle{Locutions prépositionnelles} : \eng{Because of}, \eng{In spite of}, \eng{as a family}. Pronoms : \eng{helps \textbf{me} out} (correct, Type 3).
\end{exemplebox}


\begin{savaistu}[Le phrasal verb au Baccalauréat Série A]
Les épreuves de \textbf{Grammaire} et de \textbf{Structure} du Baccalauréat A (Épreuve écrite, Partie B) comportent systématiquement 2 à 4 items sur les \cle{phrasal verbs} et \cle{locutions verbales prépositionnelles}.
\begin{itemize}
    \item \textbf{Type d'items :} \emph{Fill in the blanks} (choix de la particule), \emph{Sentence transformation} (remplacement verbe simple $\rightarrow$ phrasal verb), \emph{Error correction} (détection erreur placement pronom / mauvaise particule).
    \item \textbf{Le « Top 20 » stratégique :} Maîtriser parfaitement les 20 verbes ci-dessous couvre 85\% des occurrences aux examens : \eng{look for, look after, look forward to, look up, give up, put up with, put off, put out, take after, take off, take up, turn up, turn down, turn off, get up, get over, get on with, go on, go off, run out of}.
    \item \textbf{Astuce de révision :} Apprenez-les par \textbf{thèmes} (Famille : \eng{look after, take after, get on with} ; École/Travail : \eng{hand in, put off, go over, drop out} ; Vie quotidienne : \eng{wake up, run out of, put up with, set off}).
\end{itemize}
\end{savaistu}

\begin{center}
\begin{tabularx}{\linewidth}{|X|X|X|X|}
\hline
\rowcolor{popblueL} \textbf{Phrasal Verb / Locution} & \textbf{Type} & \textbf{Sens français} & \textbf{Exemple clé (COD Prénom)} \\
\hline
\eng{look after} & 2 (Prep) & s'occuper de & \eng{look after \textbf{him}} \\
\eng{take after} & 2 (Prep) & ressembler à & \eng{take after \textbf{her}} \\
\eng{run out of} & 4 (3-part) & manquer de & \eng{run out of \textbf{it}} \\
\eng{put up with} & 4 (3-part) & supporter & \eng{put up with \textbf{it}} \\
\eng{look forward to} & 4 (3-part) & avoir hâte de & \eng{look forward to \textbf{it / doing it}} \\
\eng{give up} & 3 (Sep) & abandonner / arrêter & \eng{give \textbf{it} up} \\
\eng{put off} & 3 (Sep) & reporter & \eng{put \textbf{it} off} \\
\eng{call off} & 3 (Sep) & annuler & \eng{call \textbf{it} off} \\
\eng{hand in} & 3 (Sep) & rendre (devoir) & \eng{hand \textbf{it} in} \\
\eng{turn off/on} & 3 (Sep) & éteindre/allumer & \eng{turn \textbf{it} off} \\
\eng{look for} & 2 (Prep) & chercher & \eng{look for \textbf{them}} \\
\eng{run into} & 2 (Prep) & rencontrer par hasard & \eng{run into \textbf{him}} \\
\eng{wake up} & 1 (Intrans) & se réveiller & — \\
\eng{go off} & 1 (Intrans) & sonner / exploser & — \\
\eng{break down} & 1 (Intrans) & tomber en panne & — \\
\hline
\end{tabularx}
\end{center}

\begin{aretenir}[Check-list finale avant l'examen]
\begin{itemize}
    \item[$\square$] Je sais distinguer les 4 types structurels (Intransitif, Prépositionnel, Séparable, 3 particules).
    \item[$\square$] Je place \textbf{obligatoirement} le pronom COD (\eng{it, them, him, her}) \textbf{entre} V et Particule pour les Type 3.
    \item[$\square$] Je ne sépare \textbf{jamais} les Type 2 (\eng{look for it}) ni Type 4 (\eng{put up with it}).
    \item[$\square$] Après \eng{look forward to}, \eng{be used to}, \eng{get used to}, \eng{object to} $\rightarrow$ \textsc{v-ing} (\eng{seeing}, \eng{doing}).
    \item[$\square$] En transformation (Bac), je garde le \textbf{même temps} et la \textbf{même voix}.
    \item[$\square$] Je connais le sens des 20 phrasal verbs « haute fréquence » du Bac A.
\end{itemize}
\end{aretenir}

\newpage

\coursec{Activité d'intégration}

\coursec{Lesson 10 — Grammar: Phrasal verbs and prepositional phrases}

\courssub{Observation et découverte}

\begin{textebox}[Dialogue : A Morning Routine in Yaoundé]
\eng{— Good morning, Mama! What time did you wake up?} « Bonjour Maman ! À quelle heure t'es-tu réveillée ? » \\
\eng{— I woke up at 5 a.m. I had to get up early to prepare the beans. Can you help me tidy up the kitchen?} « Je me suis réveillée à 5 h. J'ai dû me lever tôt pour préparer les haricots. Peux-tu m'aider à ranger la cuisine ? » \\
\eng{— Sure. I'll put on my apron and pick up the dishes. Where should I put them?} « Bien sûr. Je vais mettre mon tablier et ramasser la vaisselle. Où dois-je la mettre ? » \\
\eng{— Put them on the shelf, please. And don't forget to look after your little brother while I go out to the market.} « Mets-les sur l'étagère, s'il te plaît. Et n'oublie pas de t'occuper de ton petit frère pendant que je vais au marché. » \\
\eng{— Okay. I get on well with him, so it's easy. He usually grows up fast!} « D'accord. Je m'entends bien avec lui, donc c'est facile. Il grandit si vite ! »
\end{textebox}

\begin{methode}[Comment observer les phrasal verbs ?]
\begin{enumerate}
\item \textbf{Surlignez} dans le dialogue tous les groupes \textbf{verbe + particule} (ex. \eng{wake up}, \eng{get up}, \eng{tidy up}, \eng{put on}, \eng{pick up}, \eng{put on}, \eng{look after}, \eng{go out}, \eng{get on with}, \eng{grow up}).
\item \textbf{Classez-les} en deux colonnes : ceux où un mot peut s'intercaler entre le verbe et la particule (\eng{put \textbf{the dishes} on} / \eng{put \textbf{them} on}) et ceux où c'est impossible (\eng{look after \textbf{him}}, jamais \eng{look \textbf{him} after}).
\item \textbf{Repérez} les groupes \textbf{préposition + nom} qui indiquent le temps, le lieu ou la manière (\eng{at 5 a.m.}, \eng{in the kitchen}, \eng{on the shelf}, \eng{to the market}, \eng{with him}).
\end{enumerate}
\end{methode}

\begin{exemplebox}[Analyse guidée]
\begin{itemize}
\item \eng{wake up}, \eng{get up}, \eng{go out}, \eng{grow up} $\rightarrow$ \textbf{Intransitifs} (pas de complément d'objet).
\item \eng{tidy up}, \eng{put on}, \eng{pick up} $\rightarrow$ \textbf{Transitifs séparables} (complément nom après la particule \textbf{ou} pronom \textbf{entre} verbe et particule).
\item \eng{look after}, \eng{get on with} $\rightarrow$ \textbf{Transitifs inséparables} (complément toujours \textbf{après} la particule/préposition).
\end{itemize}
\end{exemplebox}

\courssub{Les quatre types de phrasal verbs}

\begin{definition}[Phrasal verb]
Un \cle{phrasal verb} est un verbe lexical suivi d'une \cle{particule} (adverbe prépositionnel : \eng{up, down, on, off, out, away...}) qui modifie le sens du verbe de base. Le sens global est souvent \textbf{idiomatique} (non déductible de la somme des parties).
\end{definition}

\begin{propriete}[Les quatre types structurels]
\begin{center}
\begin{tabularx}{\linewidth}{|X|X|X|X|}
\hline
\rowcolor{popblueL} \textbf{Type} & \textbf{Structure} & \textbf{Placement du COD (nom)} & \textbf{Placement du COD (pronom)} \\
\hline
\textbf{Type 1 : Intransitif} & V + Particule & --- & --- \\
\hline
\textbf{Type 2 : Transitif séparable} & V + Particule + COD & V + COD + Particule \textbf{OU} V + Particule + COD & \textbf{V + Pronom + Particule} (obligatoire) \\
\hline
\textbf{Type 3 : Transitif inséparable (Prepositional Verb)} & V + Préposition + COD & V + Préposition + COD & V + Préposition + Pronom \\
\hline
\textbf{Type 4 : Trois parties (Phrasal-Prepositional)} & V + Particule + Préposition + COD & V + Particule + Préposition + COD & V + Particule + Préposition + Pronom \\
\hline
\end{tabularx}
\end{center}
\end{propriete}

\begin{exemplebox}[Exemples des 4 types — Contexte familial]
\begin{tabular}{|l|l|l|}
\hline
\textbf{Type} & \textbf{Phrasal Verb} & \textbf{Exemple traduit} \\
\hline
1. Intransitif & \eng{wake up} (se réveiller) & \eng{The baby woke up early.} « Le bébé s'est réveillé tôt. » \\
\hline
1. Intransitif & \eng{go out} (sortir) & \eng{My parents go out on Sundays.} « Mes parents sortent le dimanche. » \\
\hline
2. Séparable & \eng{tidy up} (ranger) & \eng{I tidied up my room.} / \eng{I tidied it up.} « J'ai rangé ma chambre. / Je l'ai rangée. » \\
\hline
2. Séparable & \eng{put on} (mettre/enfiler) & \eng{She puts on her uniform.} / \eng{She puts it on.} « Elle met son uniforme. / Elle le met. » \\
\hline
3. Inséparable & \eng{look after} (s'occuper de) & \eng{He looks after his sister.} / \eng{He looks after her.} « Il s'occupe de sa sœur. / Il s'occupe d'elle. » \\
\hline
3. Inséparable & \eng{get on with} (s'entendre avec) & \eng{We get on well with them.} « Nous nous entendons bien avec eux. » \\
\hline
4. Trois parties & \eng{put up with} (supporter) & \eng{I can't put up with the noise.} « Je ne peux pas supporter le bruit. » \\
\hline
4. Trois parties & \eng{look forward to} (avoir hâte de) & \eng{She looks forward to the holidays.} « Elle a hâte des vacances. » \\
\hline
\end{tabular}
\end{exemplebox}

\begin{attention}[Piège : \eng{to} = particule ou préposition ?]
Avec les verbes à trois parties (Type 4), le dernier élément est une \textbf{préposition} suivie d'un \textbf{gérondif} (V-ing) si c'est un verbe, jamais d'un infinitif.
\eng{look forward to \textbf{seeing} you} (correct) — \eng{look forward to \textbf{see} you} (incorrect).
\end{attention}

\courssub{Phrasal verbs fréquents : vie familiale et sociale}

\begin{aretenir}[Lexique thématique — Family \& Social Life]
\begin{center}
\begin{tabularx}{\linewidth}{|l|X|l|X|}
\hline
\rowcolor{popblueL} \textbf{Phrasal Verb} & \textbf{Sens français} & \textbf{Type} & \textbf{Exemple contextuel (Cameroun)} \\
\hline
\eng{wake up} & se réveiller & 1 Intrans. & \eng{We wake up at dawn for morning prayers.} \\
\eng{get up} & se lever & 1 Intrans. & \eng{My father gets up at 5 a.m. for work.} \\
\eng{go out} & sortir & 1 Intrans. & \eng{Teenagers like going out with friends.} \\
\eng{grow up} & grandir & 1 Intrans. & \eng{Many children grow up in extended families.} \\
\hline
\eng{put on} & mettre (vêtement) & 2 Séparable & \eng{Put on your school uniform.} / \eng{Put it on.} \\
\eng{take off} & enlever (vêtement) & 2 Séparable & \eng{Take off your shoes at the door.} / \eng{Take them off.} \\
\eng{tidy up} & ranger, mettre en ordre & 2 Séparable & \eng{Please tidy up the living room.} / \eng{Tidy it up.} \\
\eng{pick up} & ramasser / aller chercher & 2 Séparable & \eng{I'll pick up the children from school.} / \eng{I'll pick them up.} \\
\eng{help out} & donner un coup de main & 2 Séparable & \eng{She helps out at the family shop.} / \eng{She helps out.} \\
\eng{turn off} & éteindre (lumière, gaz) & 2 Séparable & \eng{Turn off the gas after cooking.} / \eng{Turn it off.} \\
\hline
\eng{look after} & s'occuper de & 3 Inséparable & \eng{Aunties often look after nephews.} \\
\eng{get on with} & bien s'entendre avec & 3 Inséparable & \eng{Do you get on with your in-laws?} \\
\eng{run into} & rencontrer par hasard & 3 Inséparable & \eng{I ran into an old friend at the market.} \\
\eng{take after} & ressembler à (famille) & 3 Inséparable & \eng{He takes after his grandfather.} \\
\hline
\eng{put up with} & supporter, tolérer & 4 Trois parties & \eng{We put up with power cuts.} \\
\eng{look forward to} & avoir hâte de & 4 Trois parties & \eng{We look forward to the Christmas feast.} \\
\eng{run out of} & manquer de, épuiser & 4 Trois parties & \eng{We ran out of firewood.} \\
\hline
\end{tabularx}
\end{center}
\end{aretenir}

\begin{savaistu}[Le bendskin et la vie sociale]
Au Cameroun, \eng{to hop on a bendskin} (monter sur un taxi-moto) est une action quotidienne. On utilise souvent le phrasal verb \eng{get on} (monter dans/ sur un véhicule) et \eng{get off} (descendre de) : \eng{“Get on the bendskin, we're late!”} « Monte sur le bendskin, nous sommes en retard ! » — \eng{“Get off at the next junction.”} « Descends au prochain carrefour. »
\end{savaistu}

\courssub{Les prepositional phrases}

\begin{definition}[Prepositional phrase]
Une \cle{prepositional phrase} (groupe prépositionnel) est un groupe de mots commençant par une \textbf{préposition} et se terminant par un \textbf{nom} (ou pronom / gérondif). Elle fonctionne comme \textbf{complément circonstanciel} (temps, lieu, manière, cause) ou \textbf{complément d'objet} (après un verbe prépositionnel Type 3).
\end{definition}

\begin{propriete}[Principales prépositions de temps, lieu, manière — Pièges pour francophones]
\begin{center}
\begin{tabularx}{\linewidth}{|X|X|X|}
\hline
\rowcolor{popblueL} \textbf{Contexte} & \textbf{Anglais correct} & \textbf{Erreur fréquente (calque FR)} \\
\hline
\textbf{Temps : moment de la journée} & \eng{in the morning}, \eng{in the afternoon}, \eng{in the evening} & \eng{*at the morning}, \eng{*on the evening} \\
\textbf{Temps : heure précise} & \eng{at 6 o'clock}, \eng{at noon}, \eng{at midnight} & \eng{*in 6 o'clock} \\
\textbf{Temps : nuit (général)} & \eng{at night} & \eng{*in the night} (sauf nuit précise : \eng{in the night of July 14}) \\
\textbf{Temps : jour / date} & \eng{on Monday}, \eng{on 14 July}, \eng{on Christmas Day} & \eng{*at Monday}, \eng{*in 14 July} \\
\textbf{Lieu : ville / pays / quartier} & \eng{in Yaoundé}, \eng{in Cameroon}, \eng{in Mile 4} & \eng{*at Yaoundé} (sauf petit village/arrêt précis) \\
\textbf{Lieu : bâtiment / école / travail} & \eng{at school}, \eng{at home}, \eng{at work}, \eng{at the market} & \eng{*in school} (US), \eng{*in home} \\
\textbf{Lieu : moyen de transport} & \eng{on foot}, \eng{on a bike}, \eng{in a car}, \eng{in a taxi}, \eng{by bendskin}, \eng{by bus}, \eng{by train} & \eng{*by foot}, \eng{*with a car} \\
\textbf{Manière / Compagnie} & \eng{with my parents}, \eng{with a smile}, \eng{in silence} & \eng{*by my parents}, \eng{*with silence} \\
\hline
\end{tabularx}
\end{center}
\end{propriete}

\begin{exemplebox}[Prepositional phrases dans le récit familial]
\begin{itemize}
\item \eng{\textbf{In the morning}, I wake up \textbf{at 5 a.m.}} (Temps)
\item \eng{I go to school \textbf{on foot} \textbf{with my friend}.} (Moyen + Compagnie)
\item \eng{We eat dinner \textbf{at home} \textbf{in the evening}.} (Lieu + Temps)
\item \eng{She looks after the baby \textbf{with patience}.} (Manière)
\end{itemize}
\end{exemplebox}

\courssub{Comparaison français/anglais et pièges}

\begin{attention}[Faux amis et erreurs de traduction mot à mot]
\begin{center}
\begin{tabularx}{\linewidth}{|X|X|X|}
\hline
\rowcolor{poporangeL} \textbf{Français} & \textbf{Anglais correct (Phrasal Verb / Structure)} & \textbf{Erreur à bannir} \\
\hline
S'occuper de (qn) & \eng{look after} / \eng{take care of} & \eng{*occupy}, \eng{*take care} (sans \eng{of}) \\
Supporter (qn/qch) & \eng{put up with} & \eng{*support somebody} (sens : financer/encourager) \\
Ressembler à (famille) & \eng{take after} & \eng{*resemble to}, \eng{*look like} (physique seulement) \\
Se disputer / Se quereller & \eng{argue with} / \eng{fall out with} & \eng{*dispute}, \eng{*quarrel} (trop formel/rare) \\
Rencontrer par hasard & \eng{run into} / \eng{bump into} & \eng{*meet by chance} (lourd), \eng{*find} \\
Avoir hâte de & \eng{look forward to} + V-ing & \eng{*wait for}, \eng{*hope to} (sens différent) \\
Manquer de (ressource) & \eng{run out of} & \eng{*miss}, \eng{*lack} (sens : regretter / ne pas avoir assez) \\
\hline
\end{tabularx}
\end{center}
\end{attention}

\begin{attention}[Règles d'or du placement — Récapitulatif]
\begin{enumerate}
\item \textbf{Séparable + Nom :} \eng{put on your shoes} \textbf{OU} \eng{put your shoes on}. (Les deux corrects).
\item \textbf{Séparable + Pronom :} \eng{put \textbf{them} on} — \textbf{JAMAIS} \eng{*put on them}.
\item \textbf{Inséparable (V + Prep) :} \eng{look after \textbf{her}} — \textbf{JAMAIS} \eng{*look her after}.
\item \textbf{Trois parties :} \eng{put up with \textbf{it}} — \textbf{JAMAIS} \eng{*put it up with}.
\item \textbf{Particule accentuée à l'oral :} Dans \eng{put \textbf{ON}}, \eng{come \textbf{BACK}}, la particule porte l'accent tonique (stress).
\end{enumerate}
\end{attention}

\courssub{Entraînement guidé et réinvestissement}

\begin{exercice}[Exercice 1 : Classification et transformation]
Classez les phrasal verbs suivants dans le bon type (1, 2, 3 ou 4). Puis réécrivez la phrase en remplaçant le COD nom par le pronom indiqué.
\begin{enumerate}
\item \eng{She \textbf{looks after} her grandmother.} (her) $\rightarrow$ Type \_\_\_
\item \eng{They \textbf{tidied up} the compound.} (it) $\rightarrow$ Type \_\_\_
\item \eng{We \textbf{ran out of} petrol.} (it) $\rightarrow$ Type \_\_\_
\item \eng{The plane \textbf{took off} on time.} $\rightarrow$ Type \_\_\_
\item \eng{He \textbf{put on} his jacket.} (it) $\rightarrow$ Type \_\_\_
\item \eng{I \textbf{get on well with} my cousins.} (them) $\rightarrow$ Type \_\_\_
\end{enumerate}
\end{exercice}

\begin{corrige}[Exercice 1]
\begin{enumerate}
\item Type 3 (Inséparable). \eng{She looks after \textbf{her}.}
\item Type 2 (Séparable). \eng{They tidied \textbf{it} up.} (ou \eng{tidied up it} incorrect avec pronom).
\item Type 4 (Trois parties). \eng{We ran out of \textbf{it}.}
\item Type 1 (Intransitif). Pas de COD.
\item Type 2 (Séparable). \eng{He put \textbf{it} on.}
\item Type 4 (Trois parties : V + Part + Prep). \eng{I get on well with \textbf{them}.}
\end{enumerate}
\end{corrige}

\begin{exercice}[Exercice 2 : Choix de la préposition correcte]
Complétez avec la préposition adéquate (\eng{in, on, at, by, with, to, for}).
\begin{enumerate}
\item We usually have breakfast \_\_\_\_\_ 7 a.m. \_\_\_\_\_ the morning.
\item My brother goes to school \_\_\_\_\_ foot, but I go \_\_\_\_\_ bendskin.
\item \_\_\_\_\_ night, the compound is quiet. Everyone is \_\_\_\_\_ home.
\item She takes \_\_\_\_\_ her mother: they have the same smile.
\item We are looking forward \_\_\_\_\_ the end-of-year party.
\item Don't shout \_\_\_\_\_ your little sister; talk \_\_\_\_\_ her gently.
\end{enumerate}
\end{exercice}

\begin{corrige}[Exercice 2]
\begin{enumerate}
\item \eng{at} 7 a.m. \eng{in} the morning.
\item \eng{on} foot, \eng{by} bendskin.
\item \eng{At} night, \eng{at} home.
\item \eng{after} (phrasal verb \eng{take after}).
\item \eng{to} (\eng{look forward to}).
\item \eng{at} (shout \eng{at} = crier sur), \eng{to} (talk \eng{to} = parler à).
\end{enumerate}
\end{corrige}

\begin{exercice}[Exercice 3 : Correction d'erreurs (Production d'élèves)]
Corrigez les erreurs dans les phrases suivantes (grammaire, placement, préposition, faux ami).
\begin{enumerate}
\item I wake up at the morning and get up at 6 hours.
\item My mother occupies my baby brother when she cooks.
\item Please put on them on the table.
\item We support the noise of the generator every day.
\item He looks his father after in the village.
\item I am waiting for to see my friends this weekend.
\item We go to school with a taxi.
\item She takes off her shoes and puts off her socks.
\end{enumerate}
\end{exercice}

\begin{corrige}[Exercice 3]
\begin{enumerate}
\item \eng{in the morning} / \eng{at 6 o'clock} (pas \eng{at the morning}, pas \eng{6 hours}).
\item \eng{looks after} / \eng{takes care of} (pas \eng{occupy}).
\item \eng{Put them on} (pronom entre verbe et particule).
\item \eng{put up with} (pas \eng{support} = subvenir aux besoins / encourager).
\item \eng{looks after his father} (inséparable : complément après la particule).
\item \eng{looking forward to seeing} (\eng{to} préposition + V-ing).
\item \eng{by taxi} (moyen de transport \eng{by} + nom sans article, sauf \eng{on foot}).
\item \eng{takes off her shoes and takes off her socks} (\eng{put off} = reporter/remettre à plus tard ; \eng{take off} = enlever).
\end{enumerate}
\end{corrige}

\courssub{Activité d'intégration : Concours « A Day in My Family »}

\courssub{Objectifs de l'activité}

À l'issue de cette activité d'intégration, l'élève sera capable de :
\begin{itemize}
\item rédiger un court récit personnel en anglais (environ 8 lignes) en réinvestissant au moins \cle{6 phrasal verbs} et \cle{3 prepositional phrases} étudiés dans la leçon ;
\item placer correctement le complément d'objet des phrasal verbs séparables (\eng{put your shoes on} / \eng{put them on}) ;
\item repérer et corriger les erreurs de structure dans la production d'un camarade (relecture croisée) ;
\item poser des questions orales et y répondre en phrases complètes en utilisant des phrasal verbs de la vie familiale (\eng{get on with}, \eng{look after}, \eng{help out}, \eng{grow up}) ;
\item prendre la parole devant la classe avec une prononciation claire et un débit naturel.
\end{itemize}

\courssub{Situation}

\begin{textebox}[English Club Newsletter — Government High School, Buea]
\textbf{English Club — Writing Corner}

Dear members,

This month, our school magazine is preparing a special page called \textbf{“A Day in My Family”}. We want short, lively texts (about eight lines) in which students describe a typical day at home: waking up, helping their parents, going to school, sharing meals and relaxing in the evening.

\textbf{Rules of the competition:}
\begin{itemize}
\item Use at least \textbf{six phrasal verbs} (for example: \eng{get up}, \eng{look after}, \eng{tidy up}, \eng{go out}, \eng{come back}).
\item Use at least \textbf{three prepositional phrases} (for example: \eng{in the morning}, \eng{with my parents}, \eng{at home}).
\item Write in the \textbf{simple present} or \textbf{simple past}, but be consistent.
\end{itemize}

The best texts will be read aloud during the English Day and published in the magazine. Send your paragraph to the club secretary before Friday.

The English Club Committee
\end{textebox}

Votre classe de Terminale A a décidé de participer à ce concours organisé par le club d'anglais d'un lycée de Buea. Chaque élève rédigera son texte, le fera relire par un binôme, puis s'entraînera à l'oral grâce à une interview croisée. Les meilleures productions seront enregistrées et envoyées au club.

\courssub{Consignes}

\begin{enumerate}
\item \textbf{Préparation du lexique (5 min).} En silence, listez sur votre brouillon au moins huit phrasal verbs de la leçon liés à la vie familiale et sociale, avec leur traduction. Entourez ceux qui sont \cle{séparables} (\eng{put on}, \eng{tidy up}, \eng{pick up}...) et rappelez-vous la règle de placement du pronom complément.
\item \textbf{Rédaction individuelle (15 min).} Rédigez un paragraphe d'environ huit lignes intitulé \eng{A Day in My Family}. Contraintes : au moins six phrasal verbs différents, au moins trois prepositional phrases, un seul temps principal (présent ou passé simple). Soulignez vos phrasal verbs au crayon pour vous auto-contrôler.
\item \textbf{Relecture croisée (10 min).} Échangez votre copie avec votre voisin. Le relecteur doit : (a) souligner tous les phrasal verbs et vérifier qu'ils sont bien formés (verbe + particule) ; (b) vérifier le placement des compléments, surtout les pronoms (\eng{put them on}, jamais \eng{put on them}) ; (c) entourer les prepositional phrases et vérifier la préposition (\eng{in the morning}, \eng{at school}, \eng{with my parents}) ; (d) signaler toute répétition ou tout verbe français traduit mot à mot.
\item \textbf{Interview croisée (10 min).} À l'oral, posez à votre binôme trois questions construites avec des phrasal verbs, par exemple : \eng{Do you get on with your brothers and sisters?} \eng{Who looks after the baby when your mother is at the market?} \eng{What time do you usually come back from school?} Répondez en phrases complètes, en réutilisant un phrasal verb dans chaque réponse.
\item \textbf{Production finale (10 min).} Corrigez votre paragraphe en tenant compte des remarques du relecteur, puis recopiez-le proprement.
\item \textbf{Restitution collective (10 min).} Deux ou trois élèves lisent leur texte à la classe. Les auditeurs lèvent la main à chaque phrasal verb entendu et le notent ; la classe vérifie ensuite que les contraintes du concours sont respectées.
\end{enumerate}

\courssub{Ressources à mobiliser}

\begin{aretenir}[Boîte à outils de la leçon]
\textbf{Phrasal verbs utiles (vie familiale et sociale) :}
\begin{itemize}
\item \eng{get up} (se lever), \eng{wake up} (se réveiller), \eng{put on} (mettre, enfiler — séparable), \eng{take off} (enlever — séparable) ;
\item \eng{look after} (s'occuper de — inséparable), \eng{help out} (donner un coup de main), \eng{tidy up} (ranger — séparable), \eng{pick up} (ramasser, aller chercher — séparable) ;
\item \eng{go out} (sortir), \eng{come back} (revenir), \eng{get on with} (bien s'entendre avec — inséparable), \eng{grow up} (grandir), \eng{set off} (partir), \eng{sit down} (s'asseoir), \eng{turn off} (éteindre), \eng{run out of} (manquer de).
\end{itemize}
\textbf{Prepositional phrases utiles :} \eng{in the morning / afternoon / evening}, \eng{at six o'clock}, \eng{at school}, \eng{at home}, \eng{with my parents}, \eng{by bendskin}, \eng{on foot}, \eng{after lunch}, \eng{in front of the television}, \eng{at night}.
\end{aretenir}

\begin{attention}[Rappel des pièges — Check-list avant de rendre]
\begin{itemize}
\item Phrasal verb séparable + pronom : le pronom se place \textbf{obligatoirement} entre le verbe et la particule : \eng{I put them on} (« je les mets »), jamais \eng{I put on them}.
\item Phrasal verb inséparable : le complément suit toujours la particule/préposition : \eng{She looks after her little sister} (« elle s'occupe de sa petite sœur »).
\item Ne traduisez pas le français mot à mot : « s'occuper de » se dit \eng{look after}, pas \eng{occupy} ; « supporter quelqu'un » se dit \eng{put up with}, pas \eng{support somebody} ; « ressembler à » se dit \eng{take after}, pas \eng{resemble to}.
\item Prépositions figées : \eng{in the morning} mais \eng{at night} ; \eng{on foot} mais \eng{by taxi/bendskin} ; \eng{at home/school} mais \eng{in Yaoundé/Cameroon}.
\item \eng{look forward to} + V-\textbf{ing} (gérondif), pas infinitif.
\end{itemize}
\end{attention}

\courssub{Proposition de production et corrigé commenté}

\begin{exoresolu}[A Day in My Family — modèle de production]
Rédigez un paragraphe d'environ huit lignes respectant les contraintes du concours (au moins six phrasal verbs, au moins trois prepositional phrases), puis répondez oralement à la question : \eng{Do you get on with your brothers and sisters?}
\tcblower
\textbf{Modèle de rédaction :}

\eng{A Day in My Family}

\eng{In the morning, I usually wake up at half past five and get up immediately, because our house in Mile 4 is far from school. I put on my uniform quickly and help my mother tidy up the kitchen before breakfast. My father sets off for work at seven o'clock, and I look after my little brother while my mother prepares his food. I go to school on foot with my best friend, and we come back together in the afternoon. In the evening, the whole family sits down in front of the television, and I get on very well with my sisters because we never argue about the programmes. At night, I take off my shoes, say goodnight to my parents and go to bed early.}

« Le matin, je me réveille d'habitude à cinq heures et demie et je me lève immédiatement, car notre maison à Mile 4 est loin de l'école. J'enfile vite mon uniforme et j'aide ma mère à ranger la cuisine avant le petit-déjeuner. Mon père part au travail à sept heures, et je m'occupe de mon petit frère pendant que ma mère prépare sa nourriture. Je vais à l'école à pied avec mon meilleur ami, et nous rentrons ensemble l'après-midi. Le soir, toute la famille s'assoit devant la télévision, et je m'entends très bien avec mes sœurs car nous ne nous disputons jamais pour les programmes. La nuit, j'enlève mes chaussures, je dis bonne nuit à mes parents et je me couche tôt. »

\textbf{Commentaire des choix de langue :}
\begin{itemize}
\item \textbf{Dix phrasal verbs} sont employés : \eng{wake up}, \eng{get up}, \eng{put on}, \eng{tidy up}, \eng{set off}, \eng{look after}, \eng{come back}, \eng{sit down}, \eng{get on with}, \eng{take off} — soit plus que le minimum exigé, ce qui sécurise la copie.
\item \textbf{Placement des compléments} : dans \eng{I put on my uniform}, le nom peut suivre la particule ; avec un pronom, on écrirait \eng{I put it on}. Dans \eng{I look after my little brother}, le verbe est inséparable : le complément suit toujours \eng{after}.
\item \textbf{Six prepositional phrases} : \eng{in the morning}, \eng{at half past five}, \eng{at seven o'clock}, \eng{on foot}, \eng{in front of the television}, \eng{at night} — prépositions toutes correctes et idiomatiques.
\item \textbf{Cohérence temporelle} : le présent simple est maintenu du début à la fin, avec les adverbes de fréquence bien placés (\eng{usually wake up}, \eng{never argue}).
\item \textbf{Ancrage culturel} : la mention de \eng{Mile 4} (quartier de Bamenda) et la vie familiale camerounaise rendent le texte authentique sans sacrifier l'anglais naturel.
\end{itemize}

\textbf{Réponse orale modèle :} \eng{Yes, I get on very well with my brothers and sisters. When my parents go out, I look after them, and we help each other out with our homework.} (« Oui, je m'entends très bien avec mes frères et sœurs. Quand mes parents sortent, je m'occupe d'eux, et nous nous aidons mutuellement pour les devoirs. ») — Notez la réutilisation de trois phrasal verbs dans une réponse complète et fluide.
\end{exoresolu}

\courssub{Bilan de l'activité}

Cette activité d'intégration a permis de mobiliser simultanément les quatre savoir-faire de la leçon : la \textbf{compréhension écrite} (lecture de la consigne du concours et du texte du binôme), le \textbf{lexique} (choix et traduction des phrasal verbs), la \textbf{grammaire} (placement des particules et des compléments, choix des prépositions, cohérence des temps) et la \textbf{production orale et écrite} (récit personnel et interview croisée).

\begin{methode}[Grille d'auto-évaluation]
Avant de rendre votre texte, vérifiez chaque point :
\begin{enumerate}
\item Ai-je utilisé au moins six phrasal verbs différents et correctement formés ?
\item Les pronoms compléments sont-ils placés entre le verbe et la particule (\eng{put them on}) ?
\item Ai-je au moins trois prepositional phrases avec la bonne préposition (\eng{in the morning}, \eng{at night}, \eng{on foot}) ?
\item Mon temps principal est-il constant du début à la fin ?
\item Ai-je évité les traductions mot à mot du français (\eng{look after} et non \eng{occupy}, \eng{put up with} et non \eng{support}) ?
\item Suis-je capable de poser et de comprendre trois questions orales avec des phrasal verbs ?
\end{enumerate}
\end{methode}

Pour aller plus loin, les élèves peuvent enregistrer leur récit au téléphone, réécouter leur prononciation des particules (qui portent souvent l'accent : \eng{put them ON}, \eng{come BACK}) et transformer leur texte au passé simple pour raconter « \eng{Last Sunday in My Family} » — un excellent entraînement à la conjugaison des verbes irréguliers (\eng{woke up}, \eng{set off}, \eng{came back}, \eng{sat down}, \eng{took off}).

\newpage

\coursec{Méthodes et exercices résolus}

\coursec{Lesson 10 — Grammar: Phrasal verbs and prepositional phrases}

\courssub{Fiche méthode 1 : Identifier un phrasal verb dans un texte}

\begin{textebox}[Texte support : vie familiale]
“When my parents separated, my elder sister brought me up almost alone. She gave up her studies to look after me, and we got on very well. Last year, she moved out to Douala, but we still keep in touch every week.”
\end{textebox}

\begin{methode}[Comment repérer un phrasal verb]
\begin{enumerate}
\item Repère dans la phrase un \cle{verbe} suivi d'une \cle{particule} (adverbe ou préposition) : \eng{up, down, in, out, on, off, after, over, with...}
\item Teste le sens : remplace mentalement le groupe par un verbe simple de sens équivalent. Si le sens du groupe est différent du sens du verbe seul, c'est un \cle{phrasal verb}. \eng{She looks after her grandmother.} « Elle s'occupe de sa grand-mère. » (le sens n'est pas « regarder »).
\item Distingue du verbe + préposition ordinaire : dans \eng{She looked at the photo} « Elle a regardé la photo », \eng{look} garde son sens littéral ; ce n'est pas un phrasal verb.
\item Note le phrasal verb dans ton cahier avec sa traduction et un exemple complet, jamais seul.
\item Vérifie dans le contexte : beaucoup de phrasal verbs ont plusieurs sens (\eng{take off} = « décoller » ou « enlever »).
\end{enumerate}
\end{methode}

\begin{exoresolu}[Repérage dans un texte]
\textbf{Énoncé.} Read the text above and underline all the phrasal verbs, then give their meaning in French.

\tcblower
\textbf{Solution détaillée.}
\begin{enumerate}
\item \eng{brought me up} : verbe \eng{bring} + particule \eng{up} ; sens = « élever (un enfant) » $\neq$ « apporter ». \textbf{Phrasal verb} (séparable) — « m'a élevé ».
\item \eng{gave up} : \eng{give} + \eng{up} ; sens = « abandonner, renoncer à ». \textbf{Phrasal verb} (séparable) — « elle a abandonné ».
\item \eng{look after} : \eng{look} + \eng{after} ; sens = « s'occuper de, prendre soin de » $\neq$ « regarder ». \textbf{Phrasal verb} (verbe + préposition, inséparable, type 3).
\item \eng{got on} : \eng{get} + \eng{on} ; ici \eng{get on (well)} (intransitif) = « bien s'entendre, bien se passer ». \textbf{Phrasal verb} (intransitif, type 1).
\item \eng{moved out} : \eng{move} + \eng{out} ; sens = « déménager, quitter le domicile ». \textbf{Phrasal verb} (intransitif, type 1).
\item \eng{keep in touch} : locution figée (verbe + nom + préposition) = « rester en contact ». \textbf{Expression à retenir} comme un tout.
\end{enumerate}
\textbf{Conclusion.} Cinq phrasal verbs et une expression figée ; aucun ne garde le sens littéral du verbe de base, ce qui confirme l'identification.
\end{exoresolu}

\courssub{Fiche méthode 2 : Construire correctement un phrasal verb (place du complément)}

\begin{methode}[Séparable ou inséparable ? Les quatre types]
\begin{enumerate}
\item \textbf{Type 1 — Intransitif} (pas de complément) : \eng{The plane took off.} « L'avion a décollé. » Rien à placer.
\item \textbf{Type 2 — Transitif séparable} : le complément nominal peut se mettre avant ou après la particule : \eng{She brought up the children} / \eng{She brought the children up.} Mais un \cle{pronom} se place \textbf{obligatoirement} entre le verbe et la particule : \eng{She brought them up.} (jamais \eng{She brought up them.}).
\item \textbf{Type 3 — Transitif inséparable} (verbe + préposition) : le complément (nom ou pronom) suit toujours la préposition : \eng{He looks after his sister} / \eng{He looks after her.} (jamais \eng{He looks her after.}).
\item \textbf{Type 4 — Verbe + adverbe + préposition} (trois mots) : inséparable : \eng{I get on with my neighbours} / \eng{I get on with them.}
\item Réflexe d'examen : si le complément est un \cle{pronom} (\eng{it, them, her, him}), demande-toi si le verbe est séparable ; dans le doute, choisis un verbe inséparable ou place le pronom au milieu pour les séparables.
\end{enumerate}
\end{methode}

\begin{exoresolu}[Transformation : remplacer le nom par un pronom]
\textbf{Énoncé.} Replace the underlined noun with a pronoun and rewrite each sentence.

1. Please \textbf{turn off} \underline{the lights} before leaving.
2. My aunt \textbf{looks after} \underline{the children} on Saturdays.
3. They \textbf{put off} \underline{the meeting} until Monday.

\tcblower
\textbf{Solution détaillée.}
\begin{enumerate}
\item \eng{turn off} est \textbf{séparable} (type 2) ; le pronom \eng{them} (\eng{the lights}, pluriel) se place \textbf{entre} verbe et particule : \eng{Please turn them off before leaving.} « Éteins-les avant de partir. »
\item \eng{look after} est \textbf{inséparable} (type 3) ; le pronom suit la préposition : \eng{My aunt looks after them on Saturdays.} (jamais \eng{looks them after}).
\item \eng{put off} (« remettre, reporter ») est \textbf{séparable} (type 2) : \eng{They put it off until Monday.} (\eng{the meeting} est singulier $\to$ \eng{it}).
\end{enumerate}
\textbf{Conclusion.} La place du pronom dépend du type du phrasal verb : au milieu pour les séparables, après la préposition pour les inséparables.
\end{exoresolu}

\courssub{Fiche méthode 3 : Employer les phrasal verbs de la vie familiale et sociale}

\begin{methode}[Choisir le bon phrasal verb : champs sémantiques]
\begin{enumerate}
\item Classe mentalement les phrasal verbs par \cle{champ sémantique} :
\begin{itemize}
\item \textbf{Enfance et éducation} : \eng{bring up} (élever), \eng{grow up} (grandir), \eng{look after} (s'occuper de), \eng{take after} (ressembler à un parent — physique ou caractère).
\item \textbf{Relations} : \eng{get on (well) with} (bien s'entendre avec), \eng{fall out with} (se brouiller avec), \eng{make up} (se réconcilier), \eng{break up / split up} (se séparer, rompre), \eng{put up with} (supporter, tolérer).
\item \textbf{Vie quotidienne} : \eng{move out} (quitter le domicile parental), \eng{settle down} (se poser, se marier, fonder un foyer), \eng{keep in touch} (rester en contact), \eng{come round} (passer rendre visite).
\end{itemize}
\item Identifie dans la phrase française l'idée exprimée (« s'entendre », « élever », « se réconcilier ») et choisis le phrasal verb correspondant.
\item Conjugue le verbe de base au temps demandé (attention aux irréguliers : \eng{bring – brought – brought}, \eng{grow – grew – grown}, \eng{fall – fell – fallen}, \eng{break – broke – broken}, \eng{make – made – made}).
\item Vérifie la place du complément (fiche 2).
\end{enumerate}
\end{methode}

\begin{exoresolu}[Complétion à trous]
\textbf{Énoncé.} Complete the sentences with the correct phrasal verb in the correct tense: \emph{bring up, get on with, fall out with, grow up, take after, settle down}.

1. I \_\_\_\_\_ in Bamenda, but I now live in Yaoundé.
2. My brother \_\_\_\_\_ his best friend last week; they no longer talk to each other.
3. She \_\_\_\_\_ her mother: they are both very patient.
4. After years of travelling, he finally decided to \_\_\_\_\_ and start a family.
5. Do you \_\_\_\_\_ your new neighbours?
6. Their grandparents \_\_\_\_\_ them after their parents died.

\tcblower
\textbf{Solution détaillée.}
\begin{enumerate}
\item \eng{grew up} — « grandir » ; le contexte (\eng{now}) impose le prétérit ; irrégulier \eng{grow – grew – grown}. « J'ai grandi à Bamenda. »
\item \eng{fell out with} — « se brouiller avec » ; \eng{last week} $\to$ prétérit ; irrégulier \eng{fall – fell – fallen}.
\item \eng{takes after} — « ressembler à (un parent) » ; vérité générale $\to$ présent simple, 3\textsuperscript{e} personne : \eng{takes}.
\item \eng{settle down} — « se poser, fonder un foyer » ; après \eng{decided to}, infinitif sans modification.
\item \eng{get on with} — « t'entends-tu avec » ; question au présent avec l'auxiliaire \eng{do}.
\item \eng{brought up} — « ont élevé » ; prétérit ; irrégulier \eng{bring – brought – brought}.
\end{enumerate}
\textbf{Conclusion.} Choisir le sens, puis conjuguer correctement le verbe de base : les deux opérations sont indispensables pour obtenir le point.
\end{exoresolu}

\courssub{Fiche méthode 4 : Construire et employer les prepositional phrases}

\begin{methode}[Les locutions prépositionnelles (prepositional phrases)]
\begin{enumerate}
\item Une \cle{prepositional phrase} = préposition + groupe nominal (déterminant + nom) : \eng{in the morning, on foot, at school, by bus, for a long time}. Elle fonctionne comme un complément de temps, de lieu, de manière, de cause, etc.
\item Apprends les \cle{prépositions figées} par cœur, en contexte :
\begin{itemize}
\item Temps : \eng{at 7 o'clock} (heure précise), \eng{on Monday / on 15th May} (jour/date), \eng{in January / in 2024 / in the morning / in the afternoon / in the evening} (mois, année, saison, partie de journée) \textbf{mais} \eng{at night}.
\item Lieu : \eng{at home / at school / at work}, \eng{in the village square}, \eng{on the bus / on the train / on foot} (mais \eng{in my car}).
\item Moyens/Manière : \eng{by bus / by car / by train / by plane} (sans article), \eng{on foot}.
\item Locutions complexes : \eng{in charge of} (responsable de), \eng{on behalf of} (au nom de), \eng{in spite of} (malgré), \eng{according to} (selon), \eng{depend on} (dépendre de).
\end{itemize}
\item Place : en général en fin de phrase ; en tête de phrase, elle est suivie d'une virgule : \eng{In my family, everyone helps with the cooking.}
\item Ne traduis jamais mot à mot depuis le français : « à pied » = \eng{on foot} (pas \eng{by foot}) ; « dans le bus » = \eng{on the bus} (pas \eng{in the bus}) ; « assister à » = \eng{attend} (pas \eng{assist}).
\end{enumerate}
\end{methode}

\begin{exoresolu}[Choix de la bonne préposition]
\textbf{Énoncé.} Fill in the gaps with the correct preposition: \emph{at, on, in, by, for}.

1. The family meeting starts \_\_\_\_\_ 4 p.m. \_\_\_\_\_ Saturday.
2. My father goes to work \_\_\_\_\_ car, but my mother goes \_\_\_\_\_ foot.
3. We have lived in Douala \_\_\_\_\_ ten years.
4. \_\_\_\_\_ spite of the rain, the wedding took place \_\_\_\_\_ the village square.
5. She spoke \_\_\_\_\_ behalf of the whole family.

\tcblower
\textbf{Solution détaillée.}
\begin{enumerate}
\item \eng{at 4 p.m.} (heure précise $\to$ \cle{at}) ; \eng{on Saturday} (jour $\to$ \cle{on}).
\item \eng{by car} (moyen de transport sans article $\to$ \cle{by}) ; \eng{on foot} (locution figée, exception : « à pied »).
\item \eng{for ten years} (durée $\to$ \cle{for} ; \eng{since} s'emploie avec un point de départ : \eng{since 2014}).
\item \eng{In spite of} (locution figée = « malgré » ; ne pas oublier le \eng{of}) ; \eng{in the village square} (lieu $\to$ \cle{in}).
\item \eng{on behalf of} (locution figée = « au nom de »).
\end{enumerate}
\textbf{Conclusion.} Les locutions prépositionnelles sont figées : elles s'apprennent comme des blocs et ne se déduisent pas du français.
\end{exoresolu}

\courssub{Fiche méthode 5 : Traduire sans calquer le français}

\begin{methode}[Éviter les calques dans les traductions]
\begin{enumerate}
\item Repère dans la phrase française le verbe ou la locution qui cache un phrasal verb anglais : « s'occuper de » $\to$ \eng{look after} ; « élever » $\to$ \eng{bring up} ; « supporter » $\to$ \eng{put up with} ; « se réconcilier » $\to$ \eng{make up}.
\item Méfie-toi des \cle{faux amis} : « assister à » $\neq$ \eng{assist} (c'est \eng{attend}) ; « supporter » $\neq$ \eng{support} dans le sens « tolérer » ($\to$ \eng{put up with}) ; « occuper de » $\neq$ \eng{occupy of} (c'est \eng{look after}).
\item Traduis d'abord le \cle{sens global}, puis choisis la structure anglaise idiomatique ; ne construis jamais un phrasal verb imaginaire par calque.
\item Vérifie le temps : le passé composé français se rend souvent par le prétérit simple (\eng{yesterday, last week}) ou le present perfect (\eng{since, for, already, never}).
\item Relis en contrôlant : place du pronom complément (fiche 2), préposition finale éventuelle, 3\textsuperscript{e} personne du singulier (\eng{-s} au présent).
\end{enumerate}
\end{methode}

\begin{exoresolu}[Traduction de phrases]
\textbf{Énoncé.} Translate into English.

1. Ma grand-mère nous a élevés dans un petit village près de Buea.
2. Je ne supporte plus le comportement de mon voisin.
3. Après leur dispute, ils se sont réconciliés le lendemain.
4. Elle s'occupe de ses petits frères depuis deux ans.

\tcblower
\textbf{Solution détaillée.}
\begin{enumerate}
\item « a élevés » = \eng{bring up} au prétérit ; pronom complément « nous » $\to$ \eng{us}, placé \textbf{entre} verbe et particule (séparable, type 2) : \eng{My grandmother brought us up in a small village near Buea.}
\item « ne supporte plus » = \eng{put up with} (inséparable, type 4) au présent + négation : \eng{I can no longer put up with my neighbour's behaviour.} (ou \eng{I can't put up with... any longer.}) — ne pas utiliser \eng{support} seul.
\item « se sont réconciliés » = \eng{make up} au prétérit (\eng{made up}) ; « le lendemain » = \eng{the next day} : \eng{After their quarrel, they made up the next day.}
\item « s'occupe de » = \eng{look after} ; « depuis deux ans » + action qui continue $\to$ present perfect + \eng{for} : \eng{She has looked after her little brothers for two years.} (ou present perfect continu : \eng{She has been looking after her little brothers for two years.}).
\end{enumerate}
\textbf{Conclusion.} Une bonne traduction repose sur trois contrôles : choix du phrasal verb idiomatique, temps correct, place du complément.
\end{exoresolu}

\begin{attention}[Les 5 erreurs qui coûtent des points au Bac]
\begin{enumerate}
\item \textbf{Pronom mal placé} : \eng{She brought up them} est faux ; avec un pronom, le phrasal verb séparable (type 2) exige la position médiane : \eng{She brought them up.}
\item \textbf{Calque du français} : \eng{to assist a meeting} pour « assister à une réunion » (correct : \eng{to attend a meeting}) ; \eng{to occupy of the children} (correct : \eng{to look after the children}).
\item \textbf{Prépositions figées confondues} : \eng{in Saturday}, \eng{by foot}, \eng{depend of} sont faux ; écris \eng{on Saturday}, \eng{on foot}, \eng{depend on}.
\item \textbf{Verbes irréguliers mal conjugués} : \eng{bringed up}, \eng{falled out}, \eng{growed up} n'existent pas ; apprends \eng{bring – brought – brought}, \eng{fall – fell – fallen}, \eng{grow – grew – grown}.
\item \textbf{Temps verbal avec « depuis »} : \eng{She looks after them since two years} est un calque ; l'anglais exige le present perfect + \eng{for} : \eng{She has looked after them for two years.}
\end{enumerate}
\end{attention}

\newpage

\coursec{Exercices}

\courssub{Je vérifie mes connaissances}

\begin{exercice}[QCM : Structure des phrasal verbs]
Choisissez la bonne réponse parmi les propositions A, B, C ou D.

\begin{enumerate}
\item In a phrasal verb like \eng{``give up''}, the word \eng{``up''} is grammatically classified as a \dots
      \begin{itemize}
      \item[A.] preposition
      \item[B.] adverb (particle)
      \item[C.] conjunction
      \item[D.] adjective
      \end{itemize}

\item Which of the following sentences shows a \textbf{separable} phrasal verb used correctly with a pronoun object?
      \begin{itemize}
      \item[A.] She turned off it.
      \item[B.] She turned it off.
      \item[C.] She turned off the light.
      \item[D.] She turned the light off.
      \end{itemize}

\item Identify the \textbf{prepositional verb} in the following list:
      \begin{itemize}
      \item[A.] \eng{run into} (meet by chance)
      \item[B.] \eng{look after} (take care of)
      \item[C.] \eng{depend on} (rely on)
      \item[D.] \eng{put off} (postpone)
      \end{itemize}

\item Complete the sentence with the correct particle: \eng{``The meeting has been put \_\_\_\_ until next Monday.''}
      \begin{itemize}
      \item[A.] away
      \item[B.] off
      \item[C.] out
      \item[D.] up
      \end{itemize}
\end{enumerate}
\end{exercice}

\begin{exercice}[Vrai ou Faux justifié]
Indiquez si les affirmations suivantes sont \textbf{Vraies} ou \textbf{Fausses}. Justifiez votre réponse par une brève explication grammaticale ou un contre-exemple.

\begin{enumerate}
\item \eng{``Look up the word in the dictionary''} and \eng{``Look the word up in the dictionary''} are both correct.
\item With the phrasal verb \eng{``look after''} (take care of), you can say: \eng{``She looks after them.''}
\item \eng{``He ran up a huge bill''} and \eng{``He ran up the hill''} have the same grammatical structure.
\item A prepositional phrase like \eng{``in front of''} functions as a single preposition and is always followed by a noun phrase or a gerund.
\end{enumerate}
\end{exercice}

\begin{exercice}[Vocabulaire : Phrasal verbs de la vie familiale et sociale]
Associez chaque phrasal verb (1--6) à sa définition (A--F). Écrivez les paires correspondantes (ex. 1--C).

\begin{center}
\begin{tabular}{|c|l|c|l|}
\hline
\textbf{Phrasal Verb} & & \textbf{Définition} & \\
\hline
1. \eng{get along (with)} & & A. To stop a bad habit & \\
2. \eng{grow up} & & B. To have a good relationship with & \\
3. \eng{settle down} & & C. To become an adult & \\
4. \eng{fall out (with)} & & D. To start living a stable, routine life (marriage, job, house) & \\
5. \eng{bring up} & & E. To raise a child & \\
6. \eng{give up} & & F. To argue and stop being friends & \\
\hline
\end{tabular}
\end{center}
\end{exercice}

\begin{exercice}[Complétion : Particules et prépositions]
Complétez les phrases suivantes avec la particule ou la préposition appropriée (\eng{up, on, off, out, after, into, with, for, at, in}).

\begin{enumerate}
\item My elder brother really looks \_\_\_\_ my grandmother since she broke her hip.
\item We ran \_\_\_\_ an old friend from primary school at the market in Buea yesterday.
\item You cannot depend \_\_\_\_ him; he always forgets his promises.
\item The plane took \_\_\_\_ late because of the heavy rain in Douala.
\item She takes \_\_\_\_ her mother; they have the same smile and the same temper.
\item Please fill \_\_\_\_ this application form in capital letters.
\item The teacher asked the students to hand \_\_\_\_ their assignments before Friday.
\item He apologized \_\_\_\_ being late to the family reunion.
\end{enumerate}
\end{exercice}

\courssub{Je m'entraîne}

\begin{exercice}[Traduction : Français $\rightarrow$ Anglais]
Traduisez les phrases suivantes en anglais en utilisant obligatoirement un \textbf{phrasal verb} ou une \textbf{prepositional phrase} (surlignés dans la correction). Le vocabulaire entre parenthèses peut vous aider.

\begin{enumerate}
\item Je m'occupe de ma petite sœur tous les soirs pendant que ma mère travaille. (\eng{look after / take care of})
\item Ils se sont réconciliés après leur dispute au sujet de l'héritage. (\eng{make up})
\item Nous attendons les grandes vacances avec impatience. (\eng{look forward to})
\item Le match a été annulé à cause de la pluie torrentielle. (\eng{call off})
\item Elle a refusé la proposition de mariage, ce qui a surpris tout le village. (\eng{turn down})
\item Il faut que tu arrêtes de fumer si tu veux jouer au football correctement. (\eng{give up})
\item Ma tante a élevé trois neveux après le décès de leur mère. (\eng{bring up})
\item Le prix du cacao a augmenté de 20\% cette année. (\eng{go up})
\end{enumerate}
\end{exercice}

\begin{exercice}[Transformation de phrases : Remplacement par un phrasal verb]
Remplacez le verbe souligné par un \textbf{phrasal verb} approprié mis à la forme correcte. Conservez le sens et le temps de la phrase.

\begin{enumerate}
\item The children \textbf{resemble} their father in character and appearance. (\eng{take after})
\item The government has decided to \textbf{postpone} the elections until next year. (\eng{put off})
\item She \textbf{discovered} the truth about her adoption only last week. (\eng{find out})
\item The firemen \textbf{extinguished} the fire in the market after three hours. (\eng{put out})
\item We must \textbf{reduce} our expenses if we want to save money for the house. (\eng{cut down on})
\item The meeting \textbf{continued} until midnight. (\eng{go on})
\item He \textbf{invented} a ridiculous excuse for not doing his homework. (\eng{make up})
\item Please \textbf{submit} your application before the deadline. (\eng{hand in})
\end{enumerate}
\end{exercice}

\begin{exercice}[Réécriture avec pronoms : Position de la particule]
Réécrivez les phrases suivantes en remplaçant l'objet (nom) par le pronom personnel correspondant (\eng{it} ou \eng{them}). Attention à l'ordre des mots (séparation obligatoire).

\begin{enumerate}
\item Turn off \textbf{the lights} before leaving the classroom.
\item Pick up \textbf{your books} from the floor.
\item Look after \textbf{the children} while I go to the market.
\item Fill in \textbf{the form} carefully.
\item Give back \textbf{my dictionary} immediately.
\item Throw away \textbf{these old papers}.
\item Call up \textbf{your father} to wish him a happy birthday.
\item Put on \textbf{your raincoat}; it is pouring outside.
\end{enumerate}
\end{exercice}

\begin{exercice}[Choix de la préposition / Correction d'erreurs]
\textbf{Partie A :} Choisissez la préposition correcte entre parenthèses.
\textbf{Partie B :} Les phrases suivantes contiennent chacune une erreur (préposition manquante, superflue, mauvaise particule, ordre des mots). Soulignez l'erreur et réécrivez la phrase correcte.

\begin{enumerate}
\item \textbf{A.} She is married \textbf{(to / with)} a teacher from Bamenda.
\item \textbf{A.} We depend \textbf{(on / of)} our parents for financial support.
\item \textbf{A.} He apologized \textbf{(for / to)} his rude behaviour.
\item \textbf{A.} My father insisted \textbf{(on / to)} paying the bill.
\item \textbf{B.} I picked up \textbf{it} at the post office this morning.
\item \textbf{B.} She is interested \textbf{by} traditional music and dance.
\item \textbf{B.} Thank you for \textbf{to come} to my graduation ceremony.
\item \textbf{B.} The baby looks \textbf{like} his father, but he takes \textbf{after} his mother in character.
\end{enumerate}
\end{exercice}

\courssub{Je me prépare au Bac}

\begin{exercice}[Compréhension et langue : Texte original -- ``The Family Meeting'']
\begin{textebox}[Source : OnBuch+ Original]
Last Saturday, the atmosphere in the Njie household in Buea was tense. Mr. Njie had called a family meeting to \textbf{sort out} a recurring issue: the children's addiction to their smartphones. ``We need to \textbf{cut down on} screen time,'' he announced, ``or your grades will \textbf{go down} further.''

His eldest son, Ekema, tried to \textbf{talk his father out of} the new rule. ``Dad, we use phones for research too. You can't just \textbf{cut us off} from the world.'' His sister, Ngonde, \textbf{backed him up}, adding, ``Besides, you \textbf{check your messages} during meals all the time.''

Mr. Njie \textbf{put his foot down}. ``I \textbf{put up with} enough disrespect. The rule starts Monday. Phones in the basket by 8 p.m. No exceptions.'' He \textbf{looked around} the table. ``Does anyone want to \textbf{bring up} any other matter?''

Silence fell. Ekema \textbf{gave in} first. ``Okay, Dad. But can we \textbf{catch up on} messages at weekends?'' Mr. Njie smiled, relieved the conflict was over. ``We'll \textbf{work something out}.''
\end{textebox}

\textbf{Questions :}

\begin{enumerate}
\item \textbf{Vocabulary in context:} Explain the meaning of the following phrasal verbs \textbf{as used in the text} (do not give a dictionary definition, explain the specific meaning here):
      \begin{itemize}
      \item \eng{sort out} (l. 2)
      \item \textbf{cut down on} (l. 3)
      \item \textbf{talk ... out of} (l. 5)
      \item \textbf{back ... up} (l. 6)
      \item \textbf{put ... foot down} (l. 7)
      \item \textbf{put up with} (l. 7)
      \item \textbf{gave in} (l. 10)
      \item \textbf{catch up on} (l. 10)
      \item \textbf{work ... out} (l. 11)
      \end{itemize}

\item \textbf{Grammar analysis:}
      \begin{enumerate}
      \item Identify \textbf{three separable phrasal verbs} used with a noun object in the text. Show where the object is placed.
      \item Identify \textbf{two inseparable phrasal verbs} (or prepositional verbs) in the text.
      \item In the sentence ``\eng{You can't just cut us off from the world}'', why is the pronoun \eng{us} placed \textbf{between} \eng{cut} and \eng{off}?
      \item In the sentence ``\eng{I put up with enough disrespect}'', why can we not say ``\eng{I put up it with}''?
      \end{enumerate}

\item \textbf{Transformation:} Rewrite the following sentences using a phrasal verb from the text (or the list studied in the lesson) to replace the underlined words. Keep the same tense.
      \begin{enumerate}
      \item The children \textbf{resemble} their father in their stubbornness.
      \item The headmaster \textbf{tolerates} no noise in the library.
      \item They \textbf{resolved} their disagreement peacefully.
      \item She \textbf{supported} her brother's argument.
      \end{enumerate}
\end{enumerate}
\end{exercice}

\begin{exercice}[Traduction thématique : Français $\rightarrow$ Anglais]
Traduisez le paragraphe suivant en anglais. Vous devez impérativement utiliser au moins \textbf{six phrasal verbs ou prepositional phrases} différents parmi ceux étudiés (liste non exhaustive : \eng{get along with, grow up, settle down, fall out, bring up, give up, look after, depend on, look forward to, put up with, run into, take after, cut down on, find out, hand in, apply for}).

\begin{quote}
« Au Cameroun, la vie familiale reste un pilier essentiel. Les grands-parents élèvent souvent leurs petits-enfants pendant que les parents travaillent. Les enfants apprennent très tôt à respecter les aînés et à s'entendre avec leurs cousins. Quand un jeune veut se marier et fonder un foyer, il doit compter sur le soutien de toute la communauté. Parfois, des disputes éclatent à cause de l'argent, mais on finit toujours par se réconcilier car la famille est sacrée. Personnellement, j'ai hâte de retrouver ma famille au village pour les fêtes de fin d'année. »
\end{quote}
\end{exercice}

\begin{exercice}[Sujet de composition -- Type Bac]
\textbf{Choose ONE of the following topics and write an essay of about \textbf{250--300 words}.}

\begin{enumerate}
\item \textbf{Argumentative Essay:} ``Modern technology (smartphones, social media) is destroying family communication.'' Do you agree or disagree? Give reasons and examples from your experience or observation in Cameroon.
\item \textbf{Descriptive / Narrative Essay:} Describe a memorable family gathering (wedding, funeral, end-of-year party, \eng{``njangi''} meeting, etc.) you attended. Focus on the atmosphere, the interactions between generations, and a specific incident that \textbf{brought the family closer} or \textbf{caused a temporary fall-out}.
\item \textbf{Article for School Magazine:} Write an article entitled: ``\textbf{The Role of Elders in Bringing Up Responsible Youths in Cameroon Today.}'' Discuss the challenges (generation gap, western influence, economic pressure) and suggest solutions.
\end{enumerate}

\textbf{Instructions :}
\begin{itemize}
\item Plan your essay (Introduction, Body Paragraphs, Conclusion).
\item Use a variety of \textbf{phrasal verbs} and \textbf{prepositional phrases} correctly (underline them in your final copy if writing on paper).
\item Pay attention to tenses, subject-verb agreement, and spelling.
\item Respect the word count.
\end{itemize}
\end{exercice}

\newpage

\coursec{Corrigés des exercices}

\begin{corrige}[Exercice 1 — QCM : Structure des phrasal verbs]
\begin{enumerate}
\item \textbf{Bonne réponse : B — adverb (particle).} Dans un phrasal verb comme \eng{give up}, \eng{up} est une particule adverbiale, et non une préposition : elle modifie le sens du verbe et ne gouverne pas directement un complément. C'est d'ailleurs pourquoi on peut dire \eng{give it up} (le pronom se place entre le verbe et la particule), ce qui serait impossible avec une vraie préposition.
\item \textbf{Bonne réponse : B — \eng{She turned it off.}} Avec un phrasal verb \textbf{séparable}, un objet \emph{pronominal} doit obligatoirement se placer \textbf{entre} le verbe et la particule : \eng{turn it off}. La réponse A (\eng{turned off it}) est fautive. Les réponses C et D sont grammaticalement correctes, mais elles utilisent un \emph{nom} (\eng{the light}), pas un pronom : elles ne répondent donc pas à la question posée.
\item \textbf{Bonne réponse : C — \eng{depend on}.} \eng{Depend on} est un \emph{prepositional verb} : la préposition \eng{on} est inséparable du verbe et introduit obligatoirement le complément (\eng{depend on him}). \eng{Run into} et \eng{look after} sont des phrasal verbs inséparables ; \eng{put off} est un phrasal verb séparable.
\item \textbf{Bonne réponse : B — off.} \eng{The meeting has been put off until next Monday.} « La réunion a été reportée à lundi prochain. » \eng{Put off} signifie « reporter » ; \eng{put away} = ranger, \eng{put out} = éteindre, \eng{put up} = héberger/afficher.
\end{enumerate}
\end{corrige}

\begin{corrige}[Exercice 2 — Vrai ou Faux justifié]
\begin{enumerate}
\item \textbf{Vrai.} \eng{Look up} (chercher dans un dictionnaire) est un phrasal verb \textbf{séparable} : avec un objet nominal, la particule peut se placer après le verbe ou après l'objet. \eng{Look up the word} et \eng{look the word up} sont donc tous deux corrects. (Avec un pronom, seule \eng{look it up} serait possible.)
\item \textbf{Vrai.} \eng{Look after} est un phrasal verb \textbf{inséparable} (phrasal-prepositional verb) : l'objet, nom ou pronom, se place toujours \emph{après} la particule : \eng{She looks after them.} est correct. En revanche, \eng{She looks them after} serait fautif.
\item \textbf{Faux.} Dans \eng{He ran up a huge bill} (« il a accumulé une énorme facture »), \eng{run up} est un phrasal verb : \eng{up} est une particule qui change le sens du verbe. Dans \eng{He ran up the hill} (« il a couru en haut de la colline »), \eng{up} est une vraie préposition introduisant le complément \eng{the hill}. Test : on peut dire \eng{he ran it up} pour la facture, mais pas \eng{he ran up it} pour la colline de manière équivalente.
\item \textbf{Vrai.} Une \emph{prepositional phrase} comme \eng{in front of} fonctionne comme une préposition complexe unique ; elle est suivie d'un groupe nominal (\eng{in front of the house}) ou d'un gérondif (\eng{in front of everyone watching}), jamais d'un infinitif.
\end{enumerate}
\end{corrige}

\begin{corrige}[Exercice 3 — Vocabulaire : Phrasal verbs de la vie familiale et sociale]
\begin{center}
\begin{tabular}{|c|l|l|}
\hline
\rowcolor{popblueL} \textbf{Paire} & \textbf{Phrasal verb} & \textbf{Traduction} \\
\hline
1--B & \eng{get along (with)} & s'entendre (avec) \\
\hline
2--C & \eng{grow up} & grandir, devenir adulte \\
\hline
3--D & \eng{settle down} & se caser, mener une vie stable \\
\hline
4--F & \eng{fall out (with)} & se brouiller (avec) \\
\hline
5--E & \eng{bring up} & élever (un enfant) \\
\hline
6--A & \eng{give up} & abandonner (une mauvaise habitude) \\
\hline
\end{tabular}
\end{center}
\textbf{Exemples de réinvestissement :} \eng{My cousins get along very well with each other.} « Mes cousins s'entendent très bien entre eux. » ; \eng{She grew up in Bamenda.} « Elle a grandi à Bamenda. » ; \eng{They fell out over money but made up later.} « Ils se sont brouillés à cause de l'argent mais se sont réconciliés plus tard. »
\end{corrige}

\begin{corrige}[Exercice 4 — Complétion : Particules et prépositions]
\begin{enumerate}
\item \eng{My elder brother really looks \textbf{after} my grandmother since she broke her hip.} (\eng{look after} = s'occuper de — inséparable.)
\item \eng{We ran \textbf{into} an old friend from primary school at the market in Buea yesterday.} (\eng{run into} = rencontrer par hasard.)
\item \eng{You cannot depend \textbf{on} him; he always forgets his promises.} (\eng{depend on} = compter sur — prepositional verb.)
\item \eng{The plane took \textbf{off} late because of the heavy rain in Douala.} (\eng{take off} = décoller.)
\item \eng{She takes \textbf{after} her mother; they have the same smile and the same temper.} (\eng{take after} = ressembler à, tenir de.)
\item \eng{Please fill \textbf{in} this application form in capital letters.} (\eng{fill in} = remplir un formulaire.)
\item \eng{The teacher asked the students to hand \textbf{in} their assignments before Friday.} (\eng{hand in} = rendre, remettre un devoir.)
\item \eng{He apologized \textbf{for} being late to the family reunion.} (\eng{apologize for + V-ing} = s'excuser de faire quelque chose.)
\end{enumerate}
\end{corrige}

\begin{corrige}[Exercice 5 — Traduction : Français $\rightarrow$ Anglais]
\begin{enumerate}
\item \eng{I \textbf{look after} my little sister every evening while my mother is working.} (ou \eng{I take care of my little sister...})
\item \eng{They \textbf{made up} after their quarrel about the inheritance.} (Attention : \eng{make up} = se réconcilier ; temps : prétérit.)
\item \eng{We are \textbf{looking forward to} the long holidays.} (ou \eng{...to the summer holidays.} Rappel : \eng{to} est ici une préposition, donc suivi d'un nom ou d'un gérondif.)
\item \eng{The match was \textbf{called off} because of the torrential rain.} (Voix passive au prétérit : \eng{was called off}.)
\item \eng{She \textbf{turned down} the marriage proposal, which surprised the whole village.} (\eng{turn down} = refuser ; prétérit irrégulier : \eng{turned}.)
\item \eng{You must \textbf{give up} smoking if you want to play football properly.} (\eng{give up + V-ing} = arrêter de faire.)
\item \eng{My aunt \textbf{brought up} three nephews after their mother's death.} (Prétérit irrégulier : \eng{bring -- brought -- brought}.)
\item \eng{The price of cocoa has \textbf{gone up} by 20\% this year.} (Present perfect avec \eng{this year} ; participe passé irrégulier : \eng{go -- went -- gone}.)
\end{enumerate}
\end{corrige}

\begin{corrige}[Exercice 6 — Transformation de phrases : Remplacement par un phrasal verb]
\begin{enumerate}
\item \eng{The children \textbf{take after} their father in character and appearance.} (Présent simple, 3\textsuperscript{e} pers. pluriel : pas de \eng{-s}.)
\item \eng{The government has decided to \textbf{put off} the elections until next year.} (Après \eng{to}, infinitif de base.)
\item \eng{She \textbf{found out} the truth about her adoption only last week.} (Prétérit irrégulier : \eng{find -- found -- found}.)
\item \eng{The firemen \textbf{put out} the fire in the market after three hours.} (Prétérit irrégulier : \eng{put -- put -- put}.)
\item \eng{We must \textbf{cut down on} our expenses if we want to save money for the house.} (Après le modal \eng{must}, base verbale.)
\item \eng{The meeting \textbf{went on} until midnight.} (Prétérit irrégulier : \eng{go -- went -- gone}.)
\item \eng{He \textbf{made up} a ridiculous excuse for not doing his homework.} (Prétérit irrégulier : \eng{make -- made -- made} ; ici \eng{make up} = inventer.)
\item \eng{Please \textbf{hand in} your application before the deadline.} (Impératif : base verbale.)
\end{enumerate}
\end{corrige}

\begin{corrige}[Exercice 7 — Réécriture avec pronoms : Position de la particule]
\begin{enumerate}
\item \eng{Turn \textbf{them} off before leaving the classroom.}
\item \eng{Pick \textbf{them} up from the floor.}
\item \eng{Look after \textbf{them} while I go to the market.} (\eng{Look after} est \textbf{inséparable} : le pronom reste après la particule. Ne pas écrire \eng{look them after}.)
\item \eng{Fill \textbf{it} in carefully.}
\item \eng{Give \textbf{it} back immediately.}
\item \eng{Throw \textbf{them} away.}
\item \eng{Call \textbf{him} up to wish him a happy birthday.} (\eng{Your father} est une personne masculine : pronom \eng{him}, pas \eng{it}.)
\item \eng{Put \textbf{it} on; it is pouring outside.}
\end{enumerate}
\textbf{Règle rappelée :} avec un phrasal verb séparable, le pronom objet se place \emph{toujours} entre le verbe et la particule (\eng{turn them off}, jamais \eng{turn off them}).
\end{corrige}

\begin{corrige}[Exercice 8 — Choix de la préposition / Correction d'erreurs]
\textbf{Partie A :}
\begin{enumerate}
\item \eng{She is married \textbf{to} a teacher from Bamenda.} (\eng{married to}, jamais \eng{married with} — influence du français « marié avec ».)
\item \eng{We depend \textbf{on} our parents for financial support.}
\item \eng{He apologized \textbf{for} his rude behaviour.} (\eng{apologize for something} ; \eng{apologize to somebody}.)
\item \eng{My father insisted \textbf{on} paying the bill.} (\eng{insist on + V-ing}.)
\end{enumerate}
\textbf{Partie B :}
\begin{enumerate}
\setcounter{enumi}{4}
\item Erreur : \eng{picked up it}. Correction : \eng{I picked \textbf{it} up at the post office this morning.} (Pronom obligatoirement avant la particule.)
\item Erreur : \eng{interested by}. Correction : \eng{She is interested \textbf{in} traditional music and dance.} (\eng{interested in}, pas \eng{by} — influence du français « intéressé par ».)
\item Erreur : \eng{for to come}. Correction : \eng{Thank you for \textbf{coming} to my graduation ceremony.} (Après une préposition, toujours le gérondif, jamais l'infinitif.)
\item \textbf{Phrase correcte :} \eng{The baby \textbf{looks like} his father, but he \textbf{takes after} his mother in character.}
      \textbf{Explication :} \eng{look like} exprime une ressemblance \emph{physique} ; \eng{take after} exprime une ressemblance de \emph{caractère} (ou physique) avec un ascendant. La phrase originale utilise correctement les deux verbes. L'erreur fréquente est de les confondre (ex. : *\eng{He looks like his mother in character}).
\end{enumerate}
\end{corrige}

\begin{corrige}[Exercice 9 — Compréhension et langue : ``The Family Meeting'']
\textbf{1. Vocabulary in context :}
\begin{itemize}
\item \eng{sort out} : to resolve, to find a solution to the recurring problem of smartphone addiction.
\item \eng{cut down on} : to reduce the amount of time spent on screens.
\item \eng{talk his father out of} : to persuade his father not to apply the new rule.
\item \eng{backed him up} : supported her brother, agreed with him publicly.
\item \eng{put his foot down} : acted firmly, refused to negotiate any longer.
\item \eng{put up with} : tolerated, endured the children's disrespect.
\item \eng{gave in} : stopped resisting, accepted his father's decision.
\item \eng{catch up on} : read and answer the messages accumulated during the week.
\item \eng{work something out} : find a compromise, a practical solution together.
\end{itemize}
\textbf{2. Grammar analysis :}
\begin{enumerate}
\item Trois phrasal verbs séparables avec objet : \eng{talk \textbf{his father} out of} (objet nominal après le verbe) ; \eng{backed \textbf{him} up} (pronom) ; \eng{cut \textbf{us} off} (pronom). Avec un nom, les deux positions sont possibles : \eng{sort out \textbf{a recurring issue}} ou \eng{sort \textbf{a recurring issue} out}.
\item Deux phrasal-prepositional verbs inséparables : \eng{cut down on screen time} ; \eng{put up with enough disrespect}. On peut ajouter \eng{catch up on messages} et \eng{looked around the table}.
\item Dans \eng{cut us off}, \eng{cut off} est séparable : un objet \emph{pronominal} (\eng{us}) doit obligatoirement se placer entre le verbe et la particule. \eng{Cut off us} serait fautif.
\item \eng{Put up with} est un phrasal-prepositional verb \textbf{inséparable} : le complément suit toujours l'ensemble verbe + particules. On dit \eng{put up with it}, jamais \eng{put up it with}.
\end{enumerate}
\textbf{3. Transformation :}
\begin{enumerate}
\item \eng{The children \textbf{take after} their father in their stubbornness.}
\item \eng{The headmaster \textbf{puts up with} no noise in the library.}
\item \eng{They \textbf{sorted out} their disagreement peacefully.} (ou \eng{worked it out})
\item \eng{She \textbf{backed up} her brother's argument.} (ou \eng{backed her brother up})
\end{enumerate}
\textbf{Barème indicatif (sur 20) :} vocabulaire en contexte 6 pts ; analyse grammaticale 6 pts ; transformations 4 pts ; exactitude linguistique 4 pts.
\end{corrige}

\begin{corrige}[Exercice 10 — Traduction thématique : Français $\rightarrow$ Anglais]
\textbf{Proposition de traduction :}
\begin{quote}
\eng{In Cameroon, family life remains an essential pillar. Grandparents often \textbf{bring up} their grandchildren while the parents are working. Children learn very early to respect their elders and to \textbf{get along with} their cousins. When a young man wants to marry and \textbf{settle down}, he must \textbf{depend on} the support of the whole community. Sometimes, relatives \textbf{fall out} over money, but they always end up \textbf{making up} because the family is sacred. Personally, I am \textbf{looking forward to} joining my family in the village for the end-of-year celebrations.}
\end{quote}
\textbf{Points de vigilance :} \eng{elders} = les aînés (ne pas confondre avec \eng{elderly} = personnes âgées) ; \eng{make up} = se réconcilier ; \eng{look forward to + V-ing/nom} ; \eng{settle down} = fonder un foyer ; \eng{fall out over money} = se brouiller à cause de l'argent. Six structures cibles minimum sont bien présentes (\eng{bring up, get along with, settle down, depend on, fall out, look forward to}).
\end{corrige}

\begin{corrige}[Exercice 11 — Sujet de composition — Type Bac]
\textbf{Barème indicatif (sur 20) :} respect du sujet et de la consigne 4 pts ; organisation (introduction, développement, conclusion) 4 pts ; richesse et correction du vocabulaire, dont au moins 5 phrasal verbs/prepositional phrases bien employés 5 pts ; correction grammaticale (temps, accords, ordre des mots) 5 pts ; respect du nombre de mots et présentation 2 pts.

\textbf{Proposition de rédaction modèle (sujet 1) :}
\begin{quote}
\eng{In many Cameroonian homes today, smartphones have become uninvited guests at the dinner table. Some people claim that modern technology is destroying family communication. I partly agree, although I believe the problem lies in how we use these devices rather than in the devices themselves.}

\eng{On the one hand, it is true that screens often \textbf{cut families off} from genuine conversation. Teenagers \textbf{check} their phones during meals, parents \textbf{bring work up} at bedtime, and everyone ends up sitting together without speaking. In Yaoundé, I have seen children who can no longer \textbf{get along with} their grandparents because they never \textbf{look up} from their screens to listen to family stories. Little by little, mutual understanding \textbf{goes down} and conflicts \textbf{build up}.}

\eng{On the other hand, technology can also \textbf{bring relatives together}. Thanks to WhatsApp, my aunt in Douala \textbf{keeps in touch with} the family in Bamenda every single day. Video calls allow students abroad to \textbf{catch up on} family news. The real issue is discipline: families must \textbf{cut down on} screen time and \textbf{work out} clear rules, such as phone-free meals.}

\eng{To conclude, smartphones do not destroy family communication by themselves; careless habits do. If parents \textbf{put their foot down} and everyone learns to \textbf{give up} excessive screen time, technology can become a bridge instead of a wall.}
\end{quote}
\textbf{Points de vigilance :} éviter \eng{discuss about} (dire \eng{discuss something}) ; \eng{look forward to + V-ing} ; ne pas traduire « actuellement » par \eng{actually} (faux ami : dire \eng{nowadays, today}) ; accorder le verbe avec son sujet (\eng{technology allows}, pas \eng{allow}) ; soigner les paragraphes et les connecteurs (\eng{on the one hand, however, to conclude}).
\end{corrige}

\newpage

\coursec{Fiche bilan}

\begin{popfigure}
\begin{tikzpicture}[
  mindmap,
  grow cyclic,
  every node/.style={concept, execute at begin node=\hskip0pt},
  concept color=popblue,
  level 1/.append style={level distance=4.5cm, sibling angle=360/7},
  level 2/.append style={level distance=3cm, sibling angle=360/4},
  branch/.style={concept color=#1, edge from parent/.style={draw=#1, thick}},
]
\node[concept, text width=3.2cm, align=center, font=\bfseries\large, fill=popblue!90!popdark] {Phrasal Verbs \&\\ Prepositional Phrases}
  child[branch=poporange] { node[concept, text width=2.8cm, align=center] {Les 4 Types\\ (Intrans./Trans. Sep./Insep./3-part)} }
  child[branch=popgreen] { node[concept, text width=2.5cm, align=center] {Séparables vs\\ Inséparables\\ (Pronoms = séparation)} }
  child[branch=poppurple] { node[concept, text width=2.5cm, align=center] {Verbes fréquents\\ Famille \& Social\\ (get on, look after...)} }
  child[branch=poppink] { node[concept, text width=2.5cm, align=center] {Prepositional\\ Phrases\\ (phrases prépositives)} }
  child[branch=popgold] { node[concept, text width=2.5cm, align=center] {Règles clés\\ Ordre, Temps, Passif} }
  child[branch=popblue!60!poppurple] { node[concept, text width=2.5cm, align=center] {Pièges FR/EN\\ Faux amis, Ordre mots} }
  child[branch=popmuted] { node[concept, text width=2.5cm, align=center] {Checklist Bac\\ 10 points à cocher} };
\end{tikzpicture}
\legende{Carte mentale récapitulative : Phrasal Verbs et Prepositional Phrases (Terminale A)}
\end{popfigure}

\begin{bilan}
\end{bilan}

\courssub{Lexique clé — Phrasal Verbs : Famille \& Vie Sociale}
\begin{center}
\begin{tabularx}{\linewidth}{|X|X|X|}
\hline
\rowcolor{popblueL} \textbf{Phrasal Verb} & \textbf{Sens français} & \textbf{Exemple / Contexte} \\
\hline
\eng{get on (with)} & s'entendre (avec) & \eng{I get on well with my siblings.} (« Je m'entends bien avec mes frères et sœurs. ») \\
\hline
\eng{look after} & s'occuper de, garder & \eng{She looks after her nephew.} (« Elle s'occupe de son neveu. ») \\
\hline
\eng{bring up} & élever (enfant) & \eng{He was brought up in Buea.} (« Il a été élevé à Buea. ») \\
\hline
\eng{grow up} & grandir & \eng{They grew up in Douala.} (« Ils ont grandi à Douala. ») \\
\hline
\eng{take after} & ressembler à (famille) & \eng{She takes after her mother.} (« Elle ressemble à sa mère. ») \\
\hline
\eng{fall out (with)} & se fâcher, se brouiller & \eng{They fell out over money.} (« Ils se sont brouillés pour de l'argent. ») \\
\hline
\eng{make up} & se réconcilier & \eng{They made up yesterday.} (« Ils se sont réconciliés hier. ») \\
\hline
\eng{move out} & quitter le domicile & \eng{He moved out at 20.} (« Il a quitté la maison à 20 ans. ») \\
\hline
\eng{settle down} & s'installer, fonder foyer & \eng{They want to settle down soon.} (« Ils veulent s'installer bientôt. ») \\
\hline
\eng{count on} & compter sur & \eng{You can count on me.} (« Tu peux compter sur moi. ») \\
\hline
\eng{run into} & rencontrer par hasard & \eng{I ran into an old friend.} (« J'ai rencontré un vieux copain par hasard. ») \\
\hline
\eng{put up with} & supporter, tolérer & \eng{I can't put up with the noise.} (« Je ne supporte pas le bruit. ») \\
\hline
\end{tabularx}
\end{center}

\courssub{Règles de grammaire à connaître par cœur}
\begin{center}
\begin{tabularx}{\linewidth}{|>{\bfseries}p{3.2cm}|X|>{\itshape}X|}
\hline
\rowcolor{popblueL} \textbf{Structure / Règle} & \textbf{Explication} & \textbf{Exemple} \\
\hline
Type 1 : Intransitif & Pas d'objet. Verbe + Particule. & \eng{The plane took off.} \\
\hline
Type 2 : Transitif séparable & Objet nom \rightarrow place libre. Objet pronom \rightarrow \textbf{obligatoirement au milieu}. & \eng{Turn off the light. / Turn it off.} (\textbf{pas} *Turn off it) \\
\hline
Type 3 : Transitif inséparable & Objet toujours \textbf{après} la particule (même pronom). & \eng{Look after the baby. / Look after him.} \\
\hline
Type 4 : Three-part (Verb + Part + Prep) & Toujours inséparables. Objet après la préposition finale. & \eng{get on with, put up with, look forward to} \\
\hline
Passif possible ? & Seulement Types 2, 3, 4 (transitifs). Sujet = objet actif. & \eng{The baby is looked after by the nanny.} \\
\hline
Ordre avec adverbe & Adverbe souvent \textbf{avant} la particule (surtout Type 2). & \eng{He quickly put on his coat.} \\
\hline
Nominalisation & Verbe + Particule \rightarrow Nom (souvent trait d'union). & \eng{A break-up, a check-up, a grown-up.} \\
\hline
\end{tabularx}
\end{center}

\courssub{Méthodes express}
\begin{enumerate}
\item \textbf{Identifier le type :} Chercher l'objet (nom/pronom). Tester le pronom : \eng{Turn it off} (séparable) vs \eng{Look after it} (inséparable).
\item \textbf{Conjuguer :} Le verbe principal porte la marque du temps ; la particule/préposition ne bouge jamais. \eng{She \textbf{looked} after him.}
\item \textbf{Traduire :} Éviter le mot à mot. Chercher le verbe français unique (\eng{bring up} = élever ; \eng{put up with} = supporter).
\item \textbf{Transformer (Actif $\leftrightarrow$ Passif) :} Vérifier transitivité. \eng{They brought up the child.} $\rightarrow$ \eng{The child was brought up by them.}
\item \textbf{Rédiger :} Intégrer 3 à 4 phrasal verbs naturels dans un paragraphe sur la famille (ex. \eng{grow up, take after, fall out, make up}).
\end{enumerate}


\begin{aretenir}[Checklist avant le Bac]
$\square$ Je distingue les 4 types de phrasal verbs (intransitif, trans. séparable, trans. inséparable, three-part).\\
$\square$ Je place \textbf{correctement le pronom objet} avec les verbes séparables (\eng{turn it off}, jamais \eng{turn off it}).\\
$\square$ Je sais que les verbes à trois particules (\eng{look forward to, put up with}) sont \textbf{toujours inséparables}.\\
$\square$ Je connais par cœur la liste des 12 phrasal verbs « Famille/Vie sociale » (sens + type).\\
$\square$ Je forme le \textbf{passif} uniquement avec les verbes transitifs (Types 2, 3, 4).\\
$\square$ J'évite le piège « \eng{look for} $\neq$ regarder » (chercher) et « \eng{look after} $\neq$ regarder après » (s'occuper de).\\
$\square$ Je ne confonds pas \textbf{phrasal verb} (verbe + particule adverbiale) et \textbf{verbe + préposition} (\eng{listen to}, \eng{wait for}).\\
$\square$ J'utilise la nominalisation correcte (\eng{a break-up}, \eng{a check-up}) dans mon expression écrite.\\
$\square$ Je sais remplacer un phrasal verb par un verbe simple (registre formel) : \eng{put off} $\rightarrow$ \eng{postpone}, \eng{find out} $\rightarrow$ \eng{discover}.\\
$\square$ À l'oral, je prononce la \textbf{particule accentuée} (stress on particle) pour les verbes intransitives/séparables (\eng{turn \textbf{OFF}}).\\
$\square$ Je relis mes phrases pour vérifier l'ordre : Sujet + Verbe (+ Objet/Pronom) + Particule + (Préposition + Objet).
\end{aretenir}

\findecours

\end{document}
